% Kinematics of Particles in a plane — Pure Mathematics with Mechanics Upper Sixth — Cours signé OnBuch+
% Official MINESEC syllabus. Compiler avec : tectonic PMM11-kinematics-of-particles-in-a-plane.tex
\newcommand{\DOCMATIERE}{Pure Mathematics with Mechanics}
\newcommand{\DOCNIVEAU}{Upper Sixth}
\newcommand{\DOCMODULE}{Upper Sixth — Pure mathematics with mechanics}
\newcommand{\DOCLECON}{11}
\newcommand{\DOCTITRE}{Kinematics of Particles in a plane}
% OnBuch+ — préambule « cours signés » ENRICHI (compilé avec tectonic / XeLaTeX)
% Système visuel « Pop » d'OnBuch+ : fond crème (jamais blanc), bordures encre
% épaisses, ombres « dures » décalées, titres Archivo Black en capitales.
% Version MATHÉMATIQUES : graphes (pgfplots/tikz), boîtes pédagogiques
% (définition, propriété, méthode, attention, le savais-tu, exercice résolu,
% exercice, TP, bilan), page de garde et signature OnBuch+ ancrée en bas.
%
% Macros attendues dans main.tex AVANT \input{preamble} :
%   \DOCMATIERE  \DOCNIVEAU  \DOCMODULE  \DOCLECON (numéro)  \DOCTITRE
\documentclass[11pt]{article}

\usepackage{fontspec}
\usepackage[english]{babel}
\usepackage[a4paper,margin=2cm,top=3.4cm,bottom=2.8cm,headheight=1.4cm,footskip=1.3cm]{geometry}
\usepackage{amsmath,amssymb,amsfonts}
\usepackage{mathtools} % \xmapsto, \coloneqq... souvent utilisés par les modèles
\usepackage{cancel} % simplifications barrées (bilans de forces)
% \abs{z} (module, valeur absolue) et \norm{u} (norme), utilisés par les modèles
\providecommand{\abs}{}\let\abs\relax\DeclarePairedDelimiter{\abs}{\lvert}{\rvert}
\providecommand{\norm}{}\let\norm\relax\DeclarePairedDelimiter{\norm}{\lVert}{\rVert}
\usepackage[table]{xcolor} % [table] : fournit aussi \rowcolors (alternance de lignes)
\usepackage{pagecolor}
\usepackage{tikz}
\usetikzlibrary{arrows.meta,positioning,calc,shapes.geometric,shapes.misc,shapes.arrows,decorations.pathmorphing,decorations.pathreplacing,decorations.markings,patterns,shadows,mindmap,trees,fit,backgrounds,matrix,intersections,angles,quotes}
\usepackage{pgfplots}
\usepackage[version=4]{mhchem} % équations nucléaires \ce{^{235}_{92}U}
\usepackage[european,siunitx]{circuitikz} % schémas électriques
\pgfplotsset{compat=1.18}
\usepgfplotslibrary{fillbetween,groupplots} % aires entre courbes (calcul intégral)
\usepackage{enumitem}
\usepackage{fancyhdr}
% [most] charge toutes les bibliothèques tcolorbox (listings, minted, raster,
% documentation...) et gonfle inutilement la mémoire TeX sur un document long
% (~300 boîtes) : on ne charge que ce qui est réellement utilisé.
\usepackage[skins,breakable]{tcolorbox}
\usepackage{graphicx}
\usepackage{colortbl}
\usepackage{booktabs}
\usepackage{tabularx}
\usepackage{diagbox} % case d'en-tête diagonale des tableaux à double entrée
% Colonnes extensibles centrée / gauche / droite (C, L, R), souvent utilisées par les modèles
\newcolumntype{C}{>{\centering\arraybackslash}X}
\newcolumntype{L}{>{\raggedright\arraybackslash}X}
\newcolumntype{R}{>{\raggedleft\arraybackslash}X}
\usepackage{array}
\usepackage{multicol}
\usepackage{siunitx}
% parse-numbers=false : accepte aussi \SI{2,5 \times 10^{-3}}{...}
\sisetup{locale=UK,output-decimal-marker={.},group-minimum-digits=4,parse-numbers=false}
\DeclareSIUnit\ohm{\ensuremath{\symup{\Omega}}} % U+2126 absent de la police
\DeclareSIUnit\division{div} % réglages d'oscilloscope (ms/div, V/div)
\DeclareSIUnit\tr{tr} % tours (vitesse de rotation, ex. \SI{1500}{\tr\per\minute})
\usepackage{qrcode}
% \degree : commande fréquemment utilisée par les modèles pour un angle ou une
% température en dehors de \SI{}{} (non fournie par les paquets chargés).
\providecommand{\degree}{^{\circ}}
% \qed (fin de démonstration), utilisé par les modèles sans amsthm
\providecommand{\qed}{\hfill$\square$}
% \barre : double barre des valeurs interdites dans un tableau de variations (tkz-tab)
\providecommand{\barre}{\ensuremath{\|}}
% environnement proof (amsthm non chargé) utilisé par les modèles
\ifdefined\proof\else\newenvironment{proof}[1][Proof]{\par\noindent\textit{#1.}\ }{\qed\par}\fi
\usepackage{lastpage}
\usepackage{hyperref}
\hypersetup{hidelinks}

% ============================================================
% Palette « Pop » OnBuch+ — mêmes hex que la classe P (plus_ui.dart)
% ============================================================
\definecolor{poporange}{HTML}{F59321}
\definecolor{popdark}{HTML}{E07A0C}
\definecolor{popcream}{HTML}{FFF6E8}
\definecolor{popink}{HTML}{1C1714}
\definecolor{poppurple}{HTML}{7A5AE0}
\definecolor{popgreen}{HTML}{2E9E6A}
\definecolor{poppink}{HTML}{E0457B}
\definecolor{popblue}{HTML}{2F7DE1}
\definecolor{popgold}{HTML}{C98A12}
\definecolor{popmuted}{HTML}{8A7F76}
\definecolor{popline}{HTML}{EADFD2}
\definecolor{popgray}{HTML}{8A7F76}
\colorlet{popgrey}{popgray}
% Alias tolérants (le modèle nomme parfois les couleurs différemment)
\colorlet{popwhite}{white}
\colorlet{popblack}{popink}
\colorlet{popred}{poppink}
\colorlet{popyellow}{popgold}
% Teintes claires pour fonds de tableaux / zones
\colorlet{popblueL}{popblue!10!white}
\colorlet{poporangeL}{poporange!14!white}
\colorlet{popgreenL}{popgreen!12!white}
\colorlet{poppurpleL}{poppurple!12!white}
\colorlet{poppinkL}{poppink!12!white}
\colorlet{popgoldL}{popgold!14!white}

\pagecolor{popcream}

% ============================================================
% Typographie
% ============================================================
\defaultfontfeatures{Ligatures=TeX}
\setmainfont{PlusJakartaSans-Medium.ttf}[Path=../../../fonts/,BoldFont=PlusJakartaSans-Bold.ttf]
% Maths en sans-serif (Fira Math) avec chiffres et lettres droites de Plus
% Jakarta, pour que les formules se fondent dans le texte.
\usepackage[math-style=ISO,bold-style=ISO]{unicode-math}
\setmathfont{FiraMath-Regular.otf}[Path=../../../fonts/]
\setmathfont{PlusJakartaSans-Medium.ttf}[Path=../../../fonts/,range={up/num,up/{Latin,latin},\mathup}]
% (icomma retiré : décimales avec point en anglais)
% Caractères mathématiques Unicode absents des polices de texte (Plus Jakarta,
% Archivo Black) : rendus par leur équivalent mathématique, en texte comme en
% formule — sinon ils s'affichent en carré vide (ex. « Arithmétique dans ℤ »).
\usepackage{newunicodechar}
\newunicodechar{ℝ}{\ensuremath{\mathbb{R}}}
\newunicodechar{ℤ}{\ensuremath{\mathbb{Z}}}
\newunicodechar{ℂ}{\ensuremath{\mathbb{C}}}
\newunicodechar{ℕ}{\ensuremath{\mathbb{N}}}
\newunicodechar{ℚ}{\ensuremath{\mathbb{Q}}}
\newunicodechar{𝔻}{\ensuremath{\mathbb{D}}}
\newunicodechar{α}{\ensuremath{\alpha}}
\newunicodechar{β}{\ensuremath{\beta}}
\newunicodechar{γ}{\ensuremath{\gamma}}
\newunicodechar{δ}{\ensuremath{\delta}}
\newunicodechar{ε}{\ensuremath{\varepsilon}}
\newunicodechar{θ}{\ensuremath{\theta}}
\newunicodechar{λ}{\ensuremath{\lambda}}
\newunicodechar{μ}{\ensuremath{\mu}}
\newunicodechar{σ}{\ensuremath{\sigma}}
\newunicodechar{τ}{\ensuremath{\tau}}
\newunicodechar{φ}{\ensuremath{\varphi}}
\newunicodechar{ψ}{\ensuremath{\psi}}
\newunicodechar{ω}{\ensuremath{\omega}}
\newunicodechar{Σ}{\ensuremath{\Sigma}}
% majuscules produites par \MakeUppercase dans les titres (α-aminés -> Α)
\newunicodechar{Α}{\ensuremath{\alpha}}
\newunicodechar{Β}{\ensuremath{\beta}}
\newunicodechar{Γ}{\ensuremath{\Gamma}}
\newunicodechar{Δ}{\ensuremath{\Delta}}
\newunicodechar{Ε}{\ensuremath{\varepsilon}}
\newunicodechar{Θ}{\ensuremath{\Theta}}
\newunicodechar{Λ}{\ensuremath{\Lambda}}
\newunicodechar{Μ}{\ensuremath{\mu}}
\newunicodechar{Τ}{\ensuremath{\tau}}
\newunicodechar{Φ}{\ensuremath{\Phi}}
\newunicodechar{Ψ}{\ensuremath{\Psi}}
\newunicodechar{Ω}{\ensuremath{\Omega}}
\newunicodechar{↦}{\ensuremath{\mapsto}}
\newunicodechar{∈}{\ensuremath{\in}}
\newunicodechar{∉}{\ensuremath{\notin}}
\newunicodechar{∧}{\ensuremath{\wedge}}
\newunicodechar{⇔}{\ensuremath{\Leftrightarrow}}
\newunicodechar{⟺}{\ensuremath{\Longleftrightarrow}}
\newunicodechar{⇒}{\ensuremath{\Rightarrow}}
\newunicodechar{⟹}{\ensuremath{\Longrightarrow}}
\newunicodechar{∘}{\ensuremath{\circ}}
\newunicodechar{∪}{\ensuremath{\cup}}
\newunicodechar{∩}{\ensuremath{\cap}}
\newunicodechar{≡}{\ensuremath{\equiv}}
\newunicodechar{⊂}{\ensuremath{\subset}}
\newunicodechar{⊥}{\ensuremath{\perp}}
\newunicodechar{∃}{\ensuremath{\exists}}
\newunicodechar{∀}{\ensuremath{\forall}}
\newunicodechar{∛}{\ensuremath{\sqrt[3]{\,}}}
\newunicodechar{∜}{\ensuremath{\sqrt[4]{\,}}}
\newunicodechar{⃗}{\ensuremath{{}^{\to}}}
\newunicodechar{ⁿ}{\ifmmode^{n}\else\textsuperscript{\ensuremath{n}}\fi}
\newunicodechar{ˣ}{\ifmmode^{x}\else\textsuperscript{\ensuremath{x}}\fi}
\newunicodechar{ᵇ}{\ifmmode^{b}\else\textsuperscript{\ensuremath{b}}\fi}
\newunicodechar{ʳ}{\ifmmode^{r}\else\textsuperscript{\ensuremath{r}}\fi}
\newunicodechar{ᵘ}{\ifmmode^{u}\else\textsuperscript{\ensuremath{u}}\fi}
\newunicodechar{ʸ}{\ifmmode^{y}\else\textsuperscript{\ensuremath{y}}\fi}
\newunicodechar{ᵃ}{\ifmmode^{a}\else\textsuperscript{\ensuremath{a}}\fi}
\newunicodechar{ᵉ}{\ifmmode^{e}\else\textsuperscript{\ensuremath{e}}\fi}
\newunicodechar{ᵏ}{\ifmmode^{k}\else\textsuperscript{\ensuremath{k}}\fi}
\newunicodechar{ᵖ}{\ifmmode^{p}\else\textsuperscript{\ensuremath{p}}\fi}
\newunicodechar{ᴺ}{\ifmmode^{N}\else\textsuperscript{\ensuremath{N}}\fi}
\newunicodechar{ᶜ}{\ifmmode^{c}\else\textsuperscript{\ensuremath{c}}\fi}
\newunicodechar{ʰ}{\ifmmode^{h}\else\textsuperscript{\ensuremath{h}}\fi}
\newunicodechar{ᵠ}{\ifmmode^{q}\else\textsuperscript{\ensuremath{q}}\fi}
\newunicodechar{ₙ}{\ifmmode_{n}\else\textsubscript{\ensuremath{n}}\fi}
\newunicodechar{ₐ}{\ifmmode_{a}\else\textsubscript{\ensuremath{a}}\fi}
\newunicodechar{ᵢ}{\ifmmode_{i}\else\textsubscript{\ensuremath{i}}\fi}
\newunicodechar{ᵣ}{\ifmmode_{r}\else\textsubscript{\ensuremath{r}}\fi}
\newunicodechar{ₖ}{\ifmmode_{k}\else\textsubscript{\ensuremath{k}}\fi}
\newunicodechar{ᵦ}{\ifmmode_{\beta}\else\textsubscript{\ensuremath{\beta}}\fi}
\newunicodechar{ₓ}{\ifmmode_{x}\else\textsubscript{\ensuremath{x}}\fi}
\newfontfamily\popdisplay{ArchivoBlack-Regular.ttf}[Path=../../../fonts/]
\newfontfamily\poplogo{SpaceGrotesk-Bold.ttf}[Path=../../../fonts/]
\newfontfamily\popsemi{PlusJakartaSans-SemiBold.ttf}[Path=../../../fonts/]
\newfontfamily\popxbold{PlusJakartaSans-ExtraBold.ttf}[Path=../../../fonts/]
\newfontfamily\popxboldls{PlusJakartaSans-ExtraBold.ttf}[Path=../../../fonts/,LetterSpace=3.0]

\setlist[enumerate]{leftmargin=1.7em,itemsep=5pt,topsep=4pt}
\setlist[itemize]{leftmargin=1.7em,itemsep=4pt,topsep=4pt,label={\color{poporange}\textbullet}}
\setlength{\parindent}{0pt}
\setlength{\parskip}{5pt}
\renewcommand{\arraystretch}{1.25}

% pgfplots : style Pop par défaut
\pgfplotsset{
  every axis/.append style={axis line style={popink,line width=1pt},tick style={popink},
    grid style={popline,line width=0.5pt},label style={font=\small},tick label style={font=\footnotesize}},
  popcurve/.style={poporange,line width=1.8pt,no markers,smooth},
}
% Même style disponible dans un \draw TikZ simple (hors pgfplots)
\tikzset{popcurve/.style={poporange,line width=1.8pt}}
\tikzset{popgrid/.style={popline,very thin}}
\colorlet{popcurve}{poporange} % le nom du style est parfois utilisé comme couleur

% ============================================================
% Pill / squiggle
% ============================================================
\newcommand{\poppill}[3]{%
  \tikz[baseline=-0.55ex]\node[draw=popink,line width=1.2pt,rounded corners=2.6pt,%
    fill=#2,inner xsep=6.5pt,inner ysep=3pt]{%
    \popxboldls\fontsize{7}{7}\selectfont\color{#3}\MakeUppercase{#1}};%
}
\newcommand{\popsquiggle}[1]{%
  \tikz[baseline=-0.4ex]{
    \draw[#1,line width=2.6pt,line cap=round]
      plot [smooth] coordinates {(0,0.09) (0.34,0.01) (0.68,0.21) (1.02,0.01) (1.36,0.10)};
  }%
}
% Mot-clé surligné (marqueur orange sous le texte)
\newcommand{\cle}[1]{{\popxbold\color{popdark}#1}}

% ============================================================
% En-tête / pied
% ============================================================
\pagestyle{fancy}
\fancyhf{}
\renewcommand{\headrulewidth}{0pt}
\renewcommand{\footrulewidth}{0pt}
\lhead{\raisebox{-0.15ex}{\poplogo\fontsize{13}{13}\selectfont\color{popink}OnBuch\color{poporange}+}}
\rhead{\poppill{\DOCMATIERE\ \textbullet\ \DOCNIVEAU\ \textbullet\ Chapter \DOCLECON}{white}{popink}}
\lfoot{\popxbold\fontsize{7.6}{7.6}\selectfont\color{popmuted}\MakeUppercase{Signed course OnBuch+}\ \textbullet\ {\popsemi\DOCTITRE}}
\rfoot{\tikz[baseline=-0.6ex]\node[rounded corners=8.5pt,fill=popink,minimum height=17pt,minimum width=17pt,inner xsep=5pt,inner sep=0]{\popxbold\fontsize{8}{8}\selectfont\color{popcream}\thepage\,/\,\pageref*{LastPage}};}

% ============================================================
% Titres
% ============================================================
\newcommand{\chaptitre}[2]{%
  \begin{tcolorbox}[enhanced,colback=poporange,colframe=popink,boxrule=2.6pt,arc=11pt,
      drop shadow=popink,left=15pt,right=15pt,top=12pt,bottom=11pt,before skip=3pt,after skip=18pt]
    \poppill{#2}{popink}{popcream}\par\vspace{9pt}
    {\raggedright\hyphenpenalty=10000\exhyphenpenalty=10000\popdisplay\fontsize{22}{24}\selectfont\color{popcream}\MakeUppercase{#1}\par}
  \end{tcolorbox}
}
% Section numérotée : pastille orange avec le numéro + titre Archivo
\newcounter{popsec}
\newcommand{\coursec}[1]{%
  \stepcounter{popsec}\setcounter{popsub}{0}%
  \par\vspace{16pt}\noindent
  \begin{minipage}{\linewidth}
  \tikz[baseline=-0.7ex]\node[draw=popink,line width=1.6pt,fill=poporange,rounded corners=4pt,minimum size=22pt,inner sep=0]{\popdisplay\fontsize{12}{12}\selectfont\color{popcream}\thepopsec};%
  \hspace{8pt}{\popdisplay\fontsize{15}{17}\selectfont\color{popink}\MakeUppercase{#1}}\par
  \vspace{4pt}\noindent{\color{poporange}\rule{36pt}{3.6pt}}
  \end{minipage}\par\nopagebreak\vspace{8pt}
}
\newcounter{popsub}
\newcommand{\courssub}[1]{%
  \stepcounter{popsub}%
  \par\vspace{9pt}\noindent{\popxbold\fontsize{11.5}{13}\selectfont\color{popink}{\color{poporange}\thepopsec.\thepopsub}\hspace{6pt}#1}\par\nopagebreak\vspace{4pt}
}

% ============================================================
% Boîtes pédagogiques — même recette : fond blanc, bordure épaisse colorée,
% ombre dure encre, badge plein. Titre optionnel [#1].
% ============================================================
\tcbset{popbase/.style={breakable,enhanced,colback=white,boxrule=2.3pt,arc=9pt,
  shadow={3pt}{-3pt}{0pt}{popink},left=14pt,right=14pt,top=11pt,bottom=11pt,before skip=11pt,after skip=15pt,
  fontupper=\popsemi\fontsize{10.6}{14.5}\selectfont\color{popink},
  fontlower=\popsemi\fontsize{10.6}{14.5}\selectfont\color{popink}}}
% Coche dessinée (le glyphe ✓ n'existe pas dans les polices du document)
\newcommand{\popcheck}[1]{\tikz[baseline=-0.5ex]\draw[#1,line width=1.6pt,line cap=round,line join=round](-0.13,0.0)--(-0.04,-0.09)--(0.13,0.1);}
\providecommand{\coche}{\popcheck{popgreen}}
\providecommand{\resultat}[1]{\par\smallskip\noindent\fcolorbox{popgreen}{popgreenL}{\parbox{\dimexpr\linewidth-2\fboxsep-2\fboxrule\relax}{\textbf{Result.} #1}}\par\smallskip}
\newcommand{\popboxhead}[3]{\poppill{#1}{#2}{white}\ifx\relax#3\relax\else\hspace{6pt}{\popxbold\fontsize{10.6}{12}\selectfont\color{popink}#3}\fi\par\vspace{7pt}}

\newtcolorbox{aretenir}[1][]{popbase,colframe=popgreen,before upper={\popboxhead{Key points}{popgreen}{#1}}}
\newtcolorbox{exemplebox}[1][]{popbase,colframe=poporange,before upper={\popboxhead{Exemple}{poporange}{#1}}}
\newtcolorbox{definition}[1][]{popbase,colframe=popblue,before upper={\popboxhead{Definition}{popblue}{#1}}}
\newtcolorbox{propriete}[1][]{popbase,colframe=poppurple,before upper={\popboxhead{Property}{poppurple}{#1}}}
\newtcolorbox{methode}[1][]{popbase,colframe=popgold,before upper={\popboxhead{Method}{popgold}{#1}}}
\newtcolorbox{attention}[1][]{popbase,colframe=poppink,before upper={\popboxhead{Warning}{poppink}{#1}}}
\newtcolorbox{experience}[1][]{popbase,colframe=popblue,colback=popblueL,before upper={\popboxhead{Experiment}{popblue}{#1}}}
% « Le savais-tu ? » : carte violette pleine (contexte, culture, Cameroun)
\newtcolorbox{savaistu}[1][]{popbase,colback=poppurple,colframe=popink,
  fontupper=\popsemi\fontsize{10.4}{14.2}\selectfont\color{white},
  before upper={\poppill{Did you know?}{white}{poppurple}\ifx\relax#1\relax\else\hspace{6pt}{\popxbold\color{white}#1}\fi\par\vspace{7pt}}}
% Exercice résolu : énoncé en haut, \tcblower puis la solution sur fond crème
\newcounter{popexo}\newcounter{popexr}
\newtcolorbox{exoresolu}[1][]{popbase,colframe=poporange,colbacklower=popcream,
  segmentation style={popink,line width=1.2pt,dashed},
  before upper={\stepcounter{popexr}\popboxhead{Worked example \thepopexr}{poporange}{#1}},
  before lower={\poppill{Solution}{popink}{popcream}\par\vspace{6pt}}}
% Exercice (non résolu en place) — la correction est dans la partie Corrigés
\newtcolorbox{exercice}[1][]{popbase,colframe=popink,
  before upper={\stepcounter{popexo}\popboxhead{Exercise \thepopexo}{popink}{#1}}}
% Corrigé
\newtcolorbox{corrige}[1][]{popbase,colframe=popgreen,colback=popgreenL,
  before upper={\popboxhead{Solution}{popgreen}{#1}}}
% Objectifs / prérequis / bilan
\newtcolorbox{objectifs}{popbase,colframe=popink,colback=poporangeL,
  before upper={\popboxhead{Chapter objectives}{popink}{}}}
\newtcolorbox{prerequis}{popbase,colframe=popink,colback=popblueL,
  before upper={\popboxhead{Prerequisites}{popblue}{}}}
\newtcolorbox{bilan}{popbase,colframe=popink,colback=white,boxrule=2.8pt,
  before upper={\popboxhead{Fiche bilan}{poporange}{L'essentiel à connaître pour l'examen}}}
% Figure encadrée avec légende
\newtcolorbox{popfigure}[1][]{enhanced,breakable=false,colback=white,colframe=popline,boxrule=1.4pt,arc=8pt,
  left=10pt,right=10pt,top=10pt,bottom=8pt,before skip=10pt,after skip=12pt,halign=center,
  lower separated=false,halign lower=center,
  fontlower=\popsemi\fontsize{9}{11.5}\selectfont\color{popmuted},
  #1}
\newcommand{\legende}[1]{\par\vspace{6pt}{\popsemi\fontsize{9}{11.5}\selectfont\color{popmuted}#1}}

% ============================================================
% Signature OnBuch+
% ============================================================
\newcommand{\poplogoinline}[1]{{\poplogo\fontsize{#1}{#1}\selectfont\color{popink}OnBuch\color{poporange}+}}
\newcommand{\signatureonbuch}{%
  \begin{tcolorbox}[enhanced,colback=popink,colframe=popink,boxrule=2.2pt,arc=10pt,
      drop shadow=poporange,left=14pt,right=14pt,top=10pt,bottom=10pt,before skip=0pt,after skip=0pt]
    \tikz[baseline=-0.7ex]\node[circle,fill=popgreen,draw=popcream,line width=1.3pt,minimum size=22pt,inner sep=0]{\popcheck{white}};%
    \hspace{9pt}\begin{minipage}[c]{0.62\linewidth}
      {\popxbold\fontsize{12}{13}\selectfont\color{popcream}Signed\ \textbullet\ {\poplogo OnBuch\color{poporange}+}}\par\vspace{2pt}
      {\raggedright\popsemi\fontsize{8}{10.5}\selectfont\color{popline}\MakeUppercase{OnBuch+ teaching team}\\\MakeUppercase{Follows the official MINESEC syllabus}\par}
    \end{minipage}\hfill
    {\poplogo\fontsize{22}{22}\selectfont\color{popcream}OnBuch\color{poporange}+}
  \end{tcolorbox}%
}

% ============================================================
% Page de garde
% #1 titre · #2 matière · #3 niveau · #4 module · #5 n° leçon · #6 durée ·
% #7 description / objectifs (texte ou itemize)
% ============================================================
\newcommand{\pagegarde}[7]{%
  \thispagestyle{empty}
  \newgeometry{a4paper,margin=1.7cm,top=1.6cm,bottom=1.4cm}%
  \begin{tikzpicture}[remember picture,overlay]
    \fill[poporange!18] ([xshift=-3.2cm,yshift=-3.6cm]current page.north east) circle (4.6cm);
    \fill[poppurple!14] ([xshift=2.4cm,yshift=7.6cm]current page.south west) circle (3.8cm);
  \end{tikzpicture}%
  {\poplogo\fontsize{30}{30}\selectfont\color{popink}OnBuch\color{poporange}+}\hfill
  \raisebox{0.6ex}{\poppill{Signed course \textbullet\ Premium}{popink}{popcream}}\par
  \vspace{1pt}\popsquiggle{poppurple}\par
  \vspace{8pt}
  \begin{tcolorbox}[enhanced,colback=poporange,colframe=popink,boxrule=2.8pt,arc=13pt,
      drop shadow=popink,left=18pt,right=18pt,top=12pt,bottom=13pt,after skip=10pt]
    \poppill{#2 \textbullet\ #3}{popink}{popcream}\hspace{5pt}\poppill{#4}{white}{popink}\par\vspace{8pt}
    {\popdisplay\fontsize{14}{14}\selectfont\color{popink}CHAPTER #5}\par\vspace{4pt}
    {\raggedright\hyphenpenalty=10000\exhyphenpenalty=10000\popdisplay\fontsize{25}{27}\selectfont\color{popcream}\MakeUppercase{#1}\par}
    \vspace{8pt}
    {\popxbold\fontsize{9.5}{9.5}\selectfont\color{popink}\MakeUppercase{Suggested time: #6}}
  \end{tcolorbox}
  \begin{tcolorbox}[enhanced,colback=white,colframe=popink,boxrule=2.2pt,arc=10pt,
      drop shadow=popink,left=16pt,right=16pt,top=10pt,bottom=10pt,after skip=8pt]
    \poppill{In this chapter}{poppurple}{white}\par\vspace{5pt}
    \popsemi\fontsize{9.3}{12.6}\selectfont\color{popink}#7
  \end{tcolorbox}
  \noindent\begin{minipage}[t]{0.487\linewidth}
    \begin{tcolorbox}[enhanced,colback=popgreen,colframe=popink,boxrule=2.2pt,arc=10pt,
        drop shadow=popink,left=11pt,right=11pt,top=10pt,bottom=10pt,height=3.1cm,valign=center]
      \tcbox[colback=white,colframe=popink,boxrule=1.3pt,arc=5pt,boxsep=0pt,nobeforeafter,
        left=3pt,right=3pt,top=3pt,bottom=3pt,valign=center]{\qrcode[height=1.9cm]{https://play.google.com/store/apps/details?id=com.luvvix.onbuch&hl=en}}%
      \hspace{8pt}\begin{minipage}[c]{0.5\linewidth}
        {\popdisplay\fontsize{11}{12.5}\selectfont\color{white}\MakeUppercase{Download OnBuch}}\par\vspace{4pt}
        {\raggedright\popsemi\fontsize{8.4}{11}\selectfont\color{white}Free \textbullet\ Google Play\\AI tutor, past papers, results\par}
      \end{minipage}
    \end{tcolorbox}
  \end{minipage}\hfill
  \begin{minipage}[t]{0.487\linewidth}
    \begin{tcolorbox}[enhanced,colback=poppurple,colframe=popink,boxrule=2.2pt,arc=10pt,
        drop shadow=popink,left=11pt,right=11pt,top=10pt,bottom=10pt,height=3.1cm,valign=center]
      \tikz[baseline=-0.7ex]\node[circle,fill=white,draw=popink,line width=1.3pt,minimum size=1.6cm,inner sep=0]{\popdisplay\fontsize{24}{24}\selectfont\color{poppurple}+};%
      \hspace{8pt}\begin{minipage}[c]{0.6\linewidth}
        {\popdisplay\fontsize{11}{12.5}\selectfont\color{white}\MakeUppercase{Go OnBuch+}}\par\vspace{4pt}
        {\raggedright\popsemi\fontsize{8.4}{11}\selectfont\color{white}All signed courses, unlimited AI tutor, PDF sheets — OnBuch+ tab in the app\par}
      \end{minipage}
    \end{tcolorbox}
  \end{minipage}
  \vspace{8pt}
  \popcontenu
  \vfill
  \signatureonbuch
  \restoregeometry
  \newpage
}

% Rangée « Ce cours contient » (page de garde)
\newcommand{\poptile}[3]{% #1 couleur #2 gros texte #3 légende
  \begin{tcolorbox}[enhanced,colback=#1,colframe=popink,boxrule=2pt,arc=9pt,shadow={2.5pt}{-2.5pt}{0pt}{popink},
      left=6pt,right=6pt,top=9pt,bottom=9pt,halign=center,valign=center,height=1.75cm,width=\linewidth,nobeforeafter]
    {\popdisplay\fontsize{12.5}{13}\selectfont\color{white}\MakeUppercase{#2}}\par\vspace{3pt}
    {\centering\popsemi\fontsize{7.8}{9.5}\selectfont\color{white}#3\par}
  \end{tcolorbox}}
\newcommand{\popcontenu}{%
  \par\noindent{\popxboldls\fontsize{7.5}{7.5}\selectfont\color{popmuted}THIS SIGNED COURSE CONTAINS}\par\vspace{7pt}
  \noindent\begin{minipage}[t]{0.235\linewidth}\poptile{popblue}{Course}{complete, illustrated, step by step}\end{minipage}\hfill
  \begin{minipage}[t]{0.235\linewidth}\poptile{poporange}{Methods}{and worked examples}\end{minipage}\hfill
  \begin{minipage}[t]{0.235\linewidth}\poptile{poppink}{Exercises}{exam style + solutions}\end{minipage}\hfill
  \begin{minipage}[t]{0.235\linewidth}\poptile{popgreen}{Summary sheet}{the essentials to remember}\end{minipage}\par}

% Sommaire
\newtcolorbox{sommaire}{enhanced,colback=white,colframe=popink,boxrule=2.2pt,arc=10pt,
  drop shadow=popink,left=18pt,right=18pt,top=13pt,bottom=13pt,before skip=6pt,after skip=20pt,
  before upper={\poppill{Contents}{poppurple}{white}\par\vspace{9pt}\popsemi\fontsize{10.6}{14.5}\selectfont\color{popink}}}

% Signature de fin de cours — poussée tout en bas de la dernière page
\newcommand{\findecours}{%
  \par\vfill\nopagebreak
  \begin{center}\popsquiggle{poporange}\end{center}\vspace{4pt}
  \signatureonbuch
}
\AtBeginDocument{\def\llbracket{\mathopen{[\mkern-3.5mu[}}\def\rrbracket{\mathclose{]\mkern-3.5mu]}}} % intervalles d'entiers (glyphe ⟦ absent de la police)
% Glyphes absents de Fira Math : redéfinis à partir de glyphes disponibles
\AtBeginDocument{%
  \def\setminus{\mathbin{\backslash}}\let\smallsetminus\setminus
  \def\vdots{\mathord{\vbox{\baselineskip3.2pt\lineskiplimit0pt\kern4pt\hbox{.}\hbox{.}\hbox{.}}}}%
  \def\ddots{\mathinner{\mkern1mu\raise6.5pt\hbox{.}\mkern2mu\raise3.6pt\hbox{.}\mkern2mu\raise0.7pt\hbox{.}\mkern1mu}}%
  \def\triangle{\mathord{\scalebox{0.9}{$\symup{\Delta}$}}}%
  \def\nabla{\mathord{\raisebox{\depth}{\scalebox{1}[-1]{$\symup{\Delta}$}}}}%
  \def\checkmark{\popcheck{popdark}}%
  \def\star{\mathbin{\ast}}\def\bigstar{\mathord{\ast}}%
  \def\diamond{\mathbin{\scalebox{0.6}{$\Diamond$}}}%
  \def\bigcup{\mathop{\mathchoice{\raisebox{-0.3em}{\scalebox{1.7}{$\cup$}}}{\scalebox{1.25}{$\cup$}}{\cup}{\cup}}\displaylimits}%
  \def\bigcap{\mathop{\mathchoice{\raisebox{-0.3em}{\scalebox{1.7}{$\cap$}}}{\scalebox{1.25}{$\cap$}}{\cap}{\cap}}\displaylimits}%
  \def\lesseqqgtr{\mathrel{\lessgtr}}\def\bigoplus{\mathop{\oplus}\displaylimits}%
}
\providecommand{\ssi}{if and only if} % abréviation fréquente
\AtBeginDocument{\providecommand{\wideparen}{\overparen}} % arc de cercle

%% FIGFIX-BEGIN v5 — correctifs globaux des figures (docs/onbuch-generation/tools/figfix)
\usepackage{adjustbox}\usepackage{etoolbox}\tcbuselibrary{raster}
% toute figure TikZ/pgfplots plus large que la ligne est réduite à la largeur disponible
\BeforeBeginEnvironment{tikzpicture}{\begin{adjustbox}{max width=\linewidth}}
\AfterEndEnvironment{tikzpicture}{\end{adjustbox}}
% légendes pgfplots lisibles par-dessus les courbes
\pgfplotsset{every axis/.append style={legend style={fill=white,fill opacity=0.92,text opacity=1,draw=black!15,inner sep=3pt}}}
% étiquettes d'arêtes (syntaxe edge["..."]) avec fond blanc
\tikzset{every edge quotes/.append style={fill=white,inner sep=1.2pt}}
%% FIGFIX-END
\begin{document}

\pagegarde{Kinematics of Particles in a plane}{Pure Mathematics with Mechanics}{Upper Sixth}{Upper Sixth — Pure mathematics with mechanics}{11}{9 periods}{This chapter extends the study of motion from a straight line to the plane: position, velocity and acceleration become vectors linked by differentiation and integration. It then applies these tools to projectile motion and vertical motion under gravity, two favourite themes of the GCE Advanced Level examination.
\vspace{4pt}
\begin{multicols}{2}\small
\begin{itemize}
\item By the end of this chapter I can represent the position of a particle in a plane by a vector r(t) = x(t)i + y(t)j and interpret its components.
\item By the end of this chapter I can obtain velocity and acceleration vectors by differentiating the position vector with respect to time.
\item By the end of this chapter I can recover velocity and position from acceleration (or velocity) by integration, using initial conditions to determine the constants.
\item By the end of this chapter I can define a projectile and state the assumptions of the standard projectile model (constant g, no air resistance, point mass).
\item By the end of this chapter I can derive and use the Cartesian equation of a projectile's trajectory and show it is a parabola.
\item By the end of this chapter I can calculate the range, maximum height and time of flight of a projectile and use them to solve targeting problems.
\end{itemize}\end{multicols}}

\begin{sommaire}
\begin{multicols}{2}
\begin{enumerate}
\item Position, Velocity and Acceleration as Vectors in a Plane
\item From Acceleration back to Velocity and Position: Integration in Kinematics
\item Projectile Motion: Definition, Model and Equations
\item Range, Maximum Height and Time of Flight
\item Applications and Problem-Solving with Projectiles
\item Free Fall, Vertical Projection and Motion Graphs
\item Integration activity
\item Methods and worked examples
\item Exercises
\item Solutions to the exercises
\item Summary sheet
\end{enumerate}
\end{multicols}
\end{sommaire}

\begin{objectifs}
\begin{itemize}
\item By the end of this chapter I can represent the position of a particle in a plane by a vector r(t) = x(t)i + y(t)j and interpret its components.
\item By the end of this chapter I can obtain velocity and acceleration vectors by differentiating the position vector with respect to time.
\item By the end of this chapter I can recover velocity and position from acceleration (or velocity) by integration, using initial conditions to determine the constants.
\item By the end of this chapter I can define a projectile and state the assumptions of the standard projectile model (constant g, no air resistance, point mass).
\item By the end of this chapter I can derive and use the Cartesian equation of a projectile's trajectory and show it is a parabola.
\item By the end of this chapter I can calculate the range, maximum height and time of flight of a projectile and use them to solve targeting problems.
\item By the end of this chapter I can solve applied projectile problems: clearing an obstacle, hitting a target, launching from a height, optimal angle of projection.
\item By the end of this chapter I can analyse free fall and vertical projection (upwards or downwards) using the constant-acceleration equations.
\item By the end of this chapter I can draw and interpret displacement–time, velocity–time and acceleration–time graphs for vertical motion.
\end{itemize}
\end{objectifs}

\begin{prerequis}
\begin{itemize}
\item Linear kinematics with constant acceleration (Lower Sixth): the suvat equations v = u + at, s = ut + ½at², v² = u² + 2as.
\item Differentiation and integration of polynomials, trigonometric and exponential functions (Pure Mathematics, Lower Sixth).
\item Vectors in component form: i and j unit vectors, magnitude and direction of a vector.
\item Trigonometry: sine, cosine, tangent, the identity sin(2θ) = 2 sinθ cosθ, solving simple trigonometric equations.
\item Quadratic equations and the Cartesian equation of a parabola.
\end{itemize}
\end{prerequis}

\newpage

\coursec{Position, Velocity and Acceleration as Vectors in a Plane}

During the inter-school football tournament at the Ahmadou Ahidjo Stadium in Yaoundé, a striker launches a free kick from 25 metres: will the ball clear the wall and dip under the crossbar? Meanwhile, on the Bamenda–Bafoussam road, a driver throws a mango peel out of a moving bus window and it lands far behind the vehicle. Even the nurse at the Mfoundi health centre wonders how high the water jet from a broken pipe will rise. All these everyday Cameroonian scenes obey the same mathematical laws: the kinematics of a particle moving in a plane. This chapter gives you the tools to describe, predict and even optimise such motions.

\courssub{The Position Vector in a Plane}

To describe the motion of a particle in a plane, we choose a fixed origin $O$ and two perpendicular axes $Ox$ and $Oy$ with unit vectors $\vec{i}$ and $\vec{j}$. At any time $t$, the particle is at a point $P$. Its \cle{position vector} is the vector $\vec{OP}$.

\begin{definition}[Position Vector]
The position vector of a particle $P$ at time $t$ is
\[
\vec{r}(t) = x(t)\,\vec{i} + y(t)\,\vec{j},
\]
where $x(t)$ and $y(t)$ are the Cartesian coordinates of $P$ at time $t$. The functions $x(t)$ and $y(t)$ are called the \cle{parametric equations} of the motion. The variable $t$ (time) is the parameter.
\end{definition}

The set of all points occupied by the particle during its motion forms its \cle{trajectory} (or path). The trajectory is a curve in the plane. Its \cle{Cartesian equation} is obtained by eliminating the parameter $t$ between $x = x(t)$ and $y = y(t)$.

\begin{methode}[Finding the Cartesian Equation of the Trajectory]
\begin{enumerate}
\item Write $x = x(t)$ and $y = y(t)$.
\item If possible, solve one equation for $t$ (express $t$ in terms of $x$ or $y$).
\item Substitute this expression into the other equation to obtain a relation between $x$ and $y$ only.
\item State any restrictions on $x$ or $y$ coming from the domain of $t$ (usually $t \ge 0$).
\end{enumerate}
\end{methode}

\begin{exemplebox}[Parametric to Cartesian]
A particle moves so that its position vector at time $t$ (seconds) is
\[
\vec{r}(t) = (3t^2)\,\vec{i} + (2t - 5)\,\vec{j}\quad (\text{metres}).
\]
Find the Cartesian equation of its trajectory.
\end{exemplebox}
\textbf{Solution}
We have $x = 3t^2$ and $y = 2t - 5$. Since $t \ge 0$, from $y = 2t - 5$ we get $t = \dfrac{y+5}{2}$. Substitute into $x$:
\[
x = 3\left(\frac{y+5}{2}\right)^2 = \frac{3}{4}(y+5)^2.
\]
This is a parabola opening in the positive $x$-direction with vertex at $(0,-5)$. Because $t \ge 0$, $y \ge -5$. The trajectory is the half-parabola $x = \frac{3}{4}(y+5)^2,\ y \ge -5$.

\courssub{Velocity: Derivative of Position}

The \cle{velocity} of a particle is the rate of change of its position vector with respect to time.

\begin{definition}[Velocity Vector]
If $\vec{r}(t) = x(t)\,\vec{i} + y(t)\,\vec{j}$, then the velocity vector is
\[
\vec{v}(t) = \frac{d\vec{r}}{dt} = \frac{dx}{dt}\,\vec{i} + \frac{dy}{dt}\,\vec{j} = \dot{x}\,\vec{i} + \dot{y}\,\vec{j}.
\]
The dot notation $\dot{x} = dx/dt$, $\dot{y} = dy/dt$ is often used for time derivatives.
\end{definition}

The \cle{speed} is the magnitude of the velocity vector:
\[
v = |\vec{v}| = \sqrt{\dot{x}^2 + \dot{y}^2}\quad (\text{units: m\,s}^{-1}).
\]

\begin{propriete}[Velocity and the Trajectory]
The velocity vector $\vec{v}(t)$ is always \cle{tangent} to the trajectory at the point $P$ and points in the direction of motion.
\end{propriete}

\begin{exemplebox}[Velocity, Speed and Direction]
For the particle with $\vec{r}(t) = (3t^2)\,\vec{i} + (2t - 5)\,\vec{j}$:
\begin{enumerate}
\item Find $\vec{v}(t)$ and the acceleration $\vec{a}(t)$.
\item Find the speed and direction of motion at $t = 2$ s.
\item Determine when the particle is moving parallel to the $x$-axis and parallel to the $y$-axis.
\end{enumerate}
\end{exemplebox}
\textbf{Solution}
\begin{enumerate}
\item $\vec{v}(t) = \frac{d}{dt}(3t^2)\,\vec{i} + \frac{d}{dt}(2t-5)\,\vec{j} = 6t\,\vec{i} + 2\,\vec{j}$.
$\vec{a}(t) = \frac{d\vec{v}}{dt} = 6\,\vec{i} + 0\,\vec{j} = 6\,\vec{i}$.
\item At $t = 2$: $\vec{v}(2) = 12\,\vec{i} + 2\,\vec{j}$.
Speed $v = \sqrt{12^2 + 2^2} = \sqrt{148} = 2\sqrt{37} \approx 12.17\ \mathrm{m\,s^{-1}}$.
Direction: the velocity vector makes an angle $\theta$ with the $x$-axis where $\tan\theta = \frac{2}{12} = \frac{1}{6}$, so $\theta \approx 9.46^\circ$ above the horizontal.
\item Parallel to $x$-axis $\Rightarrow$ vertical component zero: $\dot{y} = 2 = 0$, impossible. Never parallel to $x$-axis.
Parallel to $y$-axis $\Rightarrow$ horizontal component zero: $\dot{x} = 6t = 0 \Rightarrow t = 0$. At $t=0$, $\vec{v}(0) = 2\,\vec{j}$, moving vertically upward.
\end{enumerate}

\courssub{Acceleration: Derivative of Velocity}

The \cle{acceleration} is the rate of change of velocity with respect to time.

\begin{definition}[Acceleration Vector]
If $\vec{v}(t) = \dot{x}\,\vec{i} + \dot{y}\,\vec{j}$, then the acceleration vector is
\[
\vec{a}(t) = \frac{d\vec{v}}{dt} = \frac{d^2\vec{r}}{dt^2} = \ddot{x}\,\vec{i} + \ddot{y}\,\vec{j}.
\]
Units: $\mathrm{m\,s^{-2}}$.
\end{definition}

Unlike velocity, acceleration is \emph{not} generally tangent to the path. It can be decomposed into tangential and normal components, but in Cartesian coordinates we simply work with its $x$ and $y$ components.

\courssub{Worked Examples and GCE-Style Problems}

\begin{exoresolu}[GCE Style: Perpendicular Velocity]
A particle moves in the $xy$-plane with position vector
\[
\vec{r}(t) = (t^3 - 3t)\,\vec{i} + (2t^2 - 4)\,\vec{j},\quad t \ge 0.
\]
Find the time(s) when the velocity of the particle is perpendicular to the vector $\vec{u} = 2\,\vec{i} - \vec{j}$.
\tcblower
\textbf{Solution}
Velocity: $\vec{v}(t) = \frac{d\vec{r}}{dt} = (3t^2 - 3)\,\vec{i} + 4t\,\vec{j}$.
Two vectors are perpendicular iff their scalar (dot) product is zero.
$\vec{v}(t) \cdot \vec{u} = (3t^2 - 3)(2) + (4t)(-1) = 6t^2 - 6 - 4t = 0$.
$6t^2 - 4t - 6 = 0 \implies 3t^2 - 2t - 3 = 0$.
$t = \dfrac{2 \pm \sqrt{4 + 36}}{6} = \dfrac{2 \pm \sqrt{40}}{6} = \dfrac{1 \pm \sqrt{10}}{3}$.
Since $t \ge 0$, we take $t = \dfrac{1 + \sqrt{10}}{3} \approx 1.39\ \mathrm{s}$.
\end{exoresolu}

\begin{exoresolu}[Comprehensive Kinematics]
A particle $P$ has position vector $\vec{r}(t) = (4t - t^2)\,\vec{i} + (t^2 - 2t)\,\vec{j}$ metres, $t \ge 0$.
\begin{enumerate}
\item Find the Cartesian equation of the path of $P$.
\item Find the velocity and acceleration vectors.
\item At what time is the speed of $P$ minimum? Find this minimum speed.
\end{enumerate}
\tcblower
\textbf{Solution}
\begin{enumerate}
\item $x = 4t - t^2$, $y = t^2 - 2t$. Add: $x + y = 2t \Rightarrow t = \frac{x+y}{2}$.
Substitute into $y$: $y = \left(\frac{x+y}{2}\right)^2 - 2\left(\frac{x+y}{2}\right)$.
$4y = (x+y)^2 - 4(x+y) = x^2 + 2xy + y^2 - 4x - 4y$.
$x^2 + 2xy + y^2 - 4x - 8y = 0$. This is the Cartesian equation (a parabola).
\item $\vec{v} = (4 - 2t)\,\vec{i} + (2t - 2)\,\vec{j}$.
$\vec{a} = -2\,\vec{i} + 2\,\vec{j}$ (constant acceleration).
\item Speed squared: $v^2 = (4-2t)^2 + (2t-2)^2 = 4(2-t)^2 + 4(t-1)^2 = 4[(t-2)^2 + (t-1)^2] = 4[2t^2 - 6t + 5]$.
Minimise $v^2$: derivative $8(2t - 3) = 0 \Rightarrow t = 1.5\ \mathrm{s}$.
Minimum speed: $v_{\min} = \sqrt{4[2(1.5)^2 - 6(1.5) + 5]} = \sqrt{4[4.5 - 9 + 5]} = \sqrt{4 \times 0.5} = \sqrt{2} \approx 1.41\ \mathrm{m\,s^{-1}}$.
\end{enumerate}
\end{exoresolu}

\courssub{Graphical Representation}

\begin{popfigure}
\begin{tikzpicture}[scale=0.8, >=stealth]
% Axes
\draw[->] (-1,0) -- (7,0) node[right] {$x$};
\draw[->] (0,-1) -- (0,5) node[above] {$y$};
% Trajectory: parabola x = (y+5)^2 * 3/4, but we plot parametrically for t in [0,2.5]
\draw[popblue, thick, domain=0:2.5, smooth, variable=\t, samples=100]
plot ({3*\t*\t}, {2*\t-5});
% Points at t=0, 1, 2
\foreach \t/\pos in {0/below, 1/right, 2/above} {
    \pgfmathsetmacro{\x}{3*\t*\t}
    \pgfmathsetmacro{\y}{2*\t-5}
    \fill[poporange] (\x,\y) circle (2.5pt);
    \node[\pos] at (\x,\y) {\small $t=\t$};
}
% Position vector at t=1.5
\pgfmathsetmacro{\xt}{3*1.5*1.5}
\pgfmathsetmacro{\yt}{2*1.5-5}
\draw[->, popdark, thick] (0,0) -- (\xt,\yt) node[midway, above left] {$\vec{r}(t)$};
% Velocity vector at t=1.5 (tangent)
\pgfmathsetmacro{\vx}{6*1.5}
\pgfmathsetmacro{\vy}{2}
\draw[->, popgreen, thick] (\xt,\yt) -- (\xt+\vx*0.15,\yt+\vy*0.15) node[midway, above right] {$\vec{v}(t)$};
% Dashed lines to axes
\draw[dashed, popmuted] (\xt,0) node[below] {$x(t)$} -- (\xt,\yt) -- (0,\yt) node[left] {$y(t)$};
\end{tikzpicture}
\legende{Position vector $\vec{r}(t)$ from origin to point $P$, and velocity vector $\vec{v}(t)$ tangent to the trajectory at $P$.}
\end{popfigure}

\begin{popfigure}
\begin{tikzpicture}
\begin{axis}[
    axis equal image,
    axis lines=middle,
    xlabel=$x$, ylabel=$y$,
    xmin=-1, xmax=20,
    ymin=-6, ymax=6,
    xtick={0,5,10,15,20},
    ytick={-5,0,5},
    width=12cm, height=8cm,
    samples=100,
    domain=0:2.5,
]
% Parametric trajectory
\addplot[popblue, thick, ->] ({3*x^2}, {2*x-5});
% Velocity arrows at t=0.5, 1.5, 2.2 using pgfplotsinvokeforeach
\pgfplotsinvokeforeach{0.5,1.5,2.2}{
    \pgfmathsetmacro{\x}{3*#1*#1}
    \pgfmathsetmacro{\y}{2*#1-5}
    \pgfmathsetmacro{\vx}{6*#1}
    \pgfmathsetmacro{\vy}{2}
    \pgfmathsetmacro{\scale}{0.3}
    \addplot[popgreen, thick, -stealth, samples=2, domain=0:1] ({\x + \scale*\vx*x}, {\y + \scale*\vy*x});
    \pgfmathsetmacro{\lx}{\x+\scale*\vx*1.2}
    \pgfmathsetmacro{\ly}{\y+\scale*\vy*1.2}
    \node[popgreen, font=\small] at (axis cs:\lx,\ly) {$t=#1$};
}
\end{axis}
\end{tikzpicture}
\legende{Parametric plot of $\vec{r}(t) = (3t^2, 2t-5)$ with velocity vectors (green arrows) at three instants.}
\end{popfigure}

\courssub{Common Mistakes and Warnings}

\begin{attention}[Speed vs. Derivative of Distance]
\begin{itemize}
\item \textbf{Speed} is $|\vec{v}| = \sqrt{\dot{x}^2 + \dot{y}^2}$. It is a scalar.
\item The derivative of the distance from the origin, $\frac{d}{dt}|\vec{r}|$, is \textbf{not} the speed (unless motion is purely radial).
\item Never compute speed as $\sqrt{x^2 + y^2}$ or as $|\dot{x}| + |\dot{y}|$. Always use the Pythagorean combination of the components of $\vec{v}$.
\item When asked for ``direction of motion'', give the angle the velocity vector makes with the positive $x$-axis (or a bearing), not the angle of the position vector.
\end{itemize}
\end{attention}

\courssub{Key Points to Remember}

\begin{aretenir}[Summary: Kinematics in a Plane]
\begin{itemize}
\item \textbf{Position:} $\vec{r}(t) = x(t)\,\vec{i} + y(t)\,\vec{j}$ (metres).
\item \textbf{Velocity:} $\vec{v}(t) = \dot{\vec{r}} = \dot{x}\,\vec{i} + \dot{y}\,\vec{j}$ (m\,s$^{-1}$). Tangent to path.
\item \textbf{Speed:} $v = |\vec{v}| = \sqrt{\dot{x}^2 + \dot{y}^2}$.
\item \textbf{Acceleration:} $\vec{a}(t) = \dot{\vec{v}} = \ddot{\vec{r}} = \ddot{x}\,\vec{i} + \ddot{y}\,\vec{j}$ (m\,s$^{-2}$).
\item \textbf{Trajectory:} Eliminate $t$ from $x=x(t)$, $y=y(t)$ to get $f(x,y)=0$.
\item \textbf{Perpendicular condition:} $\vec{v} \perp \vec{u} \iff \vec{v}\cdot\vec{u} = 0$.
\item \textbf{Parallel to axes:} Parallel to $x$-axis $\iff \dot{y}=0$; parallel to $y$-axis $\iff \dot{x}=0$.
\end{itemize}
\end{aretenir}

\begin{exercice}[Practice]
\begin{enumerate}
\item A particle moves with $\vec{r}(t) = (2t^2 + 3)\,\vec{i} + (t^3 - 4t)\,\vec{j}$. Find (a) $\vec{v}(t)$ and $\vec{a}(t)$, (b) the speed at $t=1$, (c) the time when the velocity is parallel to $\vec{i} + 2\vec{j}$.
\item The position vector of a particle is $\vec{r}(t) = (5\cos t)\,\vec{i} + (5\sin t)\,\vec{j}$, $t\ge 0$. (a) Show that the path is a circle. (b) Find $\vec{v}$ and $\vec{a}$. (c) Prove that $\vec{v} \cdot \vec{r} = 0$ and $\vec{a} \cdot \vec{v} = 0$. Interpret geometrically.
\item A particle is projected from the origin with velocity $\vec{u} = 10\,\vec{i} + 20\,\vec{j}$ m/s and experiences a constant acceleration $\vec{a} = -2\,\vec{i} - 10\,\vec{j}$ m/s$^2$. Find (a) $\vec{r}(t)$, (b) the time when the particle crosses the $x$-axis again, (c) the maximum height reached.
\end{enumerate}
\end{exercice}

\begin{corrige}[Exercises 1--3]
\begin{enumerate}
\item \textbf{Exercise 1:}
\begin{enumerate}
\item $\vec{v} = 4t\,\vec{i} + (3t^2-4)\,\vec{j}$, $\vec{a} = 4\,\vec{i} + 6t\,\vec{j}$.
\item At $t=1$: $\vec{v} = 4\,\vec{i} -1\,\vec{j}$, speed $= \sqrt{16+1} = \sqrt{17}$ m/s.
\item Parallel to $\vec{i}+2\vec{j} \Rightarrow \frac{4t}{1} = \frac{3t^2-4}{2} \Rightarrow 8t = 3t^2-4 \Rightarrow 3t^2-8t-4=0 \Rightarrow (3t+2)(t-2)=0 \Rightarrow t=2$ s ($t\ge0$).
\end{enumerate}
\item \textbf{Exercise 2:}
\begin{enumerate}
\item $x^2+y^2 = 25\cos^2 t + 25\sin^2 t = 25$, circle radius 5 centred at origin.
\item $\vec{v} = -5\sin t\,\vec{i} + 5\cos t\,\vec{j}$, $\vec{a} = -5\cos t\,\vec{i} -5\sin t\,\vec{j} = -\vec{r}$.
\item $\vec{v}\cdot\vec{r} = -25\sin t\cos t + 25\cos t\sin t = 0 \Rightarrow \vec{v} \perp \vec{r}$ (velocity tangent to circle). $\vec{a}\cdot\vec{v} = 25\cos t\sin t -25\sin t\cos t = 0 \Rightarrow \vec{a} \perp \vec{v}$ (acceleration radial, no tangential component, constant speed).
\end{enumerate}
\item \textbf{Exercise 3:}
\begin{enumerate}
\item $\vec{v}(t) = \vec{u} + \vec{a}t = (10-2t)\,\vec{i} + (20-10t)\,\vec{j}$.
$\vec{r}(t) = \int \vec{v}\,dt = (10t - t^2)\,\vec{i} + (20t - 5t^2)\,\vec{j}$.
\item Cross $x$-axis $\Rightarrow y=0 \Rightarrow 20t-5t^2=0 \Rightarrow t(4-t)=0 \Rightarrow t=4$ s (excluding $t=0$).
\item Max height when $v_y=0 \Rightarrow 20-10t=0 \Rightarrow t=2$ s. $y_{\max} = 20(2)-5(4) = 20$ m.
\end{enumerate}
\end{enumerate}
\end{corrige}

\coursec{From Acceleration back to Velocity and Position: Integration in Kinematics}

In the previous section we travelled \emph{down} the kinematic chain: starting from the position vector $\vec{r}(t)$ we differentiated once to obtain the velocity $\vec{v}(t) = \dot{\vec{r}}(t)$, and once more to obtain the acceleration $\vec{a}(t) = \ddot{\vec{r}}(t)$. In real life, however, the quantity we most often \emph{know} is the acceleration, because acceleration is produced by forces (Newton's second law, $\vec{F} = m\vec{a}$). A driver on the Yaound\'e--Douala highway who presses the accelerator pedal is controlling the \emph{acceleration} of the car; the speed and the position are consequences. This section is therefore devoted to the \emph{inverse problem}: \cle{given the acceleration, recover the velocity and the position by integration}.

\courssub{The Inverse Problem: Integration as the Reverse of Differentiation}

\begin{popfigure}
\begin{center}
\begin{tikzpicture}[
  box/.style={draw=popblue, thick, rounded corners, fill=popblueL, minimum width=2.4cm, minimum height=1.1cm, font=\bfseries},
  arr/.style={-{Stealth[length=3mm]}, thick, popdark},
  lbl/.style={font=\small, popdark, align=center}
]
\node[box] (r) at (0,0) {$\vec{r}(t)$};
\node[box] (v) at (5.2,0) {$\vec{v}(t)$};
\node[box] (a) at (10.4,0) {$\vec{a}(t)$};
\draw[arr] (r) -- (v) node[lbl, midway, above=2pt] {differentiate};
\draw[arr] (v) -- (a) node[lbl, midway, above=2pt] {differentiate};
\draw[arr, poporange] (a.south) .. controls (7.8,-1.6) .. (v.south) node[lbl, midway, below=2pt] {integrate $+$ initial velocity};
\draw[arr, poporange] (v.south) .. controls (2.6,-1.6) .. (r.south) node[lbl, midway, below=2pt] {integrate $+$ initial position};
\end{tikzpicture}
\end{center}
\legende{The kinematic chain. Differentiation goes from position to acceleration; integration, together with initial conditions, takes us back.}
\end{popfigure}

Since $\dfrac{d\vec{v}}{dt} = \vec{a}$, the velocity is an \emph{antiderivative} of the acceleration, and since $\dfrac{d\vec{r}}{dt} = \vec{v}$, the position is an antiderivative of the velocity. Integration is performed \cle{component by component}: if $\vec{a}(t) = a_x(t)\,\vec{i} + a_y(t)\,\vec{j}$, then
\[
\vec{v}(t) = \left(\int a_x(t)\,dt\right)\vec{i} + \left(\int a_y(t)\,dt\right)\vec{j}, \qquad
\vec{r}(t) = \left(\int v_x(t)\,dt\right)\vec{i} + \left(\int v_y(t)\,dt\right)\vec{j}.
\]
This is exactly the reverse of the component-wise differentiation of the previous section: the unit vectors $\vec{i}$ and $\vec{j}$ are constant, so they behave as ordinary constants under integration.

\begin{propriete}[Fundamental integration formulas]
If a particle moves in a plane with acceleration $\vec{a}(t)$, velocity $\vec{v}(t)$ and position $\vec{r}(t)$, then
\[
\vec{v}(t) = \vec{u} + \int_{0}^{t}\vec{a}(\tau)\,d\tau, \qquad
\vec{r}(t) = \vec{r}_0 + \int_{0}^{t}\vec{v}(\tau)\,d\tau,
\]
where $\vec{u} = \vec{v}(0)$ is the \cle{initial velocity} and $\vec{r}_0 = \vec{r}(0)$ is the \cle{initial position}. Equivalently, the change of velocity over $[t_1,t_2]$ is $\vec{v}(t_2)-\vec{v}(t_1) = \displaystyle\int_{t_1}^{t_2}\vec{a}\,dt$ and the displacement is $\vec{r}(t_2)-\vec{r}(t_1) = \displaystyle\int_{t_1}^{t_2}\vec{v}\,dt$.
\end{propriete}

\textbf{Justification.} By the Fundamental Theorem of Calculus, if $\dot{\vec{v}} = \vec{a}$ then $\displaystyle\int_0^t \vec{a}(\tau)\,d\tau = \vec{v}(t) - \vec{v}(0) = \vec{v}(t) - \vec{u}$, which gives the first formula. Applying the same argument to $\dot{\vec{r}} = \vec{v}$ gives the second. \hfill$\blacksquare$

\begin{attention}[The constants of integration]
Each integration introduces a \textbf{constant vector} $\vec{C}$. If you write $\vec{v}(t) = \displaystyle\int \vec{a}(t)\,dt$ and forget the constant, your answer is wrong \emph{even if the integration itself is perfect}. Forgetting initial conditions is the most common source of lost marks in GCE kinematics questions. Always determine each constant \emph{immediately} after each integration, before moving on.
\end{attention}

\begin{methode}[From $\vec{a}(t)$ to $\vec{r}(t)$]
\begin{enumerate}
\item Integrate $\vec{a}(t)$ component-wise to get $\vec{v}(t) + \vec{C}_1$.
\item Use the initial velocity $\vec{v}(0) = \vec{u}$ to fix $\vec{C}_1$.
\item Integrate the (now complete) $\vec{v}(t)$ component-wise to get $\vec{r}(t) + \vec{C}_2$.
\item Use the initial position $\vec{r}(0) = \vec{r}_0$ to fix $\vec{C}_2$.
\item Answer the question asked: position, velocity, speed $|\vec{v}|$, or displacement at a given time.
\end{enumerate}
\end{methode}

\courssub{Special Case: Constant Acceleration and the Vector suvat Equations}

Suppose $\vec{a}$ is \emph{constant}. Then, with $\vec{v}(0) = \vec{u}$ and $\vec{r}(0) = \vec{r}_0$:
\begin{align*}
\vec{v}(t) &= \vec{u} + \int_0^t \vec{a}\,d\tau = \vec{u} + \vec{a}\,t,\\[2pt]
\vec{r}(t) &= \vec{r}_0 + \int_0^t \left(\vec{u} + \vec{a}\tau\right)d\tau = \vec{r}_0 + \vec{u}\,t + \tfrac{1}{2}\vec{a}\,t^2.
\end{align*}

\begin{propriete}[Vector suvat equations for constant acceleration]
For a particle moving with \textbf{constant} acceleration $\vec{a}$,
\[
\boxed{\;\vec{v} = \vec{u} + \vec{a}t\;} \qquad\text{and}\qquad \boxed{\;\vec{r} = \vec{r}_0 + \vec{u}t + \tfrac{1}{2}\vec{a}t^2\;}
\]
These are the vector forms of the familiar suvat equations; taking the $\vec{i}$ and $\vec{j}$ components separately reproduces $v = u + at$ and $s = ut + \tfrac12 at^2$ in each direction. \emph{The derivation by integration, as above, is examinable.}
\end{propriete}

\begin{attention}[Constant acceleration only!]
The formulas $\vec{v} = \vec{u} + \vec{a}t$ and $\vec{r} = \vec{r}_0 + \vec{u}t + \tfrac12\vec{a}t^2$ are valid \textbf{only when $\vec{a}$ does not depend on time}. If $\vec{a} = 6t\,\vec{i} - 2\,\vec{j}$, for example, you \emph{must} integrate; substituting into suvat is a serious error.
\end{attention}

\courssub{Worked Examples with Variable Acceleration}

\begin{exoresolu}[Integrating a time-dependent acceleration]
A particle $P$ moves in a plane. At time $t$ seconds ($t \geq 0$) its acceleration is
\[
\vec{a}(t) = 6t\,\vec{i} - 2\,\vec{j} \quad (\mathrm{m\,s^{-2}}).
\]
Initially $P$ is at the origin $O$ and has velocity $\vec{u} = (2\vec{i} + \vec{j})\ \mathrm{m\,s^{-1}}$.
\begin{enumerate}
\item[(a)] Find the velocity and the position of $P$ at time $t$.
\item[(b)] Find the speed of $P$ when $t = 3$.
\item[(c)] Show that $P$ crosses the $x$-axis again and find the time at which this happens.
\end{enumerate}
\tcblower
\textbf{Solution.}
(a) \emph{Velocity.} Integrate $\vec{a}$ component-wise:
\[
\vec{v}(t) = \int \left(6t\,\vec{i} - 2\,\vec{j}\right)dt = 3t^2\,\vec{i} - 2t\,\vec{j} + \vec{C}_1.
\]
At $t = 0$: $\vec{v}(0) = \vec{C}_1 = 2\vec{i} + \vec{j}$. Hence
\[
\vec{v}(t) = (3t^2 + 2)\,\vec{i} + (1 - 2t)\,\vec{j} \quad (\mathrm{m\,s^{-1}}).
\]
\emph{Position.} Integrate again:
\[
\vec{r}(t) = \int \vec{v}(t)\,dt = \left(t^3 + 2t\right)\vec{i} + \left(t - t^2\right)\vec{j} + \vec{C}_2.
\]
At $t = 0$: $\vec{r}(0) = \vec{C}_2 = \vec{0}$ (the particle starts at $O$). Hence
\[
\vec{r}(t) = (t^3 + 2t)\,\vec{i} + (t - t^2)\,\vec{j} \quad (\mathrm{m}).
\]

(b) At $t = 3$: $\vec{v}(3) = (27 + 2)\vec{i} + (1 - 6)\vec{j} = 29\vec{i} - 5\vec{j}$, so the speed is
\[
|\vec{v}(3)| = \sqrt{29^2 + (-5)^2} = \sqrt{856} \approx 29.3\ \mathrm{m\,s^{-1}}.
\]

(c) $P$ is on the $x$-axis when the $\vec{j}$-component of $\vec{r}$ is zero:
\[
t - t^2 = t(1 - t) = 0 \implies t = 0 \ \text{or}\ t = 1.
\]
So $P$ starts on the axis and crosses it again at $t = 1\ \mathrm{s}$, at the point $\vec{r}(1) = 3\vec{i}$, i.e. $3\ \mathrm{m}$ from $O$ along the $x$-axis. \hfill$\blacksquare$
\end{exoresolu}

\begin{exoresolu}[When does the particle return to the origin?]
A particle leaves the origin $O$ with velocity $(4\vec{i} + 2\vec{j})\ \mathrm{m\,s^{-1}}$ and moves with acceleration $\vec{a}(t) = (-2\vec{i} - \vec{j})\ \mathrm{m\,s^{-2}}$ (constant). Find the time at which it returns to $O$, and its velocity at that instant.
\tcblower
\textbf{Solution.} Since $\vec{a}$ is constant, $\vec{v} = \vec{u} + \vec{a}t = (4 - 2t)\vec{i} + (2 - t)\vec{j}$, and
\[
\vec{r}(t) = \vec{u}t + \tfrac12\vec{a}t^2 = \left(4t - t^2\right)\vec{i} + \left(2t - \tfrac12 t^2\right)\vec{j}.
\]
Return to $O$ requires \emph{both} components to vanish simultaneously:
\[
t(4 - t) = 0 \quad\text{and}\quad t\left(2 - \tfrac12 t\right) = 0.
\]
The first equation gives $t = 0$ or $t = 4$; the second gives $t = 0$ or $t = 4$. The common non-zero solution is $t = 4\ \mathrm{s}$. Then $\vec{v}(4) = (4 - 8)\vec{i} + (2 - 4)\vec{j} = -4\vec{i} - 2\vec{j}\ \mathrm{m\,s^{-1}}$: the particle returns with exactly the opposite velocity, as symmetry suggests. \hfill$\blacksquare$
\end{exoresolu}

\begin{attention}[Both components must vanish together]
To find when a particle passes through the origin (or any given point), set \emph{each} component of $\vec{r}(t)$ equal to the required value and look for a value of $t$ satisfying \textbf{all} the equations at once. A time that zeros only one component corresponds to crossing an axis, not to passing through $O$.
\end{attention}

\courssub{Displacement from Velocity: Definite Integrals, and Distance versus Displacement}

Often the velocity is given directly and we only need the \emph{change} of position between two times. No constant of integration is then required, because we use a definite integral:
\[
\text{displacement from } t_1 \text{ to } t_2:\quad \vec{r}(t_2) - \vec{r}(t_1) = \int_{t_1}^{t_2} \vec{v}(t)\,dt.
\]

\begin{popfigure}
\begin{center}
\begin{tikzpicture}
\begin{axis}[
  width=11cm, height=6.2cm,
  axis lines=middle,
  xlabel={$t\ (\mathrm{s})$}, ylabel={$v\ (\mathrm{m\,s^{-1}})$},
  xmin=0, xmax=6.4, ymin=-1.2, ymax=4.6,
  xtick={1,2,3,4,5,6}, ytick={1,2,3,4},
  domain=0:6.2, samples=100, clip=false
]
\addplot[domain=2:5, draw=none, fill=poporange, opacity=0.35] {1.5 + 1.2*sin(deg(x-0.6)) + 0.15*x} \closedcycle;
\addplot[popblue, very thick] {1.5 + 1.2*sin(deg(x-0.6)) + 0.15*x};
\draw[dashed, popdark] (axis cs:2,0) -- (axis cs:2,2.98);
\draw[dashed, popdark] (axis cs:5,0) -- (axis cs:5,1.11);
\node[popdark, font=\small] at (axis cs:2,-0.55) {$t_1$};
\node[popdark, font=\small] at (axis cs:5,-0.55) {$t_2$};
\node[popdark, font=\small, align=center] at (axis cs:3.5,0.7) {area $=$ displacement $\displaystyle\int_{t_1}^{t_2} v\,dt$};
\node[popblue, font=\small] at (axis cs:5.7,2.1) {$v(t)$};
\end{axis}
\end{tikzpicture}
\end{center}
\legende{The displacement between $t_1$ and $t_2$ equals the (signed) area under the velocity--time curve.}
\end{popfigure}

\begin{attention}[Displacement is not distance]
The displacement is $\displaystyle\int_{t_1}^{t_2} \vec{v}\,dt$: regions where a component of $\vec{v}$ is negative \emph{subtract} from the integral. The \cle{total distance travelled} requires integrating the \emph{speed}:
\[
\text{distance} = \int_{t_1}^{t_2} |\vec{v}(t)|\,dt,
\]
splitting the interval wherever the relevant component of $\vec{v}$ changes sign. For one-dimensional motion with velocity $v(t)$, distance $= \displaystyle\int_{t_1}^{t_2} |v(t)|\,dt$: integrate $v$ where $v>0$ and $-v$ where $v<0$, then add.
\end{attention}

\begin{exemplebox}[Displacement versus distance]
A particle moves along a straight line with velocity $v(t) = 3t - 6\ \mathrm{m\,s^{-1}}$. Find (a) the displacement and (b) the distance travelled between $t = 0$ and $t = 4\ \mathrm{s}$.

\textbf{Solution.} (a) Displacement $= \displaystyle\int_0^4 (3t-6)\,dt = \left[\tfrac32 t^2 - 6t\right]_0^4 = (24 - 24) - 0 = 0\ \mathrm{m}$: the particle ends where it started.

(b) $v$ changes sign at $t = 2$. On $[0,2]$, $v<0$; on $[2,4]$, $v>0$.
\[
\text{distance} = \int_0^2 (6 - 3t)\,dt + \int_2^4 (3t - 6)\,dt = \left[6t - \tfrac32 t^2\right]_0^2 + \left[\tfrac32 t^2 - 6t\right]_2^4 = 6 + 6 = 12\ \mathrm{m}.
\]
The two answers differ completely: zero displacement, but $12\ \mathrm{m}$ actually travelled.
\end{exemplebox}

\courssub{GCE-Style Synthesis}

\begin{exoresolu}[Full synthesis problem]
A particle $P$ leaves the origin $O$ at time $t = 0$ with velocity $(3\vec{i} + 5\vec{j})\ \mathrm{m\,s^{-1}}$. At time $t$ seconds its acceleration is
\[
\vec{a}(t) = \left(2t\,\vec{i} - 4\,\vec{j}\right)\ \mathrm{m\,s^{-2}}.
\]
\begin{enumerate}
\item[(a)] Find the position vector of $P$ at time $t$.
\item[(b)] Find the speed of $P$ when $t = 2$.
\item[(c)] Find the displacement of $P$ between $t = 1$ and $t = 3$.
\end{enumerate}
\tcblower
\textbf{Solution.}
(a) Integrate the acceleration:
\[
\vec{v}(t) = t^2\,\vec{i} - 4t\,\vec{j} + \vec{C}_1; \qquad \vec{v}(0) = \vec{C}_1 = 3\vec{i} + 5\vec{j},
\]
so $\vec{v}(t) = (t^2 + 3)\vec{i} + (5 - 4t)\vec{j}\ \mathrm{m\,s^{-1}}$. Integrate again:
\[
\vec{r}(t) = \left(\tfrac13 t^3 + 3t\right)\vec{i} + \left(5t - 2t^2\right)\vec{j} + \vec{C}_2; \qquad \vec{r}(0) = \vec{C}_2 = \vec{0}.
\]
\[
\therefore\ \vec{r}(t) = \left(\tfrac13 t^3 + 3t\right)\vec{i} + \left(5t - 2t^2\right)\vec{j}\ \mathrm{m}.
\]

(b) $\vec{v}(2) = (4+3)\vec{i} + (5-8)\vec{j} = 7\vec{i} - 3\vec{j}$, so the speed is
\[
|\vec{v}(2)| = \sqrt{7^2 + 3^2} = \sqrt{58} \approx 7.62\ \mathrm{m\,s^{-1}}.
\]

(c) Using the definite integral (no constants needed):
\[
\vec{r}(3) - \vec{r}(1) = \int_1^3 \vec{v}(t)\,dt = \left[\left(\tfrac13 t^3 + 3t\right)\vec{i} + \left(5t - 2t^2\right)\vec{j}\right]_1^3.
\]
At $t = 3$: $(9 + 9)\vec{i} + (15 - 18)\vec{j} = 18\vec{i} - 3\vec{j}$. At $t = 1$: $\tfrac{10}{3}\vec{i} + 3\vec{j}$.
\[
\text{Displacement} = \left(18 - \tfrac{10}{3}\right)\vec{i} + (-3 - 3)\vec{j} = \tfrac{44}{3}\vec{i} - 6\vec{j}\ \mathrm{m}.
\]
\hfill$\blacksquare$
\end{exoresolu}

\begin{methode}[Exam technique for integration problems]
\begin{enumerate}
\item Write vectors consistently in $\vec{i},\vec{j}$ (or column) form throughout the whole solution --- never mix the two.
\item Keep the units ($\mathrm{m\,s^{-2}}$, $\mathrm{m\,s^{-1}}$, $\mathrm{m}$) at every line; a unit error often reveals a missing integration step.
\item Fix each constant of integration \emph{as soon as it appears}, quoting the initial condition explicitly: ``when $t=0$, $\vec{v} = \ldots$''.
\item For ``speed'', remember to take the modulus $|\vec{v}| = \sqrt{v_x^2 + v_y^2}$; for ``velocity'', leave the vector.
\item Check your final $\vec{r}(t)$ by differentiating twice: you must recover the given $\vec{a}(t)$.
\end{enumerate}
\end{methode}

\begin{aretenir}[Key points]
\begin{itemize}
\item $\vec{v}(t) = \vec{u} + \displaystyle\int_0^t \vec{a}\,d\tau$ and $\vec{r}(t) = \vec{r}_0 + \displaystyle\int_0^t \vec{v}\,d\tau$; integration is done component by component.
\item Each integration produces a constant vector, fixed by an initial condition.
\item Constant acceleration gives the vector suvat equations $\vec{v} = \vec{u} + \vec{a}t$, $\vec{r} = \vec{r}_0 + \vec{u}t + \tfrac12\vec{a}t^2$ --- valid \emph{only} for constant $\vec{a}$.
\item Displacement $= \displaystyle\int_{t_1}^{t_2}\vec{v}\,dt$ (signed area under the $v$--$t$ curve); distance travelled $= \displaystyle\int_{t_1}^{t_2}|\vec{v}|\,dt$, splitting at sign changes.
\end{itemize}
\end{aretenir}

\begin{exercice}[Practice]
\begin{enumerate}
\item A particle starts from rest at the point with position vector $(2\vec{i} - \vec{j})\ \mathrm{m}$ and moves with acceleration $\vec{a}(t) = (4\vec{i} + 6t\,\vec{j})\ \mathrm{m\,s^{-2}}$. Find its velocity and position at time $t$, and its distance from the origin when $t = 2$.
\item A particle moves with velocity $\vec{v}(t) = (2t - 4)\vec{i} + 3\vec{j}\ \mathrm{m\,s^{-1}}$. Find (a) its displacement between $t = 0$ and $t = 5$; (b) the total distance travelled in the $x$-direction over the same interval.
\item A particle leaves $O$ with velocity $\vec{u} = (\vec{i} + 4\vec{j})\ \mathrm{m\,s^{-1}}$ and has constant acceleration $\vec{a} = (2\vec{i} - 2\vec{j})\ \mathrm{m\,s^{-2}}$. Show that it crosses the $x$-axis again at $t = 4\ \mathrm{s}$ and find its speed at that instant.
\end{enumerate}
\end{exercice}

\begin{corrige}[Exercise 1]
$\vec{v}(t) = 4t\,\vec{i} + 3t^2\,\vec{j}$ (constant $\vec{0}$ since it starts from rest). $\vec{r}(t) = (2 + 2t^2)\vec{i} + (t^3 - 1)\vec{j}$ (constant $(2\vec{i} - \vec{j})$). At $t = 2$: $\vec{r} = 10\vec{i} + 7\vec{j}$, so distance from $O$ is $\sqrt{100 + 49} = \sqrt{149} \approx 12.2\ \mathrm{m}$.
\end{corrige}

\begin{corrige}[Exercise 2]
(a) Displacement $= \displaystyle\int_0^5 \vec{v}\,dt = \left[t^2 - 4t\right]_0^5\vec{i} + \left[3t\right]_0^5\vec{j} = 5\vec{i} + 15\vec{j}\ \mathrm{m}$.
(b) $v_x = 2t - 4$ changes sign at $t = 2$: distance in the $x$-direction $= \displaystyle\int_0^2 (4 - 2t)\,dt + \int_2^5 (2t - 4)\,dt = 4 + 9 = 13\ \mathrm{m}$.
\end{corrige}

\begin{corrige}[Exercise 3]
$\vec{r}(t) = \vec{u}t + \tfrac12\vec{a}t^2 = (t + t^2)\vec{i} + (4t - t^2)\vec{j}$. On the $x$-axis the $\vec{j}$-component vanishes: $t(4 - t) = 0$ gives $t = 0$ (start) or $t = 4\ \mathrm{s}$, as required. Then $\vec{v}(4) = \vec{u} + 4\vec{a} = (\vec{i} + 4\vec{j}) + (8\vec{i} - 8\vec{j}) = 9\vec{i} - 4\vec{j}$, and the speed is $\sqrt{81 + 16} = \sqrt{97} \approx 9.85\ \mathrm{m\,s^{-1}}$.
\end{corrige}

\begin{savaistu}[Why this section matters for what comes next]
Projectile motion --- the subject of the next sections --- is nothing more than the constant-acceleration case of what you have just learned, with $\vec{a} = -g\,\vec{j}$. Every formula for the range, the greatest height and the time of flight of a projectile is obtained by exactly the two integrations you have practised here. Master this section, and the rest of the chapter becomes a routine application.
\end{savaistu}

\coursec{Projectile Motion: Definition, Model and Equations}

In the previous section we saw how to reconstruct the velocity and position of a particle from its acceleration by integration. We now apply those ideas to one of the most important and beautiful problems in mechanics: the motion of a \textbf{projectile}. A projectile is any object that is launched into the air and then moves freely under the influence of gravity alone. Think of a football kicked from the penalty spot, a stone thrown from a hilltop, or a cannonball fired from the walls of the Foumban palace. Once the object leaves the foot, the hand or the barrel, the only force acting on it (in our model) is its weight. This section builds the complete mathematical model for that motion.

\courssub{What is a Projectile?}

\begin{definition}[Projectile]
A \textbf{projectile} is a particle that is projected (launched) with an initial velocity and thereafter moves freely under the action of gravity alone, with no other forces (such as air resistance or thrust) acting on it.
\end{definition}

The key phrase is \emph{``moves freely under the action of gravity alone''}. The moment the projectile loses contact with the launcher, its acceleration is exactly the acceleration due to gravity, directed vertically downwards. This is true whether the projectile is moving upward, downward, or is instantaneously at the top of its flight.

\courssub{Modelling Assumptions}

To obtain a tractable mathematical model we make the following standard assumptions. In an examination you should be ready to state them.

\begin{aretenir}[Standard Projectile Model Assumptions]
\begin{enumerate}
    \item The projectile is treated as a \textbf{particle} (a point mass); its size, shape and rotation are ignored.
    \item The acceleration due to gravity, $g$, is \textbf{constant} in magnitude and direction. Near the Earth's surface we take $g = 9.8\ \mathrm{m\,s^{-2}}$ unless the question explicitly says ``take $g = 10\ \mathrm{m\,s^{-2}}$''.
    \item \textbf{No air resistance} (or any other resistive force). The only force is the weight $mg$ downwards.
    \item The motion takes place over distances small enough that the curvature of the Earth and the variation of $g$ with height can be neglected.
\end{enumerate}
\end{aretenir}

\begin{attention}[Assumptions in Exam Answers]
When a question asks ``State the modelling assumptions you have made'', you must list the ones relevant to the problem. The most common required answers are: ``The object is modelled as a particle'' and ``Air resistance is neglected''. Do not write ``gravity is constant'' unless the question specifically asks about the model for gravity.
\end{attention}

\courssub{Setting Up the Coordinate System}

Because the acceleration is purely vertical, it is natural to choose axes that separate the horizontal and vertical directions.

\begin{methode}[Choosing Axes for Projectile Problems]
\begin{enumerate}
    \item Place the \textbf{origin} at the point of projection (or at ground level directly below it).
    \item Let the \textbf{$x$-axis} be horizontal (positive in the direction of the horizontal component of the initial velocity).
    \item Let the \textbf{$y$-axis} be vertical (positive \textbf{upwards}).
    \item Then the acceleration vector is $\vec{a} = (0,\ -g)$.
\end{enumerate}
\end{methode}

With this choice, the horizontal and vertical motions become completely independent --- a crucial simplification.

\begin{popfigure}
\begin{tikzpicture}[scale=0.9, >=Stealth]
    % Axes
    \draw[->, thick] (-0.5,0) -- (10,0) node[right] {$x$};
    \draw[->, thick] (0,-1) -- (0,5.5) node[above] {$y$};
    % Parabolic trajectory y = 1.5x - 0.125x^2
    \draw[popblue, thick, domain=0:8.5, smooth, variable=\x] plot ({\x}, {1.5*\x - 0.125*\x*\x});
    % Initial velocity vector and components
    \draw[poporange, very thick, ->] (0,0) -- (2,1.5) node[midway, above left, poporange] {$\vec{u}$};
    \draw[poporange, dashed] (2,0) -- (2,1.5);
    \draw[poporange, dashed] (0,1.5) -- (2,1.5);
    \draw[poporange, thick, ->] (0,0) -- (2,0) node[midway, below, poporange] {$u\cos\theta$};
    \draw[poporange, thick, ->] (0,0) -- (0,1.5) node[midway, left, poporange] {$u\sin\theta$};
    % Angle
    \draw[popdark] (0.6,0) arc (0:36.87:0.6) node[midway, right, popdark] {$\theta$};
    % Labels
    \node[below left] at (0,0) {$O$};
    % Velocity at point P = (7, 4.375) on the curve, slope -0.25
    \fill[popgreen] (7,4.375) circle (2pt) node[above left, popgreen] {$P$};
    \draw[popgreen, very thick, ->] (7,4.375) -- (8.4,4.025) node[midway, above, popgreen] {$\vec{v}$};
    \draw[popgreen, dashed] (8.4,4.375) -- (8.4,4.025);
    \draw[popgreen, dashed] (7,4.375) -- (8.4,4.375);
    \draw[popgreen, thick, ->] (7,4.375) -- (8.4,4.375) node[midway, below, popgreen] {$v_x$};
    \draw[popgreen, thick, ->] (8.4,4.375) -- (8.4,4.025) node[midway, right, popgreen] {$v_y$};
    % Gravity vector
    \draw[poppink, very thick, ->] (9.3,4.5) -- (9.3,3.5) node[midway, right, poppink] {$\vec{g}$};
\end{tikzpicture}
\legende{Projectile launched from $O$ with speed $u$ at angle $\theta$. The horizontal velocity remains constant; the vertical velocity changes due to gravity.}
\end{popfigure}

\courssub{Initial Velocity Components}

The initial velocity $\vec{u}$ has magnitude $u$ (often denoted $V$ in GCE questions) and makes an angle $\theta$ (or $\alpha$) with the horizontal. Resolving into components:

\[
u_x = u\cos\theta, \qquad u_y = u\sin\theta.
\]

\begin{propriete}[Independence of Horizontal and Vertical Motions]
Because $\vec{a} = (0,\ -g)$, the horizontal acceleration is zero and the vertical acceleration is $-g$. Consequently:
\begin{itemize}
    \item The \textbf{horizontal motion} is \textbf{uniform} (constant velocity): $v_x = u\cos\theta$ for all $t$.
    \item The \textbf{vertical motion} is \textbf{uniformly accelerated} with acceleration $-g$: $v_y = u\sin\theta - gt$.
\end{itemize}
The two motions share only the common time variable $t$. This independence is the cornerstone of projectile analysis.
\end{propriete}

\courssub{Equations of Motion}

Integrating the acceleration (or using the constant-acceleration formulae) gives the parametric equations for position and velocity as functions of time $t$.

\begin{aretenir}[Parametric Equations of Projectile Motion]
\begin{align*}
\text{Horizontal:} \quad & x = (u\cos\theta)\,t, \qquad v_x = u\cos\theta. \\[4pt]
\text{Vertical:} \quad & y = (u\sin\theta)\,t - \frac{1}{2}gt^2, \qquad v_y = u\sin\theta - gt.
\end{align*}
The position vector is $\vec{r}(t) = \bigl((u\cos\theta)t,\ (u\sin\theta)t - \tfrac{1}{2}gt^2\bigr)$ and the velocity vector is $\vec{v}(t) = \bigl(u\cos\theta,\ u\sin\theta - gt\bigr)$.
\end{aretenir}

\begin{attention}[Sign of $g$]
Because we chose $y$ positive \textbf{upwards}, the acceleration due to gravity is $-g$. If you instead choose $y$ positive downwards, the acceleration becomes $+g$ and the initial vertical velocity becomes $-u\sin\theta$. Be consistent! The GCE standard is $y$ upwards, $a_y = -g$.
\end{attention}

\courssub{The Cartesian Equation of the Trajectory}

We can eliminate the parameter $t$ to obtain the path $y$ as a function of $x$. From $x = (u\cos\theta)t$ we have $t = \dfrac{x}{u\cos\theta}$. Substitute into the equation for $y$:

\[
y = (u\sin\theta)\left(\frac{x}{u\cos\theta}\right) - \frac{1}{2}g\left(\frac{x}{u\cos\theta}\right)^2
   = x\tan\theta - \frac{g x^2}{2u^2\cos^2\theta}.
\]

Using the identity $\sec^2\theta = 1 + \tan^2\theta$, this is sometimes written as

\[
y = x\tan\theta - \frac{g x^2}{2u^2}(1+\tan^2\theta).
\]

\begin{propriete}[Trajectory is a Parabola]
The Cartesian equation $y = x\tan\theta - \dfrac{g x^2}{2u^2\cos^2\theta}$ is a quadratic in $x$ with a negative coefficient for $x^2$. Hence the trajectory of a projectile (under the standard assumptions) is a \textbf{parabola} with a vertical axis, opening downwards.
\end{propriete}

\begin{exemplebox}[Finding the Trajectory Equation]
A particle is projected from ground level with a speed of $\SI{20}{m.s^{-1}}$ at an angle of $30^\circ$ to the horizontal. Taking $g = \SI{10}{m.s^{-2}}$, find the Cartesian equation of its trajectory.
\begin{align*}
u &= 20,\quad \theta = 30^\circ,\quad \tan30^\circ = \frac{1}{\sqrt{3}},\quad \cos30^\circ = \frac{\sqrt{3}}{2}. \\[4pt]
y &= x\tan30^\circ - \frac{g x^2}{2u^2\cos^2 30^\circ} \\
  &= \frac{x}{\sqrt{3}} - \frac{10\,x^2}{2 \times 400 \times \left(\frac{\sqrt{3}}{2}\right)^2} \\
  &= \frac{x}{\sqrt{3}} - \frac{10\,x^2}{800 \times \frac{3}{4}} \\
  &= \frac{x}{\sqrt{3}} - \frac{10\,x^2}{600} \\
  &= \frac{x}{\sqrt{3}} - \frac{x^2}{60}.
\end{align*}
The trajectory is $y = \dfrac{x}{\sqrt{3}} - \dfrac{x^2}{60}$ (with $x,y$ in metres).
\end{exemplebox}

\courssub{Velocity at Any Time}

The velocity vector at time $t$ is $\vec{v}(t) = (u\cos\theta,\ u\sin\theta - gt)$. Its magnitude (speed) and direction are often required.

\begin{aretenir}[Speed and Direction at Time $t$]
\begin{align*}
\text{Speed:} \quad & v = \sqrt{v_x^2 + v_y^2} = \sqrt{(u\cos\theta)^2 + (u\sin\theta - gt)^2}. \\[4pt]
\text{Angle $\phi$ to the horizontal:} \quad & \tan\phi = \frac{v_y}{v_x} = \frac{u\sin\theta - gt}{u\cos\theta}.
\end{align*}
If $\phi > 0$ the velocity points above the horizontal; if $\phi < 0$ it points below.
\end{aretenir}

\begin{attention}[Velocity at the Highest Point]
At the maximum height, the vertical velocity $v_y = 0$. The velocity is then purely horizontal: $\vec{v} = (u\cos\theta,\ 0)$. The speed is minimum there and equals $u\cos\theta$. \textbf{The acceleration is still $g$ downwards} --- a common mistake is to think the acceleration becomes zero at the top.
\end{attention}

\courssub{Worked Example: Complete Analysis}

\begin{exoresolu}[Projectile from a Cliff]
A stone is projected from the edge of a vertical cliff with a speed of $\SI{25}{m.s^{-1}}$ at an angle of $40^\circ$ above the horizontal. The cliff is $\SI{50}{m}$ high. Take $g = \SI{9.8}{m.s^{-2}}$.
\begin{enumerate}
    \item Find the maximum height reached above the cliff top.
    \item Find the time taken to reach the sea.
    \item Find the horizontal distance from the cliff base where the stone hits the sea.
    \item Find the speed and direction of the stone just before impact.
\end{enumerate}
\tcblower
\textbf{Step 1: Resolve the initial velocity.}
$u = 25$, $\theta = 40^\circ$.
$u_x = 25\cos40^\circ \approx 19.15\ \mathrm{m\,s^{-1}}$,
$u_y = 25\sin40^\circ \approx 16.07\ \mathrm{m\,s^{-1}}$.

\textbf{Step 2: Maximum height above cliff.}
At max height, $v_y = 0$. Using $v_y^2 = u_y^2 - 2g h_{\text{max}}$:
$0 = (16.07)^2 - 2(9.8)h_{\text{max}} \implies h_{\text{max}} = \dfrac{258.2}{19.6} \approx \SI{13.2}{m}$ above the cliff.

\textbf{Step 3: Time to reach the sea.}
Take $y=0$ at cliff top, so sea level is $y = -50$.
$y = u_y t - \frac{1}{2}gt^2 = -50$.
$16.07 t - 4.9 t^2 = -50 \implies 4.9 t^2 - 16.07 t - 50 = 0$.
$t = \dfrac{16.07 + \sqrt{16.07^2 + 4\times4.9\times50}}{2\times4.9}
   = \dfrac{16.07 + \sqrt{258.2 + 980}}{9.8}
   = \dfrac{16.07 + 35.19}{9.8} \approx \SI{5.23}{s}$.
(We discard the negative root, since time must be positive.)

\textbf{Step 4: Horizontal distance from cliff base.}
$x = u_x t = 19.15 \times 5.23 \approx \SI{100}{m}$ (3 s.f.: $\SI{100.2}{m}$).

\textbf{Step 5: Velocity at impact.}
$v_x = u_x = 19.15\ \mathrm{m\,s^{-1}}$.
$v_y = u_y - gt = 16.07 - 9.8(5.23) = 16.07 - 51.25 = -35.18\ \mathrm{m\,s^{-1}}$.
Speed $v = \sqrt{19.15^2 + (-35.18)^2} = \sqrt{366.7 + 1237.6} = \sqrt{1604.3} \approx \SI{40.1}{m.s^{-1}}$.
Angle below horizontal: $\tan\phi = \dfrac{|v_y|}{v_x} = \dfrac{35.18}{19.15} \approx 1.837 \implies \phi \approx 61.4^\circ$ below the horizontal.
\end{exoresolu}

\courssub{Complementary Angles and Range}

For a given projection speed $u$, two different angles of projection that sum to $90^\circ$ (complementary angles) produce the same horizontal range (on level ground). This is a classic GCE result.

\begin{propriete}[Range on Level Ground]
If the projectile lands at the same vertical level from which it was launched ($y=0$), the time of flight $T$ and the range $R$ are:
\[
T = \frac{2u\sin\theta}{g}, \qquad R = \frac{u^2\sin2\theta}{g}.
\]
Since $\sin(180^\circ-2\theta) = \sin2\theta$, the angles $\theta$ and $90^\circ-\theta$ give the same range. The maximum range is $R_{\text{max}} = \dfrac{u^2}{g}$, achieved at $\theta = 45^\circ$.
\end{propriete}

\begin{popfigure}
\begin{tikzpicture}
\begin{axis}[
    width=0.9\linewidth, height=7cm,
    axis lines=middle,
    xlabel={$x$ (m)}, ylabel={$y$ (m)},
    xmin=0, xmax=45, ymin=0, ymax=12,
    xtick={0,10,20,30,40}, ytick={0,5,10},
    grid=major, grid style={popline},
    legend style={at={(0.5,1.02)}, anchor=south, legend columns=2, draw=popmuted},
    clip=false
]
\addplot[popblue, thick, domain=0:40, samples=100] {x*tan(30) - 9.8*x^2/(2*20^2*cos(30)^2)};
\addlegendentry{$\theta = 30^\circ$}
\addplot[poporange, thick, domain=0:40, samples=100] {x*tan(60) - 9.8*x^2/(2*20^2*cos(60)^2)};
\addlegendentry{$\theta = 60^\circ$}
\addplot[popgreen, dashed, thick, domain=0:40.8] {x*tan(45) - 9.8*x^2/(2*20^2*cos(45)^2)};
\addlegendentry{$\theta = 45^\circ$ (max range)}
\end{axis}
\end{tikzpicture}
\legende{Trajectories for $u = \SI{20}{m.s^{-1}}$ at $30^\circ$, $45^\circ$ and $60^\circ$ ($g=9.8$). The $30^\circ$ and $60^\circ$ trajectories have the same range; $45^\circ$ gives the maximum range.}
\end{popfigure}

\courssub{Key Points to Remember}

\begin{aretenir}[Summary: Projectile Motion Essentials]
\begin{enumerate}
    \item \textbf{Always resolve} the initial velocity into horizontal and vertical components first.
    \item \textbf{Horizontal motion:} constant velocity $u\cos\theta$; $x = (u\cos\theta)t$.
    \item \textbf{Vertical motion:} constant acceleration $-g$; $y = (u\sin\theta)t - \frac{1}{2}gt^2$, $v_y = u\sin\theta - gt$.
    \item \textbf{Trajectory:} $y = x\tan\theta - \dfrac{g x^2}{2u^2\cos^2\theta}$ (a parabola).
    \item \textbf{Range on level ground:} $R = \dfrac{u^2\sin2\theta}{g}$; maximum at $\theta=45^\circ$.
    \item \textbf{At the highest point:} $v_y=0$, speed $= u\cos\theta$, acceleration $= g$ downwards.
    \item \textbf{Check the value of $g$} given in the question ($9.8$ or $10$) and the required accuracy (usually 3 significant figures or 1 decimal place).
\end{enumerate}
\end{aretenir}

\begin{savaistu}[Historical Note]
The parabolic trajectory of projectiles was first demonstrated mathematically by Galileo Galilei in his \emph{Discorsi e dimostrazioni matematiche intorno a due nuove scienze} (1638). Before Galileo, it was commonly believed (following Aristotle) that a projectile moved in a straight line until its ``impetus'' was exhausted, then fell vertically. Galileo's insight --- separating the motion into independent horizontal and vertical components --- is exactly the method we use today. In Cameroon, traditional hunters intuitively understood the need to aim higher for distant targets, effectively compensating for the parabolic drop long before the mathematics was formalised.
\end{savaistu}

\begin{exercice}[Practice Problems]
\begin{enumerate}
    \item A ball is thrown from ground level with a speed of $\SI{15}{m.s^{-1}}$ at $50^\circ$ to the horizontal. ($g=9.8$)
    \begin{enumerate}
        \item Find the maximum height reached.
        \item Find the range on level ground.
        \item Find the velocity (magnitude and direction) after $1.2\ \mathrm{s}$.
    \end{enumerate}
    \item A projectile is fired from the top of an $\SI{80}{m}$ high tower with a velocity of $\SI{50}{m.s^{-1}}$ at $30^\circ$ above the horizontal. ($g=10$)
    \begin{enumerate}
        \item Find the time of flight.
        \item Find the horizontal distance from the foot of the tower where it lands.
        \item Find the angle of impact with the ground.
    \end{enumerate}
    \item Show that for a given speed $u$, the maximum range on a horizontal plane is $\dfrac{u^2}{g}$ and occurs at $\theta = 45^\circ$. For a fixed range $R < \dfrac{u^2}{g}$, show that there are two possible angles of projection and that they are complementary.
\end{enumerate}
\end{exercice}

\begin{corrige}[Exercise 1]
\begin{enumerate}
    \item $u=15$, $\theta=50^\circ$, $u_x=15\cos50^\circ\approx9.64$, $u_y=15\sin50^\circ\approx11.49$.
    \begin{enumerate}
        \item $H = \dfrac{u_y^2}{2g} = \dfrac{11.49^2}{19.6} \approx \SI{6.74}{m}$.
        \item $R = \dfrac{u^2\sin100^\circ}{g} = \dfrac{225\times0.9848}{9.8} \approx \SI{22.6}{m}$.
        \item At $t=1.2$: $v_x=9.64$, $v_y=11.49-9.8(1.2)= -0.27$. Speed $\approx\SI{9.64}{m.s^{-1}}$, angle $\approx -1.6^\circ$ (just below horizontal).
    \end{enumerate}
    \item $u=50$, $\theta=30^\circ$, $u_x=25\sqrt{3}\approx43.3$, $u_y=25$. Tower height $80\ \mathrm{m}$ ($y=-80$).
    \begin{enumerate}
        \item $-80 = 25t - 5t^2 \implies 5t^2 -25t -80=0 \implies t^2-5t-16=0 \implies t = \dfrac{5+\sqrt{89}}{2} \approx \SI{7.22}{s}$.
        \item $x = 43.3 \times 7.22 \approx \SI{312.6}{m}$.
        \item $v_y = 25 - 10(7.22) = -47.2$. $\tan\phi = 47.2/43.3 \implies \phi \approx 47.5^\circ$ below horizontal.
    \end{enumerate}
    \item $R = \dfrac{u^2\sin2\theta}{g}$. Max when $\sin2\theta=1 \implies 2\theta=90^\circ \implies \theta=45^\circ$, $R_{\max}=u^2/g$. For $R < u^2/g$, $\sin2\theta = Rg/u^2$. Two solutions: $2\theta_1 = \alpha$, $2\theta_2 = 180^\circ-\alpha \implies \theta_1+\theta_2=90^\circ$.
\end{enumerate}
\end{corrige}

\coursec{Range, Maximum Height and Time of Flight}

In the previous section we modelled a projectile as a particle moving under gravity alone, and we established its equations of motion. With axes $Ox$ horizontal and $Oy$ vertically upwards, launch speed $u$ at angle $\theta$ above the horizontal from the origin, we found
\[
x(t)=u\cos\theta\; t, \qquad y(t)=u\sin\theta\; t-\tfrac{1}{2}gt^{2},
\]
\[
v_x=u\cos\theta, \qquad v_y=u\sin\theta-gt,
\]
together with the trajectory equation
\[
y=x\tan\theta-\frac{g\,x^{2}}{2u^{2}\cos^{2}\theta}.
\]
Everything in this section is a consequence of these four lines. Three questions dominate GCE questions: \emph{How long is the projectile in the air? How high does it go? How far does it go?} We answer each in turn, deriving every formula from the equations of motion rather than quoting results.

\courssub{Time of Flight on Level Ground}

\begin{definition}[Time of flight]
The \cle{time of flight} $T$ of a projectile is the total time between launch and landing. On level ground (launch and landing at the same height), it is the positive solution of $y(t)=0$.
\end{definition}

\textbf{Derivation.} The projectile lands when its height returns to zero:
\[
y(t)=u\sin\theta\; t-\tfrac{1}{2}gt^{2}=t\left(u\sin\theta-\tfrac{1}{2}gt\right)=0.
\]
The solution $t=0$ is the launch instant. The other solution gives the landing time:
\[
u\sin\theta-\tfrac{1}{2}gT=0 \quad\Longrightarrow\quad T=\frac{2u\sin\theta}{g}.
\]

\begin{propriete}[Time of flight on level ground]
For a projectile launched from level ground with speed $u$ at angle $\theta$,
\[
\boxed{T=\dfrac{2u\sin\theta}{g}}
\]
The derivation above is examinable and should be reproduced rather than quoted.
\end{propriete}

Notice that $T$ depends only on the \emph{vertical} component of the launch velocity: doubling $u\sin\theta$ doubles the time in the air. The horizontal component $u\cos\theta$ plays no role in how long the projectile stays airborne — a direct consequence of the independence of horizontal and vertical motions.

\courssub{Maximum Height and the Symmetry of the Trajectory}

Intuition: the projectile rises until its vertical velocity is exhausted, then falls. At the very top, the motion is instantaneously horizontal, so $v_y=0$ (while $v_x=u\cos\theta$ is unchanged).

\textbf{Derivation.} Setting $v_y=0$:
\[
u\sin\theta-gt=0 \quad\Longrightarrow\quad t_{\text{top}}=\frac{u\sin\theta}{g}=\frac{T}{2}.
\]
Substituting into $y(t)$:
\[
H=y\!\left(\tfrac{T}{2}\right)=u\sin\theta\cdot\frac{u\sin\theta}{g}-\frac{1}{2}g\left(\frac{u\sin\theta}{g}\right)^{2}
=\frac{u^{2}\sin^{2}\theta}{g}-\frac{u^{2}\sin^{2}\theta}{2g}
=\frac{u^{2}\sin^{2}\theta}{2g}.
\]

\begin{propriete}[Maximum height and symmetry]
The \cle{maximum height} above the launch point is
\[
\boxed{H=\dfrac{u^{2}\sin^{2}\theta}{2g}}
\]
reached at time $t=T/2$. The trajectory is \cle{symmetric} about the vertical line through the highest point: the time going up equals the time coming down, and the speed at any height on the way down equals the speed at the same height on the way up.
\end{propriete}

\begin{attention}[$H$ is measured from the launch point]
The formula $H=\dfrac{u^{2}\sin^{2}\theta}{2g}$ gives the height gained \emph{above the launch point}. If the projectile is launched from a height $h$ above the ground (a cliff, a balcony), the greatest height above the \emph{ground} is $h+H$. Mixing these two is a very common error.
\end{attention}

\courssub{Horizontal Range on Level Ground}

The \cle{range} $R$ is the horizontal distance between the launch and landing points. Since $v_x$ is constant, we simply evaluate $x$ at the landing time $T$:
\[
R=x(T)=u\cos\theta\cdot\frac{2u\sin\theta}{g}=\frac{2u^{2}\sin\theta\cos\theta}{g}.
\]
Using the double-angle identity $2\sin\theta\cos\theta=\sin 2\theta$:

\begin{propriete}[Range on level ground]
\[
\boxed{R=\dfrac{u^{2}\sin 2\theta}{g}}
\]
valid \textbf{only} when launch and landing are at the same height.
\end{propriete}

\begin{aretenir}[The three key formulae — level ground, launch from origin]
\[
T=\frac{2u\sin\theta}{g}, \qquad
H=\frac{u^{2}\sin^{2}\theta}{2g}, \qquad
R=\frac{u^{2}\sin 2\theta}{g}.
\]
Conditions of validity: flat horizontal ground, launch and landing at the same level, no air resistance, $g$ constant.
\end{aretenir}

\begin{popfigure}
\begin{center}
\begin{tikzpicture}[scale=0.9]
\draw[->,popmuted] (-0.3,0) -- (10.6,0) node[below left] {$x$};
\draw[->,popmuted] (0,-0.3) -- (0,3.0) node[below left] {$y$};
\draw[very thick,popblue,domain=0:10,samples=60] plot (\x,{0.08*\x*(10-\x)});
\fill[popdark] (0,0) circle (2pt) node[below left] {$O$};
\fill[popdark] (10,0) circle (2pt) node[below right] {landing};
\draw[dashed,poporange] (5,0) -- (5,2.0);
\fill[poporange] (5,2.0) circle (2pt);
\draw[<->,poporange] (5.2,0) -- (5.2,2.0) node[midway,right] {$H$};
\draw[<->,popgreen] (0,-0.55) -- (10,-0.55) node[midway,below] {$R$};
\draw[->,popdark,thick] (0,0) -- (2.2,1.76) node[midway,above left] {$\vec{u}$};
\draw[popdark] (1.1,0) arc (0:38.7:1.1) node[midway,right] {$\theta$};
\draw[popdark] (8.9,0) arc (180:141.3:1.1) node[midway,left] {$\theta$};
\node[poporange,align=center] at (5,2.55) {axis of symmetry, $t=T/2$};
\node[popblue] at (1.9,2.35) {$T=\dfrac{2u\sin\theta}{g}$};
\end{tikzpicture}
\end{center}
\legende{Annotated trajectory on level ground: time of flight $T$, maximum height $H$, range $R$, equal launch and landing angles, and the axis of symmetry through the summit.}
\end{popfigure}

\courssub{Maximum Range and Complementary Angles}

For a fixed launch speed $u$, the range $R=\dfrac{u^{2}}{g}\sin 2\theta$ depends on $\theta$ only through $\sin 2\theta$. Since the greatest value of the sine function is $1$, attained when $2\theta=90^{\circ}$:

\begin{propriete}[Maximum range and complementary angles]
For a given speed $u$ on level ground:
\begin{enumerate}
\item the range is greatest when $\theta=45^{\circ}$, giving $R_{\max}=\dfrac{u^{2}}{g}$;
\item two complementary angles $\theta$ and $90^{\circ}-\theta$ give the \emph{same} range, because $\sin 2(90^{\circ}-\theta)=\sin(180^{\circ}-2\theta)=\sin 2\theta$.
\end{enumerate}
\end{propriete}

So a ball kicked at $30^{\circ}$ and one kicked at $60^{\circ}$ with the same speed land at the same spot — but the $60^{\circ}$ ball stays in the air longer (larger $u\sin\theta$) and goes higher. Goalkeepers and basketball players exploit this trade-off constantly.

\begin{popfigure}
\begin{center}
\begin{tikzpicture}
\begin{axis}[
width=11cm, height=6.5cm,
xlabel={launch angle $\theta$ (degrees)},
ylabel={$R\,g/u^{2}$},
xmin=0, xmax=90, ymin=0, ymax=1.15,
xtick={0,15,30,45,60,75,90},
ytick={0,0.5,1},
grid=both, grid style={popline!40},
axis lines=left,
]
\addplot[very thick,popblue,domain=0:90,samples=80] {sin(2*x)};
\addplot[mark=*,poporange,mark size=2pt] coordinates {(45,1)};
\addplot[dashed,poporange] coordinates {(45,0) (45,1)};
\addplot[mark=*,popgreen,mark size=2pt] coordinates {(30,0.866) (60,0.866)};
\addplot[dashed,popgreen] coordinates {(30,0) (30,0.866) (60,0.866) (60,0)};
\node[poporange] at (axis cs:45,1.08) {$R_{\max}=u^{2}/g$};
\node[popgreen,align=center] at (axis cs:45,0.72) {equal ranges\\ $30^{\circ}$ and $60^{\circ}$};
\end{axis}
\end{tikzpicture}
\end{center}
\legende{Range (in units of $u^{2}/g$) against launch angle. The maximum occurs at $45^{\circ}$; complementary angles give equal ranges.}
\end{popfigure}

\begin{exoresolu}[Level-ground projectile]
A footballer on a flat pitch in Yaounde kicks a ball with speed $20\ \mathrm{m\,s^{-1}}$ at $30^{\circ}$ above the horizontal. Taking $g=10\ \mathrm{m\,s^{-2}}$, find (a) the time of flight, (b) the maximum height, (c) the range, (d) the speed and direction of the ball at impact.
\tcblower
\textbf{Solution.} Here $u=20$, $\sin\theta=\tfrac12$, $\cos\theta=\tfrac{\sqrt3}{2}$, $\sin 2\theta=\sin 60^{\circ}=\tfrac{\sqrt3}{2}$.

(a) Time of flight:
\[
T=\frac{2u\sin\theta}{g}=\frac{2(20)(0.5)}{10}=2\ \mathrm{s}.
\]

(b) Maximum height:
\[
H=\frac{u^{2}\sin^{2}\theta}{2g}=\frac{400\times 0.25}{20}=5\ \mathrm{m}.
\]

(c) Range:
\[
R=\frac{u^{2}\sin 2\theta}{g}=\frac{400\times \frac{\sqrt3}{2}}{10}=20\sqrt3\approx 34.6\ \mathrm{m}.
\]

(d) At $t=T=2$ s:
\[
v_x=20\cos 30^{\circ}=10\sqrt3\approx 17.3\ \mathrm{m\,s^{-1}},\qquad
v_y=20\sin 30^{\circ}-10(2)=10-20=-10\ \mathrm{m\,s^{-1}}.
\]
Speed at impact:
\[
v=\sqrt{v_x^{2}+v_y^{2}}=\sqrt{300+100}=20\ \mathrm{m\,s^{-1}},
\]
directed at $\tan^{-1}\!\left(\dfrac{10}{10\sqrt3}\right)=30^{\circ}$ \emph{below} the horizontal. Note the symmetry: impact speed equals launch speed, impact angle equals launch angle — exactly as the symmetry property predicts.
\end{exoresolu}

\courssub{When the Landing Level Differs from the Launch Level}

The boxed formulae above assume level ground. If the projectile lands higher or lower than it was launched, they fail — and using them is one of the most penalised mistakes at GCE.

\begin{methode}[Projectile landing at a different height]
\begin{enumerate}
\item Write $x(t)=u\cos\theta\,t$ and $y(t)=u\sin\theta\,t-\tfrac12 gt^{2}$ from scratch.
\item Set $y(t)$ equal to the \emph{landing height} (e.g. $y=-h$ for a drop of $h$) and solve the quadratic in $t$; keep the positive root.
\item Compute the horizontal distance as $x=u\cos\theta\,t$ with that value of $t$.
\item Do \textbf{not} use $R=u^{2}\sin 2\theta/g$.
\end{enumerate}
\end{methode}

\textbf{Launch from a height $h$.} Suppose the projectile is launched from a point at height $h$ above flat ground, and we measure $y$ upwards from the launch point. Landing occurs when $y=-h$:
\[
u\sin\theta\,T-\tfrac12 gT^{2}=-h
\quad\Longrightarrow\quad
\tfrac12 gT^{2}-u\sin\theta\,T-h=0,
\]
a quadratic in $T$ whose positive root is
\[
T=\frac{u\sin\theta+\sqrt{u^{2}\sin^{2}\theta+2gh}}{g}.
\]
The horizontal distance travelled is then $X=u\cos\theta\,T$, and the greatest height above the ground is $h+\dfrac{u^{2}\sin^{2}\theta}{2g}$. In an examination, prefer solving the quadratic with the actual numbers rather than quoting this last formula.

\begin{exoresolu}[Stone thrown from a cliff]
A stone is thrown horizontally with speed $15\ \mathrm{m\,s^{-1}}$ from the top of a cliff $45\ \mathrm{m}$ above the sea. Taking $g=10\ \mathrm{m\,s^{-2}}$, find the time of flight, the horizontal distance travelled, and the velocity at impact.
\tcblower
\textbf{Solution.} Launch is horizontal, so $\theta=0$: $u\sin\theta=0$, $u\cos\theta=15$. Take $y$ upwards from the top of the cliff.

Time of flight: $y(t)=-\tfrac12(10)t^{2}=-5t^{2}$. Landing: $-5t^{2}=-45$ gives $t^{2}=9$, so $T=3\ \mathrm{s}$.

Horizontal distance: $X=15\times 3=45\ \mathrm{m}$.

Velocity at impact: $v_x=15\ \mathrm{m\,s^{-1}}$, $v_y=-10(3)=-30\ \mathrm{m\,s^{-1}}$, so
\[
v=\sqrt{15^{2}+30^{2}}=\sqrt{1125}\approx 33.5\ \mathrm{m\,s^{-1}},\qquad
\tan\alpha=\frac{30}{15}=2,
\]
i.e. $\alpha\approx 63.4^{\circ}$ below the horizontal.
\end{exoresolu}

\textbf{Sloping ground (GCE-useful extension).} If the ground is a slope, say the line $y=mx$ (or $y=mx+c$), the landing point is found by substituting the trajectory equation into the equation of the slope and solving for $x$:
\[
x\tan\theta-\frac{gx^{2}}{2u^{2}\cos^{2}\theta}=mx.
\]
The non-zero root gives the landing point directly. This avoids finding the time of flight first and is often the fastest route in examination questions.

\begin{attention}[Do not quote the level-ground formulae blindly]
$R=\dfrac{u^{2}\sin 2\theta}{g}$ and $T=\dfrac{2u\sin\theta}{g}$ are valid \emph{only} when launch and landing are at the same height. On a slope, from a cliff, or onto a raised platform, return to $x(t)$ and $y(t)$ and solve.
\end{attention}

\begin{methode}[Exam technique: derive, then compute]
GCE mark schemes award method marks for the derivation. A safe routine:
\begin{enumerate}
\item State $x=u\cos\theta\,t$, $y=u\sin\theta\,t-\tfrac12 gt^{2}$.
\item For $T$: solve $y=0$ (or the appropriate landing condition).
\item For $H$: set $v_y=0$, find $t$, substitute into $y$.
\item For $R$: substitute $T$ into $x$.
\end{enumerate}
Quoting $R=u^{2}\sin 2\theta/g$ with no justification when the ground is not level scores zero for that part.
\end{methode}

\begin{exercice}[]
A projectile is launched from level ground at $25\ \mathrm{m\,s^{-1}}$. Taking $g=10\ \mathrm{m\,s^{-2}}$, find the two possible launch angles for which the range is $50\ \mathrm{m}$, and compare the two times of flight.
\end{exercice}

\begin{corrige}[Exercise 1]
$R=\dfrac{u^{2}\sin 2\theta}{g}$ gives $50=\dfrac{625\sin 2\theta}{10}$, so $\sin 2\theta=0.8$, hence $2\theta=53.13^{\circ}$ or $126.87^{\circ}$, i.e. $\theta\approx 26.6^{\circ}$ or $63.4^{\circ}$ (complementary, as expected). Times of flight: $T_1=\dfrac{2(25)\sin 26.6^{\circ}}{10}\approx 2.24\ \mathrm{s}$ and $T_2=\dfrac{2(25)\sin 63.4^{\circ}}{10}\approx 4.47\ \mathrm{s}$. Same range, but the steeper launch stays in the air about twice as long.
\end{corrige}

\begin{aretenir}[Summary of key formulae and their conditions]
\begin{center}
\begin{tabularx}{\linewidth}{|l|X|X|}
\hline
\rowcolor{popblueL} Quantity & Formula & Conditions \\
\hline
Time of flight $T$ & $\dfrac{2u\sin\theta}{g}$ & level ground \\
\hline
Maximum height $H$ & $\dfrac{u^{2}\sin^{2}\theta}{2g}$ & above launch point \\
\hline
Range $R$ & $\dfrac{u^{2}\sin 2\theta}{g}$ & level ground only \\
\hline
Maximum range & $\dfrac{u^{2}}{g}$ at $\theta=45^{\circ}$ & level ground \\
\hline
Time to summit & $T/2$ & any landing level (for the summit itself) \\
\hline
\end{tabularx}
\end{center}
When conditions fail, go back to $x(t)$ and $y(t)$: they never fail.
\end{aretenir}

\begin{savaistu}[Why $45^{\circ}$ is not always best]
On level ground $45^{\circ}$ maximises the range, but when launching from a height $h$ the optimal angle is slightly \emph{less} than $45^{\circ}$ — shot-putters, who release the ball about $2\ \mathrm{m}$ above the ground, indeed throw at angles near $40^{\circ}$ to $42^{\circ}$. The extra height gives the lower, faster trajectory more time to travel.
\end{savaistu}

\coursec{Applications and Problem-Solving with Projectiles}

In the previous section we established the three headline results of projectile motion: the time of flight $T = \dfrac{2u\sin\theta}{g}$, the maximum height $H = \dfrac{u^2\sin^2\theta}{2g}$ and the range $R = \dfrac{u^2\sin 2\theta}{g}$. Those formulas answer the question \emph{``what happens when we launch a projectile?''}. In real life and in the GCE examination, the question is usually the reverse: \emph{``how must we launch the projectile so that a given thing happens?''} — so that the ball clears the net, lands in the basket, or reaches a target. This section is devoted to that art: translating a concrete condition into an equation, and solving it.

\courssub{The General Strategy for Projectile Problems}

Almost every projectile problem is solved by the same five-step routine. Master it once and the whole chapter becomes mechanical.

\begin{methode}[The five-step strategy for projectile problems]
\begin{enumerate}
\item \textbf{Draw a diagram.} Sketch the launch point, the landing point or target, and any obstacle. Mark all given distances and heights.
\item \textbf{Choose axes.} Take the origin at the launch point (unless the problem suggests otherwise), the $x$-axis horizontal in the direction of motion and the $y$-axis vertically upwards.
\item \textbf{Resolve the initial velocity.} If the launch speed is $u$ at angle $\theta$ above the horizontal, then $u_x = u\cos\theta$ and $u_y = u\sin\theta$.
\item \textbf{Write the equations of motion.}
\[ x(t) = (u\cos\theta)\,t, \qquad y(t) = (u\sin\theta)\,t - \tfrac{1}{2}gt^2, \]
or, eliminating $t$, the \cle{trajectory equation}
\[ y = x\tan\theta - \frac{g\,x^2}{2u^2\cos^2\theta} = x\tan\theta - \frac{g\,x^2}{2u^2}\left(1+\tan^2\theta\right). \]
\item \textbf{Translate the condition of the problem into an equation.} ``Hits the target $(x_0,y_0)$'' means $y = y_0$ when $x = x_0$; ``clears the wall'' means $y \geq h$ when $x = d$; ``lands on the ground'' means $y = 0$ (or $y = -h$ if launched from a height). Then solve.
\end{enumerate}
\end{methode}

\begin{attention}[Exam tip: keep $g$ symbolic]
Do not substitute $g = 9.81\ \mathrm{m\,s^{-2}}$ (or $10$) until the very last line of your calculation. Rounding $g$ early makes errors accumulate through the working, and examiners award method marks only if your symbolic line is correct. Carry $g$ as a letter, simplify, and only then compute the numerical value.
\end{attention}

\courssub{Hitting a Given Target: the Quadratic in $\tan\theta$}

Suppose a projectile is launched from the origin with speed $u$ and we want it to pass through the point $(x_0, y_0)$. Which angle $\theta$ should we choose? Substituting $(x_0, y_0)$ into the trajectory equation gives
\[ y_0 = x_0\tan\theta - \frac{g\,x_0^{\,2}}{2u^2}\left(1+\tan^2\theta\right). \]

\begin{propriete}[Condition to pass through $(x_0,y_0)$]
Setting $m = \tan\theta$, the condition becomes a \cle{quadratic equation in} $\tan\theta$:
\[ \frac{g\,x_0^{\,2}}{2u^2}\,m^2 - x_0\,m + \left(y_0 + \frac{g\,x_0^{\,2}}{2u^2}\right) = 0. \tag{1} \]
The key algebraic step is the identity $\sec^2\theta = 1 + \tan^2\theta$ applied to $\dfrac{1}{\cos^2\theta}$.
\end{propriete}

Since a quadratic has (in general) two roots, there are usually \textbf{two} launch angles that hit the same target: a \emph{low, flat} trajectory and a \emph{high, lobbed} one. This matches experience: a footballer can reach a teammate either with a driven pass or with a chipped ball.

\begin{attention}[Two values of $\tan\theta$ usually mean two possible angles]
When the quadratic in $\tan\theta$ has two positive roots, both correspond to genuine trajectories through the target — check that both are physically acceptable in the context (a ``low'' solution may be blocked by an obstacle; a negative root would mean firing below the horizontal and is usually rejected). If the discriminant is negative, the target is \emph{out of reach} for the given speed $u$: no angle works. The boundary case (discriminant zero) gives the single angle that just reaches the target; this defines the \cle{envelope of trajectories} (parabola of safety) $y = \dfrac{u^2}{2g} - \dfrac{g x^2}{2u^2}$.
\end{attention}

\begin{exoresolu}[Hitting a target]
A projectile is launched from the origin with speed $u = 20\ \mathrm{m\,s^{-1}}$. Find the angle(s) of projection for it to pass through the point $(30\ \mathrm{m},\ 10\ \mathrm{m})$. Take $g = 10\ \mathrm{m\,s^{-2}}$.
\tcblower
\textbf{Solution.} The trajectory condition at $(x_0,y_0) = (30,10)$ is
\[ 10 = 30\tan\theta - \frac{10 \times 30^2}{2 \times 20^2}\left(1+\tan^2\theta\right). \]
Compute the coefficient: $\dfrac{10 \times 900}{800} = 11.25$. With $m = \tan\theta$:
\[ 10 = 30m - 11.25(1+m^2) \;\Longrightarrow\; 11.25\,m^2 - 30m + 21.25 = 0. \]
Multiply by $4$ to clear decimals: $45m^2 - 120m + 85 = 0$, i.e. $9m^2 - 24m + 17 = 0$. Then
\[ m = \frac{24 \pm \sqrt{576 - 612}}{18} = \frac{24 \pm \sqrt{-36}}{18}. \]
The discriminant is negative: with $u = 20\ \mathrm{m\,s^{-1}}$ the target is \textbf{unreachable}. (Indeed the maximum range on level ground is $u^2/g = 40\ \mathrm{m}$, but the point $(30,10)$ lies above the reachable envelope.) Increasing the speed to $u = 25\ \mathrm{m\,s^{-1}}$ gives coefficient $\dfrac{10 \times 900}{2 \times 625} = 7.2$, so
\[ 7.2m^2 - 30m + 17.2 = 0, \qquad m = \frac{30 \pm \sqrt{900 - 495.36}}{14.4} = \frac{30 \pm 20.12}{14.4}, \]
giving $m \approx 3.48$ ($\theta \approx 74.0°$, high lob) or $m \approx 0.686$ ($\theta \approx 34.5°$, flat throw). Both angles hit the target.
\end{exoresolu}

\courssub{Clearing an Obstacle}

A very common GCE situation: a ball must pass \emph{over} a wall, a net or a fence of height $h$ situated at horizontal distance $d$ from the launch point.

\begin{methode}[``Just clears an obstacle'' — the limiting case]
The projectile clears the obstacle if $y \geq h$ when $x = d$. The \textbf{limiting case} — minimum launch speed, or critical angle — corresponds to the trajectory passing exactly through the top of the obstacle:
\[ h = d\tan\theta - \frac{g\,d^{\,2}}{2u^2}\left(1+\tan^2\theta\right). \]
Impose this equality to find the minimum speed $u_{\min}$ (for a given angle) or the critical angle (for a given speed). Any speed above $u_{\min}$ clears the wall comfortably.
\end{methode}

\begin{popfigure}
\begin{center}
\begin{tikzpicture}[scale=0.9, >=stealth]
\draw[->, thick] (0,0) -- (8.6,0) node[below left] {$x$};
\draw[->, thick] (0,0) -- (0,4.2) node[below left] {$y$};
% wall
\fill[popmuted] (5.4,0) rectangle (5.75,2.2);
\draw[popdark, thick] (5.4,0) rectangle (5.75,2.2);
\node[popdark] at (5.57,-0.35) {$d$};
\node[popdark, right] at (5.75,1.1) {$h$};
% limiting trajectory tangent to top of wall
\draw[very thick, popblue, domain=0:8.2, samples=80] plot (\x, {0.9*\x - 0.0625*\x*\x});
\node[popblue] at (7.6,2.6) {limiting};
% safe trajectory
\draw[very thick, poporange, dashed, domain=0:8.2, samples=80] plot (\x, {1.15*\x - 0.075*\x*\x});
\node[poporange] at (2.6,2.9) {clears easily};
% launch angle
\draw[popgreen, thick] (1.4,0) arc (0:42:1.4);
\node[popgreen] at (1.85,0.42) {$\theta$};
\draw[->, popgreen, thick] (0,0) -- (1.05,0.94) node[midway, below left] {$\vec{u}$};
\fill (0,0) circle (2pt) node[below left] {$O$};
\fill[poppink] (5.575,2.2) circle (2.2pt);
\end{tikzpicture}
\end{center}
\legende{A projectile clearing a wall of height $h$ at distance $d$. The solid curve is the limiting trajectory, passing exactly through the top of the wall: $y(d) = h$.}
\end{popfigure}

\begin{exemplebox}[Ball over a wall]
A footballer kicks a ball from ground level at $\theta = 45°$. A wall of height $h = 2\ \mathrm{m}$ stands $d = 10\ \mathrm{m}$ away. What is the minimum kicking speed for the ball to clear the wall? (Take $g = 10\ \mathrm{m\,s^{-2}}$.)

\emph{Solution.} Limiting case: $y = h$ at $x = d$ with $\tan 45° = 1$:
\[ 2 = 10 - \frac{10 \times 10^2}{2u^2}(1+1) \;\Longrightarrow\; \frac{1000}{u^2} = 8 \;\Longrightarrow\; u^2 = 125, \quad u_{\min} \approx 11.2\ \mathrm{m\,s^{-1}}. \]
\end{exemplebox}

\courssub{Horizontal Projection from a Height}

If a projectile is launched \emph{horizontally} from a height $h$ — a stone rolled off a table, a package dropped from a helicopter moving horizontally, water spurting from a horizontal pipe on a bridge — then $\theta = 0$ and the equations simplify dramatically:
\[ x(t) = u\,t, \qquad y(t) = -\tfrac{1}{2}g\,t^2 \quad (\text{taking } O \text{ at the launch point, } y \text{ upwards}). \]
The time of flight is fixed by the height alone: falling through $h$ gives
\[ h = \tfrac{1}{2}g\,T^2 \;\Longrightarrow\; T = \sqrt{\frac{2h}{g}}, \qquad \text{range from the foot: } R = uT = u\sqrt{\frac{2h}{g}}. \]
Notice that $T$ does \emph{not} depend on $u$: a fast stone and a slow stone dropped from the same cliff hit the sea at the same instant.

\begin{popfigure}
\begin{center}
\begin{tikzpicture}[scale=0.85, >=stealth]
% cliff
\fill[popmuted] (0,0) rectangle (1.6,3.2);
\draw[popdark, thick] (0,0) -- (0,3.2) -- (1.6,3.2) -- (1.6,0);
\node[popdark] at (0.8,1.6) {cliff};
\draw[<->, popdark] (-0.35,0) -- (-0.35,3.2) node[midway, left] {$h$};
% sea
\fill[popblueL] (1.6,0) rectangle (8.6,-0.5);
\draw[popblue, thick] (1.6,0) -- (8.6,0);
\node[popblue] at (7.6,-0.28) {sea};
% launch
\fill (1.6,3.2) circle (2pt);
\draw[->, popgreen, very thick] (1.6,3.2) -- (2.9,3.2) node[midway, above] {$\vec{u}$};
% trajectory
\draw[very thick, poporange, domain=1.6:6.4, samples=60] plot (\x, {3.2 - 0.3125*(\x-1.6)*(\x-1.6)});
\fill[poppink] (6.4,0) circle (2.2pt);
% range
\draw[<->, popdark, thick] (1.6,-0.85) -- (6.4,-0.85) node[midway, below] {$R = u\sqrt{2h/g}$};
\draw[dashed, gray] (1.6,3.2) -- (6.4,3.2);
\end{tikzpicture}
\end{center}
\legende{Horizontal projection from a cliff of height $h$. The trajectory is a half-parabola; the time of flight depends only on $h$.}
\end{popfigure}

\begin{exoresolu}[Stone from a bridge]
A stone is thrown horizontally with speed $u = 12\ \mathrm{m\,s^{-1}}$ from a bridge $20\ \mathrm{m}$ above a river on the Douala--Yaoundé highway. Take $g = 10\ \mathrm{m\,s^{-2}}$. Find (a) the time of flight, (b) the horizontal distance travelled, (c) the speed with which the stone hits the water.
\tcblower
\textbf{Solution.} (a) $T = \sqrt{\dfrac{2h}{g}} = \sqrt{\dfrac{2 \times 20}{10}} = 2\ \mathrm{s}$.

(b) $R = uT = 12 \times 2 = 24\ \mathrm{m}$.

(c) At impact: $v_x = u = 12\ \mathrm{m\,s^{-1}}$ and $v_y = -gT = -20\ \mathrm{m\,s^{-1}}$. Hence
\[ v = \sqrt{v_x^2 + v_y^2} = \sqrt{144 + 400} = \sqrt{544} \approx 23.3\ \mathrm{m\,s^{-1}}. \]
Note that the impact speed is \emph{not} $u + gT$: the horizontal and vertical components must be combined as vectors.
\end{exoresolu}

\courssub{Two Projectiles and Collision Problems}

A useful GCE extension: two particles $A$ and $B$ are in projectile motion simultaneously. They \cle{collide} if they occupy the \textbf{same position at the same time}: one writes $x_A(t) = x_B(t)$ and $y_A(t) = y_B(t)$ and solves the system for $t$ and the unknown launch parameter. A classic result (the ``monkey and hunter''): if particle $A$ is fired directly at particle $B$ which is released from rest at the same instant, then $A$ always hits $B$ (provided the ground does not intervene), because both acquire the same downward displacement $\tfrac12 g t^2$ relative to the straight line of sight.

\begin{exemplebox}[Collision in mid-air — monkey and hunter]
From the origin $O$, particle $A$ is projected with $u = 20\ \mathrm{m\,s^{-1}}$ at $\theta = 60°$. At the same instant, particle $B$ is released from rest at the point $(10\ \mathrm{m},\ 10\sqrt{3}\ \mathrm{m})$ — i.e. on the initial line of sight of $A$. Show that they collide and find the time and position of collision. (Take $g = 10\ \mathrm{m\,s^{-2}}$.)

\emph{Solution.} For $A$: $u_x = 20\cos60° = 10$, $u_y = 20\sin60° = 10\sqrt{3}$.
\[ x_A = 10t, \quad y_A = 10\sqrt{3}\,t - 5t^2. \]
For $B$ (dropped from rest at $(10, 10\sqrt{3})$):
\[ x_B = 10, \quad y_B = 10\sqrt{3} - 5t^2. \]
Collision requires $x_A = x_B$ and $y_A = y_B$ simultaneously.
$x_A = x_B \;\Rightarrow\; 10t = 10 \;\Rightarrow\; t = 1\ \mathrm{s}$.
At $t=1$: $y_A = 10\sqrt{3} - 5$, $y_B = 10\sqrt{3} - 5$. The $y$-coordinates match automatically.
Thus they collide at $t = 1\ \mathrm{s}$ at the point $(10\ \mathrm{m},\ 10\sqrt{3}-5 \approx 12.3\ \mathrm{m})$.
\end{exemplebox}

\courssub{Real-Life Modelling and the Limits of the Model}

\begin{experience}[Modelling task: the free kick]
A footballer takes a free kick $25\ \mathrm{m}$ from goal. The crossbar is $2.44\ \mathrm{m}$ high and the defensive wall, $9\ \mathrm{m}$ away, jumps to $2\ \mathrm{m}$. Model the ball as a particle launched from ground level at $u = 22\ \mathrm{m\,s^{-1}}$, $\theta = 25°$, with $g = 9.81\ \mathrm{m\,s^{-2}}$.
\begin{itemize}
\item At the wall ($x = 9$): $y = 9\tan 25° - \dfrac{9.81 \times 81}{2 \times 484 \cos^2 25°} \approx 4.20 - 1.00 = 3.20\ \mathrm{m} > 2\ \mathrm{m}$: the ball clears the wall.
\item At the goal line ($x = 25$): $y = 25\tan 25° - \dfrac{9.81 \times 625}{2 \times 484 \cos^2 25°} \approx 11.66 - 7.71 = 3.95\ \mathrm{m}$: too high — it sails over the bar. The player must reduce the speed or the angle. Trying $\theta = 20°$: at the wall $y \approx 2.35\ \mathrm{m}$ (still clears) and at the goal $y \approx 1.93\ \mathrm{m} < 2.44\ \mathrm{m}$: \textbf{goal!}
\end{itemize}
\end{experience}

\begin{savaistu}[Modelling box: what we have ignored]
Our model treats the projectile as a \cle{particle} moving under gravity alone, in a vacuum, over a flat non-rotating Earth. Real projectiles experience \textbf{air resistance} (which shortens the range and makes the descent steeper than the ascent) and \textbf{spin} (the Magnus effect bends footballs and golf balls sideways). A struck football may lose $10$--$20\%$ of its vacuum range. In the examination, when a question says \emph{``state one modelling assumption''} or \emph{``comment on the model''}, expected answers include: air resistance is neglected; the projectile is modelled as a particle (its size is ignored); $g$ is taken as constant; wind is ignored; the ground is horizontal. Stating these clearly earns easy marks.
\end{savaistu}

\begin{attention}[Common mistakes in applied projectile problems]
\begin{itemize}
\item Using the range formula $R = u^2\sin 2\theta / g$ when the landing level differs from the launch level (cliff, raised target). That formula is valid \emph{only} on horizontal ground; otherwise solve $y(t) = -h$.
\item Forgetting that ``just clears'' means \emph{equality} $y(d) = h$, not inequality, when finding a minimum speed.
\item Mixing degrees and radians on the calculator when evaluating $\tan\theta$, $\sin\theta$.
\item Adding velocity components as scalars instead of using $v = \sqrt{v_x^2 + v_y^2}$.
\item Substituting $g = 9.81$ too early and losing accuracy (and method marks).
\end{itemize}
\end{attention}

\begin{exercice}[Practice]
\begin{enumerate}
\item A ball is thrown from ground level at $u = 14\ \mathrm{m\,s^{-1}}$, $\theta = 60°$. Does it clear a fence $3\ \mathrm{m}$ high standing $8\ \mathrm{m}$ away? (Take $g = 9.81\ \mathrm{m\,s^{-2}}$.)
\item A hose held horizontally $1.5\ \mathrm{m}$ above the ground sends water at $8\ \mathrm{m\,s^{-1}}$. How far from the nozzle does the jet strike the ground?
\item A projectile launched from $O$ with speed $28\ \mathrm{m\,s^{-1}}$ must pass through $(40\ \mathrm{m}, 15\ \mathrm{m})$. Find the two possible angles of projection. (Take $g = 9.8\ \mathrm{m\,s^{-2}}$.)
\end{enumerate}
\end{exercice}

\begin{corrige}[Exercise 1]
\begin{enumerate}
\item $y(8) = 8\tan 60° - \dfrac{9.81 \times 64}{2 \times 196 \cos^2 60°} = 13.86 - 6.41 = 7.45\ \mathrm{m} > 3\ \mathrm{m}$: yes, comfortably.
\item $T = \sqrt{2(1.5)/9.81} \approx 0.553\ \mathrm{s}$; $R = 8 \times 0.553 \approx 4.42\ \mathrm{m}$.
\item $15 = 40m - \dfrac{9.8 \times 1600}{2 \times 784}(1+m^2) = 40m - 10(1+m^2)$, so $10m^2 - 40m + 25 = 0$, i.e. $2m^2 - 8m + 5 = 0$: $m = \dfrac{8 \pm \sqrt{24}}{4} = 2 \pm \dfrac{\sqrt6}{2}$, giving $\theta \approx 72.8°$ or $\theta \approx 37.8°$.
\end{enumerate}
\end{corrige}

\begin{aretenir}[Key points of the section]
\begin{itemize}
\item Five steps: diagram, axes, resolve $\vec{u}$, write $x(t)$ and $y(t)$, translate the condition into an equation.
\item Passing through $(x_0,y_0)$ gives a quadratic in $\tan\theta$ via $\sec^2\theta = 1+\tan^2\theta$; two roots usually mean a high and a low trajectory.
\item ``Just clears an obstacle'' $\Rightarrow$ equality $y(d) = h$ (limiting case).
\item Horizontal launch from height $h$: $T = \sqrt{2h/g}$, $R = u\sqrt{2h/g}$; $T$ is independent of $u$.
\item Collisions: same position at the same time — solve $x_A = x_B$, $y_A = y_B$.
\item Keep $g$ symbolic until the last line, and always state modelling assumptions (no air resistance, particle model, constant $g$).
\end{itemize}
\end{aretenir}

\coursec{Free Fall, Vertical Projection and Motion Graphs}

In the previous section we treated projectiles as two independent motions: uniform motion horizontally and uniformly accelerated motion vertically. In this final section we strip away the horizontal part completely and concentrate on the purely vertical motion of a body moving under gravity alone: a stone dropped from a bridge over the Wouri, a mango falling from a tree, a ball thrown straight up on the playground. These \emph{one-dimensional} motions are the simplest of all, yet they contain every idea of the chapter, and they lend themselves beautifully to graphical analysis — a favourite GCE examination theme.

Throughout this section we take $g = 9.81\ \mathrm{m\,s^{-2}}$, often approximated as $g = 10\ \mathrm{m\,s^{-2}}$ when the question allows, and we neglect air resistance.

\courssub{Free Fall: Dropped from Rest}

\begin{experience}[Galileo's idea]
Drop a small stone and a heavy hammer from the same height at the same instant. Contrary to everyday intuition, they hit the ground essentially together. Air resistance, not weight, is what makes a feather fall slowly; in a vacuum a feather and a hammer fall side by side. The conclusion is profound: \emph{all} bodies in free fall have the same acceleration, whatever their mass.
\end{experience}

\begin{definition}[Free fall]
A body is in \cle{free fall} when the \emph{only} force acting on it is its \cle{weight}. By Newton's second law its acceleration is then $g$ vertically downwards, \emph{whatever its velocity} — whether the body is moving down, is momentarily at rest, or is even moving upwards.
\end{definition}

\begin{attention}[Free fall does not mean ``falling down'']
A ball thrown \emph{upwards} is in free fall from the instant it leaves the hand: the only force on it is its weight, so its acceleration is $g$ downwards throughout. ``Free fall'' describes the \emph{force situation}, not the direction of motion.
\end{attention}

Consider now a body \emph{dropped} (released from rest, $u = 0$) from a height. Take the \cle{downwards} direction as positive. Then $a = g$, $u = 0$, and the constant-acceleration formulae $v = u + at$, $s = ut + \tfrac12 at^2$, $v^2 = u^2 + 2as$ give immediately:

\begin{propriete}[Free fall from rest, downwards positive]
For a body dropped from rest and falling through a height $h$:
\[ v = gt, \qquad h = \tfrac{1}{2}gt^{2}, \qquad v^{2} = 2gh. \]
The speed reached after falling a height $h$ is $v = \sqrt{2gh}$, independent of the mass of the body. (Proof: direct substitution into the constant-acceleration formulae; examinable as a one-line derivation.)
\end{propriete}

\begin{exemplebox}[Stone dropped from a bridge]
A stone is dropped from a bridge $45\ \mathrm{m}$ above a river. Taking $g = 10\ \mathrm{m\,s^{-2}}$, the time to reach the water satisfies $45 = \tfrac12(10)t^2$, so $t^2 = 9$ and $t = 3\ \mathrm{s}$. Its speed on entering the water is $v = gt = 10 \times 3 = 30\ \mathrm{m\,s^{-1}}$ (check: $v = \sqrt{2gh} = \sqrt{2 \times 10 \times 45} = \sqrt{900} = 30\ \mathrm{m\,s^{-1}}$).
\end{exemplebox}

\courssub{Object Thrown Vertically Upwards}

Now throw a ball straight up with initial speed $u$. Choose \cle{upwards as positive}. The acceleration is $a = -g$ (gravity points down), and the constant-acceleration formulae become:

\begin{aretenir}[Upward projection, upwards positive]
\[ v = u - gt, \qquad y = ut - \tfrac{1}{2}gt^{2}, \qquad v^{2} = u^{2} - 2gy. \]
At the greatest height $v = 0$, giving:
\[ \text{time to greatest height} = \frac{u}{g}, \qquad \text{greatest height} = H = \frac{u^{2}}{2g}. \]
The total time of flight back to the launch point is $T = \dfrac{2u}{g}$.
\end{aretenir}

\textbf{Derivation of the greatest height.}\par
At the top of the flight the body is momentarily at rest: $v = 0$. From $v = u - gt$, the time to reach the top is $t = u/g$. Substituting into $y = ut - \tfrac12 gt^2$:
\[ H = u\left(\frac{u}{g}\right) - \frac{1}{2}g\left(\frac{u}{g}\right)^{2} = \frac{u^{2}}{g} - \frac{u^{2}}{2g} = \frac{u^{2}}{2g}. \]
(The same result comes instantly from $v^2 = u^2 - 2gy$ with $v = 0$.)

\begin{propriete}[The highest point]
At the greatest height the \cle{velocity} is instantaneously \emph{zero}, but the \cle{acceleration} is still $g$ directed downwards. The body does not ``stop accelerating'' at the top — it is in free fall at every instant of the flight.
\end{propriete}

\begin{attention}[Zero velocity does not mean zero acceleration]
A very common GCE error is to write $a = 0$ at the highest point ``because the ball stops''. The ball stops only \emph{for an instant}; gravity never switches off. If $a$ were zero at the top, the ball would stay there forever!
\end{attention}

\textbf{Symmetry of vertical motion.}\par
Set $y = 0$ in $y = ut - \tfrac12 gt^2$: $t(u - \tfrac12 gt) = 0$, so the body returns to the launch level at $t = \dfrac{2u}{g}$ — exactly \emph{twice} the time to reach the top. Hence \cle{time up = time down}. Moreover, at $t = 2u/g$:
\[ v = u - g\cdot\frac{2u}{g} = -u, \]
so the body returns with the \emph{same speed} $u$, directed downwards. More generally, at any given height on the way up and on the way down, the speeds are equal (from $v^2 = u^2 - 2gy$, which depends only on $y$). This perfect symmetry is a powerful checking tool in examinations.

\begin{exoresolu}[Ball thrown vertically upwards]
A ball is thrown vertically upwards with speed $24.5\ \mathrm{m\,s^{-1}}$ from ground level. Taking $g = 9.81\ \mathrm{m\,s^{-2}}$, find (a) the greatest height reached, (b) the total time of flight, (c) the speed with which it returns to the thrower's hand.
\tcblower
Take upwards as positive: $u = 24.5$, $a = -9.81$.

\textbf{(a)} At the top, $v = 0$. Using $v^2 = u^2 - 2gy$:
\[ 0 = 24.5^{2} - 2(9.81)H \;\Longrightarrow\; H = \frac{24.5^{2}}{2 \times 9.81} = \frac{600.25}{19.62} \approx 30.6\ \mathrm{m}. \]

\textbf{(b)} Time to the top: $t = \dfrac{u}{g} = \dfrac{24.5}{9.81} \approx 2.50\ \mathrm{s}$. By symmetry the total flight time is
\[ T = \frac{2u}{g} \approx 5.0\ \mathrm{s}. \]

\textbf{(c)} By symmetry the ball returns with speed $u = 24.5\ \mathrm{m\,s^{-1}}$ downwards, i.e. velocity $-24.5\ \mathrm{m\,s^{-1}}$. Check: $v = 24.5 - 9.81(5.0) \approx -24.5\ \mathrm{m\,s^{-1}}$. \checkmark
\end{exoresolu}

\courssub{Object Thrown Vertically Downwards; Release from a Moving Body}

If a body is \emph{thrown downwards} with speed $u$ from a height, take \cle{downwards as positive}; then $a = +g$ and:

\begin{propriete}[Downward projection, downwards positive]
\[ v = u + gt, \qquad y = ut + \tfrac{1}{2}gt^{2}, \qquad v^{2} = u^{2} + 2gy. \]
\end{propriete}

\begin{methode}[One sign convention per question]
\begin{enumerate}
\item Choose \emph{one} positive direction (usually upwards) and write it down.
\item Give every quantity its sign accordingly: $u$, $v$, $y$ positive if directed/measured upwards, negative if downwards; $a = -g$ always (upwards positive).
\item Substitute into $v = u + at$, $y = ut + \tfrac12 at^2$, $v^2 = u^2 + 2ay$ \emph{with the signs included}.
\item Interpret the sign of the answer: negative velocity means ``moving downwards'', negative displacement means ``below the launch point''.
\end{enumerate}
Never change convention in the middle of a question — this is the source of most sign errors.
\end{methode}

A subtle and frequently examined situation is the release of an object from a \emph{moving} body — a package dropped from a rising balloon, a stone let go from a descending lift.

\begin{attention}[The balloon trap]
``\emph{Dropped}'' means $u = 0$. But ``\emph{released} from a moving balloon'' means the object \emph{keeps the balloon's velocity at the instant of release}. A package released from a balloon rising at $5\ \mathrm{m\,s^{-1}}$ has initial velocity $u = +5\ \mathrm{m\,s^{-1}}$ (upwards): it first continues to \emph{rise}, slows, stops, and only then falls. Writing $u = 0$ here is a classic GCE trap.
\end{attention}

\begin{exoresolu}[Package released from a rising balloon]
A balloon is rising vertically at a steady $12\ \mathrm{m\,s^{-1}}$. When it is $80\ \mathrm{m}$ above the ground, a package is released. Taking $g = 10\ \mathrm{m\,s^{-2}}$, find the time taken for the package to reach the ground and its speed on impact.
\tcblower
Take upwards as positive, origin at the release point. At release the package has the balloon's velocity: $u = +12\ \mathrm{m\,s^{-1}}$, and $a = -10\ \mathrm{m\,s^{-2}}$. The ground is at $y = -80\ \mathrm{m}$.

Using $y = ut - \tfrac12 gt^2$:
\[ -80 = 12t - 5t^{2} \;\Longrightarrow\; 5t^{2} - 12t - 80 = 0. \]
\[ t = \frac{12 \pm \sqrt{144 + 1600}}{10} = \frac{12 \pm \sqrt{1744}}{10} \approx \frac{12 \pm 41.76}{10}. \]
Taking the positive root: $t \approx 5.38\ \mathrm{s}$.

Impact velocity: $v = u - gt = 12 - 10(5.38) \approx -41.8\ \mathrm{m\,s^{-1}}$, i.e. a speed of about $41.8\ \mathrm{m\,s^{-1}}$ downwards. Note that the package first rose for $u/g = 1.2\ \mathrm{s}$ before falling — exactly the behaviour the sign convention predicts automatically.
\end{exoresolu}

\courssub{Graphical Representation of Vertical Motion}

Graphs turn the formulae above into pictures, and the GCE regularly asks candidates to sketch or interpret them. Recall the three golden links from earlier work:

\begin{aretenir}[Reading motion graphs]
\begin{itemize}
\item \cle{Gradient} of the displacement--time ($s$--$t$) graph $=$ \cle{velocity}.
\item \cle{Gradient} of the velocity--time ($v$--$t$) graph $=$ \cle{acceleration}.
\item \cle{Area} under the $v$--$t$ graph $=$ \cle{displacement} (area below the time axis counts \emph{negative}).
\end{itemize}
For upward projection (upwards positive): the $v$--$t$ graph is a straight line of gradient $-g$ \emph{always}; it crosses the time axis exactly at the highest point. The $a$--$t$ graph is the horizontal line $a = -g$ for the whole flight. The $s$--$t$ graph is a parabola whose maximum occurs at the same instant.
\end{aretenir}

\begin{popfigure}
\begin{center}
\begin{tikzpicture}
\begin{axis}[
  width=0.9\linewidth, height=4.6cm,
  axis lines=middle, xlabel={$t$ (s)}, ylabel={$y$ (m)},
  xlabel style={at={(ticklabel* cs:1)}, anchor=west},
  ylabel style={at={(ticklabel* cs:1)}, anchor=south},
  xmin=0, xmax=5.4, ymin=0, ymax=35,
  xtick={2.5,5}, xticklabels={$u/g$, $2u/g$}, ytick=\empty,
  domain=0:5, samples=60, clip=false]
\addplot[popblue, very thick] {24.5*x - 4.905*x^2};
\addplot[popmuted, dashed] coordinates {(2.5,0) (2.5,30.6)};
\node[popblue, font=\small] at (axis cs:3.4,26) {parabola};
\node[popdark, font=\small, align=center] at (axis cs:2.5,33.5) {max height};
\end{axis}
\end{tikzpicture}

\vspace{0.4cm}
\begin{tikzpicture}
\begin{axis}[
  width=0.9\linewidth, height=4.6cm,
  axis lines=middle, xlabel={$t$ (s)}, ylabel={$v$ (m\,s$^{-1}$)},
  xlabel style={at={(ticklabel* cs:1)}, anchor=west},
  ylabel style={at={(ticklabel* cs:1)}, anchor=south},
  xmin=0, xmax=5.4, ymin=-30, ymax=30,
  xtick={2.5,5}, xticklabels={$u/g$, $2u/g$}, ytick=\empty,
  domain=0:5, samples=10, clip=false]
\addplot[poporange, very thick] {24.5 - 9.81*x};
\addplot[popmuted, dashed] coordinates {(2.5,-28) (2.5,28)};
\node[popdark, font=\small, align=center] at (axis cs:1.1,14) {gradient $=-g$};
\node[popdark, font=\small, align=center] at (axis cs:3.6,8) {$v=0$ at top};
\end{axis}
\end{tikzpicture}

\vspace{0.4cm}
\begin{tikzpicture}
\begin{axis}[
  width=0.9\linewidth, height=3.6cm,
  axis lines=middle, xlabel={$t$ (s)}, ylabel={$a$ (m\,s$^{-2}$)},
  xlabel style={at={(ticklabel* cs:1)}, anchor=west},
  ylabel style={at={(ticklabel* cs:1)}, anchor=south},
  xmin=0, xmax=5.4, ymin=-15, ymax=8,
  xtick={2.5,5}, xticklabels={$u/g$, $2u/g$}, ytick=\empty,
  domain=0:5, samples=2, clip=false]
\addplot[popgreen, very thick] {-9.81};
\addplot[popmuted, dashed] coordinates {(2.5,-14) (2.5,6)};
\node[popdark, font=\small] at (axis cs:3.9,-7) {$a=-g$ (constant)};
\end{axis}
\end{tikzpicture}
\end{center}
\legende{Displacement--time, velocity--time and acceleration--time graphs for a ball thrown vertically upwards ($u = 24.5\ \mathrm{m\,s^{-1}}$, $g = 9.81\ \mathrm{m\,s^{-2}}$). The instant of maximum height $t = u/g$ is marked on all three graphs: the parabola peaks, the velocity line crosses the axis, and the acceleration stays fixed at $-g$.}
\end{popfigure}

\textbf{Interpreting the sign changes.}\par
Before the highest point, $v > 0$ (moving up) and the $s$--$t$ curve rises. At the top, $v = 0$: the $v$--$t$ line crosses the axis and the $s$--$t$ parabola has a horizontal tangent. After the top, $v < 0$ (moving down) and the parabola descends. The acceleration graph \emph{never changes}: the same line $a = -g$ governs the ascent, the instant at the top, and the descent. The area under the $v$--$t$ line from $0$ to $2u/g$ is a triangle above the axis minus a congruent triangle below it — total area zero, matching the zero total displacement when the ball returns to the hand.

\textbf{Comparing a thrown-up ball and a dropped ball.}\par
For a dropped ball (downwards positive), the $v$--$t$ graph is a straight line through the \emph{origin} with gradient $+g$, and the $s$--$t$ graph is the rising half-parabola $h = \tfrac12 gt^2$. For the thrown-up ball (upwards positive) the $v$--$t$ line starts at $+u$ and slopes down with gradient $-g$, crossing the axis. Same physics, different initial conditions — the gradient is $\pm g$ in every case.

\textbf{Extension: a bouncing ball.}\par
If a ball bounces on the ground, the $v$--$t$ graph consists of straight segments of gradient $-g$ (upwards positive), with a sudden jump at each bounce (the velocity reverses rapidly from negative to a smaller positive value, since each bounce loses energy). The $s$--$t$ graph is a chain of parabola arcs of decreasing height. Examiners occasionally ask for such a sketch: the key features are the \emph{constant gradient} $-g$ of every segment and the \emph{decreasing} peaks.

\begin{popfigure}
\begin{center}
\begin{tikzpicture}[scale=0.7]
% Vertical axis for height
\draw[->, popdark, thick] (0,0) -- (0,34) node[above, font=\small] {height (m)};
% Time labels on the right
\foreach \n/\t/\y in {
  0/0/0.0,
  1/0.5/11.0,
  2/1.0/19.6,
  3/1.5/25.7,
  4/2.0/29.4,
  5/2.5/30.6,
  6/3.0/29.4,
  7/3.5/25.7,
  8/4.0/19.6,
  9/4.5/11.0,
  10/5.0/0.0
}{
  \fill[popblue] (0.5,\y) circle (3pt);
  \node[font=\tiny, popdark, anchor=west] at (0.8,\y) {$t=\t\,$s};
}
% Mark the top
\fill[poporange] (0.5,30.6) circle (4pt);
\node[font=\small, poporange, anchor=west] at (0.8,30.6) {top ($t=2.5$ s)};
% Annotations
\node[font=\small, popdark, align=left] at (3.5,28) {gaps shrink\\going up};
\node[font=\small, popdark, align=left] at (3.5,6) {gaps grow\\coming down};
% Ground line
\draw[popdark, thick] (-0.5,0) -- (1.0,0);
\node[font=\small, popdark] at (-0.5,-1.5) {ground};
\end{tikzpicture}
\end{center}
\legende{Strobe diagram: positions of a stone thrown vertically upwards ($u=24.5\ \mathrm{m\,s^{-1}}$, $g=9.81\ \mathrm{m\,s^{-2}}$) at equal time intervals $\Delta t = 0.5\ \mathrm{s}$. Going up (blue) the spacing decreases as gravity slows the stone; at the top (orange) it is instantaneously at rest; coming down (blue) the spacing increases symmetrically. The motion is perfectly symmetric about the highest point.}
\end{popfigure}

\begin{exercice}[Vertical motion and graphs]
A stone is thrown vertically upwards with speed $19.62\ \mathrm{m\,s^{-1}}$ from the top of a cliff $40\ \mathrm{m}$ high. Take $g = 9.81\ \mathrm{m\,s^{-2}}$.
\begin{enumerate}
\item Find the greatest height above the cliff top and the time taken to reach it.
\item Find the total time until the stone hits the sea at the foot of the cliff.
\item Find the velocity of the stone on impact.
\item Sketch the $v$--$t$ graph for the whole motion (upwards positive), marking the time at the greatest height and the impact time.
\end{enumerate}
\end{exercice}

\begin{corrige}[Exercise 1]
Take upwards positive, origin at the cliff top: $u = 19.62$, $a = -9.81$.

\textbf{1.} At the top $v = 0$: $t = u/g = 19.62/9.81 = 2\ \mathrm{s}$; $H = u^2/(2g) = 19.62^2/19.62 = 19.62\ \mathrm{m}$ above the cliff top.

\textbf{2.} The sea is at $y = -40$: $-40 = 19.62t - 4.905t^2$, i.e. $4.905t^2 - 19.62t - 40 = 0$. Then $t = \dfrac{19.62 + \sqrt{19.62^2 + 4 \times 4.905 \times 40}}{9.81} = \dfrac{19.62 + \sqrt{384.9 + 784.8}}{9.81} = \dfrac{19.62 + 34.21}{9.81} \approx 5.49\ \mathrm{s}$ (positive root).

\textbf{3.} $v = 19.62 - 9.81(5.49) \approx -34.2\ \mathrm{m\,s^{-1}}$, i.e. $34.2\ \mathrm{m\,s^{-1}}$ downwards.

\textbf{4.} Straight line from $(0, 19.62)$ with gradient $-9.81$, crossing the time axis at $t = 2\ \mathrm{s}$ (greatest height) and ending at $(5.49, -34.2)$. The net signed area under the line equals $-40\ \mathrm{m}$, the total displacement.
\end{corrige}

\begin{aretenir}[Chapter summary: vertical motion]
\begin{itemize}
\item Free fall: only weight acts; $a = g$ downwards at all times, even at the top of the flight.
\item Dropped from rest (down positive): $v = gt$, $h = \tfrac12 gt^2$, $v^2 = 2gh$.
\item Thrown up (up positive): $v = u - gt$, $y = ut - \tfrac12 gt^2$; top at $t = u/g$, height $u^2/(2g)$; total flight time $2u/g$; returns with speed $u$.
\item Released from a moving body: $u =$ velocity of the body at release.
\item Graphs: $s$--$t$ parabola, $v$--$t$ line of gradient $-g$ crossing the axis at the top, $a$--$t$ horizontal at $-g$; gradient of $s$--$t$ is $v$, gradient of $v$--$t$ is $a$, area under $v$--$t$ is displacement.
\end{itemize}
\end{aretenir}

\newpage

\coursec{Integration activity}

\courssub{Aims of the activity}

This integrative activity closes the chapter by mobilising, in one single realistic problem, every competence developed in Lessons 114 to 122. Working in groups of three or four, you will:

\begin{itemize}
\item resolve an initial velocity vector into horizontal and vertical components and write the parametric equations $x(t)$ and $y(t)$ of a projectile;
\item derive and use the \cle{trajectory equation} $y = x\tan\theta - \dfrac{gx^{2}}{2u^{2}}(1+\tan^{2}\theta)$;
\item translate physical constraints (clearing a wall, staying under a crossbar) into inequalities, and solve the resulting \cle{quadratic inequalities in $\tan\theta$};
\item compute a time of flight and a velocity vector at a given instant, then its magnitude;
\item sketch and interpret a \cle{velocity--time graph} for the vertical motion;
\item discuss critically the assumptions of the projectile model, as expected of a GCE Advanced Level candidate.
\end{itemize}

\begin{aretenir}[Data you may use]
Take $g = 10\ \mathrm{m\,s^{-2}}$, neglect air resistance, and treat the ball as a particle. Give angles to the nearest $0.1^{\circ}$ and lengths to $2$ decimal places.
\end{aretenir}

\courssub{Situation}

It is the last minute of the inter-school football final at the sports field of Lyc\'ee de Nkol-Eton, Yaound\'e. Your school has just been awarded a direct free kick. The ball is placed on flat ground, exactly $22\ \mathrm{m}$ from the goal line, directly in front of the centre of the goal. The defending team forms a ``wall'' of players standing $9.1\ \mathrm{m}$ from the ball (the regulation distance is about $9.15\ \mathrm{m}$). The tallest defender jumps and reaches an effective height of $2.0\ \mathrm{m}$, so to pass the wall the ball must be \emph{strictly above} $2.0\ \mathrm{m}$ when it crosses the plane of the wall. The crossbar of the goal is $2.44\ \mathrm{m}$ above the ground, so to score, the ball must cross the goal line \emph{below} $2.44\ \mathrm{m}$ (and, of course, still above the ground).

The free-kick specialist strikes the ball so that it leaves the ground with speed $u = 20\ \mathrm{m\,s^{-1}}$ at an angle $\theta$ above the horizontal. The whole stadium is silent. Your group's job is to advise the kicker: \textbf{for which angles $\theta$ does the ball clear the wall AND enter the goal below the crossbar? What is the ``optimal'' kick?}

\begin{popfigure}
\begin{center}
\begin{tikzpicture}[scale=0.42,>=stealth]
\draw[popgreenL,fill=popgreenL] (-0.5,-0.15) rectangle (24,0);
\draw[popink,thick] (0,0) -- (24,0);
\fill[popink] (0,0) circle (0.18);
\node[below] at (0,-0.2) {\small ball};
\draw[popdark,very thick] (9.1,0) -- (9.1,2.0);
\node[above,align=center,text width=2.2cm] at (9.1,2.0) {\small wall $2.0$ m};
\draw[popdark,very thick] (22,0) -- (22,2.44);
\draw[popdark,very thick] (22,2.44) -- (23.2,2.44);
\node[above,align=center,text width=2.6cm] at (22.6,2.6) {\small crossbar $2.44$ m};
\draw[<->,popmuted] (0,-0.9) -- (9.1,-0.9) node[midway,below]{\small $9.1$ m};
\draw[<->,popmuted] (0,-1.8) -- (22,-1.8) node[midway,below]{\small $22$ m};
\draw[->,poporange,very thick] (0,0) -- (3.2,1.9) node[right]{\small $\vec{u}$};
\draw[poporange] (1.6,0) arc (0:31:1.6) node[midway,right]{\small $\theta$};
% Trajectory for optimal angle theta = 21 deg: tan=0.3839, 1+tan^2=1.1474, coeff = 1.1474/80 = 0.01434
\draw[popblue,dashed,domain=0:22,samples=80] plot (\x, {0.3839*\x - 0.01434*\x*\x});
\end{tikzpicture}
\end{center}
\legende{The free kick: the ball must pass above the wall at $9.1$ m and under the crossbar at $22$ m. The dashed curve shows the trajectory for the optimal angle $\theta\approx21^{\circ}$.}
\end{popfigure}

\courssub{Tasks}

\begin{enumerate}
\item \textbf{Setting up the model.} Choose axes with origin at the ball's initial position, $Ox$ horizontal towards the goal and $Oy$ vertically upwards. Resolve $\vec{u}$ into components and, starting from $\vec{a} = -g\,\vec{j}$, show by integration that
\[ x(t) = 20t\cos\theta, \qquad y(t) = 20t\sin\theta - 5t^{2}. \]
\item \textbf{Trajectory equation.} Eliminate $t$ to obtain
\[ y = x\tan\theta - \frac{x^{2}}{80}\left(1+\tan^{2}\theta\right). \]
\item \textbf{Clearing the wall.} Write down the condition $y(9.1) > 2.0$ and show that it leads to a quadratic inequality in $T = \tan\theta$. Solve it and give the corresponding interval of angles $\theta$.
\item \textbf{Staying under the crossbar.} Write down the condition $0 < y(22) < 2.44$ and show that it also leads to quadratic inequalities in $T$. Solve them.
\item \textbf{The scoring window.} Combine the conditions of Tasks 3 and 4 to find all values of $\theta$ for which the kick scores. Which branch of solutions is physically sensible for a direct free kick? Propose an ``optimal'' angle and justify your choice (margin of safety, time available to the goalkeeper, \dots).
\item \textbf{Flight to the goal line.} For your chosen optimal angle, compute the time at which the ball crosses the goal line and the velocity vector $\vec{v}$ at that instant; deduce the speed of the ball.
\item \textbf{Graphical representation.} Sketch the graph of $v_{y}(t)$ against $t$ for the whole flight, from $t = 0$ until the ball crosses the goal line. Indicate the slope, the intercept, and the instant when $v_{y} = 0$; state what the area under the graph represents.
\item \textbf{Critical discussion.} Identify \emph{one} modelling assumption that makes a real free kick behave differently from your model, and explain qualitatively its effect on the trajectory.
\end{enumerate}

\courssub{Model answer}

\textbf{Task 1.} The only force acting on the ball (neglecting air resistance) is its weight $\vec{W}=m\vec{g}$. Hence the acceleration is constant: $\vec{a} = \begin{pmatrix} 0 \\ -g \end{pmatrix} = \begin{pmatrix} 0 \\ -10 \end{pmatrix}\ \mathrm{m\,s^{-2}}$. Integrating with respect to time and using the initial velocity $\vec{v}(0) = \begin{pmatrix} 20\cos\theta \\ 20\sin\theta \end{pmatrix}$:
\[ v_{x}(t) = 20\cos\theta, \qquad v_{y}(t) = 20\sin\theta - 10t. \]
Integrating again with initial position $\vec{r}(0)=\vec{0}$:
\[ x(t) = 20t\cos\theta, \qquad y(t) = 20t\sin\theta - 5t^{2}. \]

\textbf{Task 2.} From $x(t)$ we have $t = \dfrac{x}{20\cos\theta}$ (valid for $\theta\neq90^{\circ}$). Substitute into $y(t)$:
\[ y = x\tan\theta - 5\left(\frac{x}{20\cos\theta}\right)^{2} = x\tan\theta - \frac{x^{2}}{80\cos^{2}\theta}. \]
Using the identity $\dfrac{1}{\cos^{2}\theta} = 1 + \tan^{2}\theta$, we obtain the Cartesian equation of the trajectory:
\[ y = x\tan\theta - \frac{x^{2}}{80}\left(1+\tan^{2}\theta\right). \]

\textbf{Task 3.} Let $T = \tan\theta$. The condition to clear the wall is $y(9.1) > 2.0$:
\[ 9.1T - \frac{9.1^{2}}{80}(1+T^{2}) > 2.0. \]
Since $9.1^{2}=82.81$ and $\frac{82.81}{80}=1.035125$, this becomes
\[ 9.1T - 1.035125(1+T^{2}) > 2 \;\Longleftrightarrow\; -1.035125T^{2} + 9.1T - 3.035125 > 0. \]
Multiplying by $-1$ (which reverses the inequality):
\[ 1.035125T^{2} - 9.1T + 3.035125 < 0. \]
The discriminant is $\Delta = 9.1^{2} - 4(1.035125)(3.035125) = 82.81 - 12.570 = 70.240$, so $\sqrt{\Delta} \approx 8.381$. The roots are
\[ T_{1} = \frac{9.1 - 8.381}{2\times1.035125} \approx 0.347, \qquad T_{2} = \frac{9.1 + 8.381}{2\times1.035125} \approx 8.444. \]
The quadratic is negative \emph{between} its roots:
\[ 0.347 < \tan\theta < 8.444 \;\Longleftrightarrow\; 19.1^{\circ} \lesssim \theta \lesssim 83.2^{\circ}. \]

\textbf{Task 4.} The ball must cross the goal line ($x=22$) at a height strictly between $0$ and $2.44\ \mathrm{m}$:
\[ 0 < y(22) < 2.44. \]
With $22^{2}=484$ and $\frac{484}{80}=6.05$, the trajectory equation gives
\[ y(22) = 22T - 6.05(1+T^{2}). \]

\emph{Upper bound (crossbar):} $22T - 6.05(1+T^{2}) < 2.44$
\[ \Longleftrightarrow\; -6.05T^{2} + 22T - 8.49 < 0 \;\Longleftrightarrow\; 6.05T^{2} - 22T + 8.49 > 0. \]
Discriminant $\Delta = 484 - 4(6.05)(8.49) = 484 - 205.458 = 278.542$, $\sqrt{\Delta}\approx 16.690$. Roots:
\[ T_{3} = \frac{22 - 16.690}{12.1} \approx 0.439, \qquad T_{4} = \frac{22 + 16.690}{12.1} \approx 3.197. \]
The quadratic is positive \emph{outside} the roots:
\[ \tan\theta < 0.439 \quad\text{or}\quad \tan\theta > 3.197 \;\Longleftrightarrow\; \theta \lesssim 23.7^{\circ} \quad\text{or}\quad \theta \gtrsim 72.6^{\circ}. \]

\emph{Lower bound (ground):} $22T - 6.05(1+T^{2}) > 0$
\[ \Longleftrightarrow\; 6.05T^{2} - 22T + 6.05 < 0. \]
Discriminant $\Delta = 484 - 4(6.05)^{2} = 484 - 146.41 = 337.59$, $\sqrt{\Delta}\approx 18.374$. Roots:
\[ T_{5} = \frac{22 - 18.374}{12.1} \approx 0.300, \qquad T_{6} = \frac{22 + 18.374}{12.1} \approx 3.337. \]
The quadratic is negative between its roots:
\[ 0.300 < \tan\theta < 3.337 \;\Longleftrightarrow\; 16.7^{\circ} \lesssim \theta \lesssim 73.3^{\circ}. \]

Combining the upper and lower bounds for the crossbar condition:
\[ (0.300 < \tan\theta < 0.439) \quad\text{or}\quad (3.197 < \tan\theta < 3.337). \]

\textbf{Task 5.} Intersect the wall condition ($0.347 < \tan\theta < 8.444$) with the crossbar condition above:
\begin{itemize}
\item Low branch: $\max(0.347,\,0.300) < \tan\theta < \min(8.444,\,0.439) \;\Rightarrow\; 0.347 < \tan\theta < 0.439$.
\item High branch: $\max(0.347,\,3.197) < \tan\theta < \min(8.444,\,3.337) \;\Rightarrow\; 3.197 < \tan\theta < 3.337$.
\end{itemize}
In degrees:
\[ \theta \in [19.1^{\circ},\, 23.7^{\circ}] \quad\text{or}\quad \theta \in [72.6^{\circ},\, 73.3^{\circ}]. \]

The high branch corresponds to a very steep kick: the ball would be in the air for nearly $4\ \mathrm{s}$, giving the goalkeeper ample time to react. It is not a practical direct free kick. The \textbf{physically sensible branch} is the low one: $\theta \in [19.1^{\circ},\, 23.7^{\circ}]$.

An \textbf{optimal angle} is the midpoint of this window, maximising the vertical clearance margins both above the wall and below the crossbar:
\[ \theta_{\text{opt}} \approx \frac{19.1^{\circ}+23.7^{\circ}}{2} = 21.4^{\circ} \approx 21^{\circ}\text{ (to nearest degree)}. \]
We adopt $\theta = 21^{\circ}$ for the remaining computations.

\textbf{Task 6.} For $\theta = 21^{\circ}$: $\cos21^{\circ} \approx 0.9336$, $\sin21^{\circ} \approx 0.3584$.
Time to reach the goal line ($x=22$):
\[ t = \frac{22}{20\cos21^{\circ}} = \frac{22}{18.672} \approx 1.178\ \mathrm{s}. \]
Velocity components at that instant:
\[ v_{x} = 20\cos21^{\circ} \approx 18.67\ \mathrm{m\,s^{-1}}, \]
\[ v_{y} = 20\sin21^{\circ} - 10t \approx 7.168 - 11.78 = -4.61\ \mathrm{m\,s^{-1}}. \]
The negative sign confirms the ball is descending (dipping under the crossbar). The speed is
\[ v = \sqrt{v_{x}^{2} + v_{y}^{2}} = \sqrt{18.67^{2} + 4.61^{2}} = \sqrt{348.6 + 21.3} \approx 19.2\ \mathrm{m\,s^{-1}}. \]
Check height at goal line: $y(1.178) = 20(0.3584)(1.178) - 5(1.178)^{2} \approx 8.44 - 6.94 = 1.50\ \mathrm{m}$, safely below the bar.

\textbf{Task 7.} The vertical velocity is $v_{y}(t) = 7.17 - 10t$ (in $\mathrm{m\,s^{-1}}$). This is a straight line with:
\begin{itemize}
\item slope $-g = -10\ \mathrm{m\,s^{-2}}$,
\item vertical intercept $v_{y}(0) = 7.17\ \mathrm{m\,s^{-1}}$,
\item zero at $t = 0.717\ \mathrm{s}$ (instant of maximum height),
\item value at goal line $t=1.178\ \mathrm{s}$: $v_{y} = -4.61\ \mathrm{m\,s^{-1}}$.
\end{itemize}
The signed area between the graph and the $t$-axis represents the vertical displacement $\Delta y$.

\begin{popfigure}
\begin{center}
\begin{tikzpicture}[scale=0.9,>=stealth]
% Axes
\draw[->] (0,-1.5) -- (0,3.0) node[left]{\small $v_{y}$ (m\,s$^{-1}$)};
\draw[->] (0,0) -- (5.5,0) node[below]{\small $t$ (s)};
% Grid lines (optional, light)
\draw[popline,very thin] (0,0) grid (5.2,2.8);
% Graph: v_y = 7.17 - 10t. Scale: 1 unit = 1 cm for t, 1 m/s = 0.35 cm for v_y? 
% Let's use actual coordinates with a scaling factor for clarity.
% t_max = 1.18, v_max = 7.17, v_min = -4.61.
% Use xscale=3, yscale=0.35
\begin{scope}[xscale=3,yscale=0.35]
\draw[popblue,very thick,domain=0:1.18] plot (\x, {7.17 - 10*\x});
\fill[popink] (0,7.17) circle (0.06) node[above left]{\small $7.17$};
\fill[popink] (0.717,0) circle (0.06) node[below right]{\small $0.72$};
\fill[popink] (1.178,-4.61) circle (0.06) node[below right]{\small $-4.61$};
\draw[dashed,popmuted] (0.717,0) -- (0.717,7.17) node[midway,left]{\small $t_{\text{top}}$};
\draw[dashed,popmuted] (1.178,-4.61) -- (1.178,0) node[below]{\small $1.18$};
\end{scope}
% Tick marks on axes (unscaled)
\foreach \t/\lbl in {0.717/0.72, 1.178/1.18} {
  \draw (\t*3,-0.05) -- (\t*3,0.05) node[below,yshift=-2pt] {\small \lbl};
}
\foreach \v/\lbl in {7.17/7.17, -4.61/-4.61} {
  \draw (-0.05,\v*0.35) -- (0.05,\v*0.35) node[left,xshift=-2pt] {\small \lbl};
}
\end{tikzpicture}
\end{center}
\legende{Vertical component of velocity during the flight of the optimal free kick ($\theta=21^{\circ}$).}
\end{popfigure}

\textbf{Task 8.} Possible answers (any one well argued):
\begin{itemize}
\item \textbf{Air resistance (drag):} reduces the horizontal speed and the maximum height, causing the ball to fall shorter and steeper than predicted. The real scoring window is narrower, and the optimal angle would be slightly lower.
\item \textbf{The Magnus effect (spin):} a spinning ball experiences a force perpendicular to its velocity. Topspin makes the ball dip more sharply, helping it clear the wall and drop under the crossbar — this is exactly what elite free-kick takers exploit. Backspin would make it float, risking hitting the crossbar or going over.
\item \textbf{Finite size of the ball / initial height:} the ball is struck slightly above the ground (on the boot), so $y(0)>0$. This effectively raises the whole trajectory, slightly widening the window.
\item \textbf{Wind:} a headwind increases drag and lift (if spinning), a tailwind does the opposite; crosswind curves the trajectory laterally.
\end{itemize}

\begin{attention}[Marking points to watch]
\begin{itemize}
\item Do not forget the factor $(1+\tan^{2}\theta)$ when replacing $1/\cos^{2}\theta$ in the trajectory equation.
\item A quadratic inequality in $\tan\theta$ is solved by studying the \emph{sign} of the quadratic expression (using a sign table or sketch), not just by stating the roots.
\item Remember that a velocity is a \textbf{vector} — the speed in Task 6 requires both components via Pythagoras.
\item Always check the physical feasibility: the ball must be above ground ($y>0$) at the goal line, not just below the crossbar.
\item In the $v_y$--$t$ graph, the slope is the acceleration ($-g$), the intercept is the initial vertical velocity, and the area under the curve gives the vertical displacement.
\end{itemize}
\end{attention}

\newpage

\coursec{Methods and worked examples}

\coursec{Kinematics of Particles in a Plane}

\courssub{Introduction and Syllabus Overview}
In this chapter we extend the one-dimensional kinematics of Lower Sixth to motion in a plane. We describe the position of a particle by a vector $\vec{r}(t)$, obtain velocity and acceleration by differentiation, and recover them by integration when acceleration is known. The central application is \cle{projectile motion} — motion under uniform gravity with no air resistance — where the horizontal and vertical motions are independent and linked only by the time $t$. We shall derive the standard formulas for range, maximum height and time of flight, solve problems where the launch and landing heights differ, analyse vertical projection (free fall and upward throws), and interpret motion graphs. This covers Lessons 114–122 of the official MINESEC/GCE Board progression.

\begin{savaistu}[Why vectors in kinematics?]
In the Anglophone GCE Advanced Level, vectors $\vec{i},\vec{j}$ are the standard language for planar motion. They allow us to treat the two perpendicular components independently while keeping the physical unity of the motion. This approach is also the gateway to the study of relative motion, circular motion and simple harmonic motion in later chapters.
\end{savaistu}

% ============================================================
% 1. POSITION, VELOCITY AND ACCELERATION AS VECTORS
% ============================================================
\courssub{Position, Velocity and Acceleration as Vectors in a Plane}
\begin{definition}[Position vector]
Let $O$ be a fixed origin in the plane. The \cle{position vector} of a particle $P$ at time $t$ is $\vec{r}(t)=\overrightarrow{OP}$. In Cartesian coordinates with constant unit vectors $\vec{i}$ (horizontal) and $\vec{j}$ (vertical), we write
\[
\vec{r}(t)=x(t)\,\vec{i}+y(t)\,\vec{j},
\]
where $x(t)$ and $y(t)$ are the scalar coordinates of $P$.
\end{definition}

\begin{propriete}[Velocity and acceleration by differentiation]
If $\vec{r}(t)=x(t)\vec{i}+y(t)\vec{j}$ is differentiable, the \cle{velocity} and \cle{acceleration} vectors are
\[
\vec{v}(t)=\frac{d\vec{r}}{dt}=\dot{x}(t)\,\vec{i}+\dot{y}(t)\,\vec{j},\qquad
\vec{a}(t)=\frac{d\vec{v}}{dt}=\ddot{x}(t)\,\vec{i}+\ddot{y}(t)\,\vec{j}.
\]
The \cle{speed} is the magnitude $|\vec{v}|=\sqrt{\dot{x}^2+\dot{y}^2}$; the magnitude of acceleration is $|\vec{a}|=\sqrt{\ddot{x}^2+\ddot{y}^2}$.
\end{propriete}

\begin{methode}[Differentiating a position vector]
\begin{enumerate}
\item Write $\vec{r}(t)=x(t)\vec{i}+y(t)\vec{j}$; identify $x(t)$ and $y(t)$.
\item Differentiate each component with respect to $t$ to obtain $\vec{v}(t)=\dot{x}\vec{i}+\dot{y}\vec{j}$.
\item Differentiate again to obtain $\vec{a}(t)=\ddot{x}\vec{i}+\ddot{y}\vec{j}$.
\item For speed or magnitude of acceleration, compute the square root of the sum of squares of the components.
\item To find when the motion is parallel to an axis, set the other component of $\vec{v}$ to zero and solve for $t$.
\end{enumerate}
\textbf{Reflex:} $\vec{i}$ and $\vec{j}$ are constant vectors, so only the scalar components are differentiated.
\end{methode}

\begin{exoresolu}[Velocity and acceleration from $\vec{r}(t)$]
A particle $P$ moves in the plane so that its position vector at time $t$ seconds is
\[
\vec{r}(t)=\left(2t^3-3t\right)\vec{i}+\left(t^2+4t\right)\vec{j}\ \text{metres}.
\]
\begin{itemize}
\item[(a)] Find the velocity of $P$ at $t=2$.
\item[(b)] Find the acceleration of $P$ and show that it is never zero.
\item[(c)] Find the speed of $P$ when $t=1$.
\end{itemize}
\tcblower
\textbf{(a)} Differentiate each component:
\[
\vec{v}(t)=\frac{d\vec{r}}{dt}=\left(6t^2-3\right)\vec{i}+\left(2t+4\right)\vec{j}.
\]
At $t=2$: $\vec{v}(2)=(6\times4-3)\vec{i}+(2\times2+4)\vec{j}=21\vec{i}+8\vec{j}\ \mathrm{m\,s^{-1}}$.

\textbf{(b)} Differentiate again:
\[
\vec{a}(t)=\frac{d\vec{v}}{dt}=12t\,\vec{i}+2\vec{j}.
\]
The $\vec{j}$-component is the constant $2\neq 0$, so $\vec{a}(t)$ can never be the zero vector. Hence the acceleration is never zero.

\textbf{(c)} At $t=1$: $\vec{v}(1)=(6-3)\vec{i}+(2+4)\vec{j}=3\vec{i}+6\vec{j}$.
\[
|\vec{v}(1)|=\sqrt{3^2+6^2}=\sqrt{45}=3\sqrt{5}\approx 6.71\ \mathrm{m\,s^{-1}}.
\]
\end{exoresolu}

% ============================================================
% 2. INTEGRATION IN KINEMATICS
% ============================================================
\courssub{From Acceleration back to Velocity and Position: Integration in Kinematics}
When the acceleration (or velocity) is given as a function of time, we recover the velocity (or position) by integration. Each integration introduces a constant vector which is determined by the initial conditions.

\begin{methode}[Integrating with initial conditions]
Given $\vec{a}(t)$ (or $\vec{v}(t)$) and initial values $\vec{v}(t_0)$, $\vec{r}(t_0)$:
\begin{enumerate}
\item Integrate each component of $\vec{a}$ with respect to $t$ to obtain $\vec{v}$; a constant of integration appears in \emph{each} component (write them as $c_1\vec{i}+c_2\vec{j}$ or a constant vector $\vec{c}_1$).
\item Use the given velocity at a known time (often $t=0$) to fix the constants.
\item Integrate $\vec{v}$ component by component to obtain $\vec{r}$; use the initial position to fix the new constants.
\item Write the final answers as single vectors in $\vec{i},\vec{j}$ form.
\end{enumerate}
\textbf{Key idea:} $\displaystyle \vec{v}=\int\vec{a}\,dt+\vec{c}_1,\qquad \vec{r}=\int\vec{v}\,dt+\vec{c}_2.$ Never forget the constant vectors: they carry the physical initial conditions.
\end{methode}

\begin{exoresolu}[From acceleration back to position]
A particle moves in a plane with acceleration $\vec{a}(t)=(6t)\vec{i}-2\vec{j}\ \mathrm{m\,s^{-2}}$. Initially ($t=0$) the particle is at the point with position vector $3\vec{j}$ metres and is moving with velocity $4\vec{i}\ \mathrm{m\,s^{-1}}$.
\begin{itemize}
\item[(a)] Find the velocity of the particle at time $t$.
\item[(b)] Find the position vector of the particle at time $t$.
\end{itemize}
\tcblower
\textbf{(a)} Integrate the acceleration component by component:
\[
\vec{v}(t)=\int\vec{a}\,dt=\left(3t^2+c_1\right)\vec{i}+\left(-2t+c_2\right)\vec{j}.
\]
At $t=0$, $\vec{v}=4\vec{i}$, so $c_1=4$ and $c_2=0$. Hence
\[
\vec{v}(t)=\left(3t^2+4\right)\vec{i}-2t\,\vec{j}\ \mathrm{m\,s^{-1}}.
\]

\textbf{(b)} Integrate the velocity:
\[
\vec{r}(t)=\int\vec{v}\,dt=\left(t^3+4t+k_1\right)\vec{i}+\left(-t^2+k_2\right)\vec{j}.
\]
At $t=0$, $\vec{r}=3\vec{j}$, so $k_1=0$ and $k_2=3$. Therefore
\[
\vec{r}(t)=\left(t^3+4t\right)\vec{i}+\left(3-t^2\right)\vec{j}\ \text{metres}.
\]
\textbf{Check:} differentiating $\vec{r}$ twice returns $6t\vec{i}-2\vec{j}=\vec{a}$. \checkmark
\end{exoresolu}

% ============================================================
% 3. PROJECTILE MOTION: DEFINITION, MODEL AND EQUATIONS
% ============================================================
\courssub{Projectile Motion: Definition, Model and Equations}
\begin{definition}[Projectile]
A \cle{projectile} is a particle projected in a vertical plane and moving under the action of gravity alone. The standard model assumes:
\begin{itemize}
\item The acceleration due to gravity $\vec{g}$ is constant in magnitude ($g\approx 9.8\ \mathrm{m\,s^{-2}}$ or $10\ \mathrm{m\,s^{-2}}$ as instructed) and direction (vertically downwards).
\item Air resistance is negligible.
\item The Earth's curvature and rotation are ignored over the range of motion.
\end{itemize}
\end{definition}

\begin{propriete}[Equations of motion of a projectile launched from the origin]
A projectile is launched from the origin $O$ with initial speed $u$ at an angle $\alpha$ above the horizontal. Take $\vec{i}$ horizontal (positive in the direction of projection) and $\vec{j}$ vertically upward. Then $\vec{a}=-g\vec{j}$.
Resolving the initial velocity: $u_x=u\cos\alpha$, $u_y=u\sin\alpha$.
\begin{align*}
\text{Horizontal motion (no acceleration):}\quad & x(t)=u\cos\alpha\cdot t, \quad v_x=u\cos\alpha \ (\text{constant}).\\
\text{Vertical motion (constant acceleration $-g$):}\quad & y(t)=u\sin\alpha\cdot t-\tfrac12 gt^2, \quad v_y=u\sin\alpha-gt.
\end{align*}
The \cle{trajectory equation} (Cartesian equation) is obtained by eliminating $t$:
\[
y=x\tan\alpha-\frac{gx^2}{2u^2\cos^2\alpha}=x\tan\alpha-\frac{gx^2}{2u^2}\left(1+\tan^2\alpha\right).
\]
This is a parabola with vertical axis, opening downwards.
\end{propriete}

\begin{aretenir}[Independence of motions]
The horizontal and vertical motions are \emph{independent}; only the time $t$ links them. The horizontal velocity remains constant throughout the flight; the vertical velocity changes exactly as in one-dimensional motion under constant acceleration $-g$.
\end{aretenir}

\begin{exoresolu}[Trajectory of a stone]
A stone is thrown from the top of a cliff with speed $20\ \mathrm{m\,s^{-1}}$ at an angle of $30^{\circ}$ above the horizontal. Take $g=10\ \mathrm{m\,s^{-2}}$ and the point of projection as origin.
\begin{itemize}
\item[(a)] Write down $x(t)$ and $y(t)$.
\item[(b)] Find the Cartesian equation of the trajectory.
\item[(c)] Find the position and the velocity of the stone at $t=1\ \mathrm{s}$.
\end{itemize}
\tcblower
\textbf{(a)} $u\cos30^{\circ}=20\times\frac{\sqrt3}{2}=10\sqrt3$, $u\sin30^{\circ}=20\times\frac12=10$. Hence
\[
x(t)=10\sqrt3\,t,\qquad y(t)=10t-5t^2.
\]

\textbf{(b)} From $x=10\sqrt3\,t$, $t=\dfrac{x}{10\sqrt3}$. Substituting:
\[
y=10\cdot\frac{x}{10\sqrt3}-5\left(\frac{x}{10\sqrt3}\right)^2=\frac{x}{\sqrt3}-\frac{x^2}{60}.
\]
This is a parabola opening downwards, as expected.

\textbf{(c)} At $t=1$: $x=10\sqrt3\approx 17.3\ \mathrm{m}$, $y=10-5=5\ \mathrm{m}$. So the stone is at $(17.3\,;\,5)\ \mathrm{m}$.
Velocity: $v_x=10\sqrt3\ \mathrm{m\,s^{-1}}$ (constant), $v_y=10-10\times1=0$. The vertical velocity is zero: at $t=1\ \mathrm{s}$ the stone is at the \textbf{top of its path}, moving horizontally with speed $10\sqrt3\approx 17.3\ \mathrm{m\,s^{-1}}$.
\end{exoresolu}

% ============================================================
% 4. RANGE, MAXIMUM HEIGHT AND TIME OF FLIGHT
% ============================================================
\courssub{Range, Maximum Height and Time of Flight}
The following formulas apply \emph{only} when the projectile is launched from and lands on the same horizontal level (level ground).

\begin{propriete}[Standard projectile formulas (level ground)]
For a projectile launched from level ground at speed $u$, angle $\alpha$ ($0<\alpha<90^{\circ}$):
\begin{enumerate}
\item \cle{Time of flight}: $\displaystyle T=\frac{2u\sin\alpha}{g}$ (solve $y=0$, $t\neq 0$).
\item \cle{Maximum height}: $\displaystyle H=\frac{u^2\sin^2\alpha}{2g}$ (set $v_y=0$ and use $v_y^2=u_y^2-2gH$).
\item \cle{Range}: $\displaystyle R=\frac{u^2\sin 2\alpha}{g}$ (compute $x(T)=u\cos\alpha\cdot T$).
\item \cle{Maximum range} for a given $u$: $R_{\max}=\dfrac{u^2}{g}$ at $\alpha=45^{\circ}$.
\item \cle{Complementary angles} $\alpha$ and $90^{\circ}-\alpha$ give the \emph{same} range because $\sin(2(90^{\circ}-\alpha))=\sin(180^{\circ}-2\alpha)=\sin2\alpha$.
\end{enumerate}
\end{propriete}

\begin{methode}[Using the standard projectile formulas]
\begin{enumerate}
\item Verify that launch and landing are at the same height. If not, go back to $x(t)$, $y(t)$.
\item Identify $u$, $\alpha$, $g$.
\item Substitute into $T$, $H$, $R$ formulas.
\item For complementary angle problems, state $90^{\circ}-\alpha$ and note the different $T$ and $H$.
\end{enumerate}
\textbf{Warning:} these formulas require launch and landing at the \emph{same} height. Otherwise return to the equations of motion.
\end{methode}

\begin{exoresolu}[A footballer's kick]
A footballer kicks a ball with speed $25\ \mathrm{m\,s^{-1}}$ at $40^{\circ}$ above the horizontal from level ground. Take $g=10\ \mathrm{m\,s^{-2}}$, $\sin40^{\circ}\approx 0.643$, $\cos40^{\circ}\approx 0.766$, $\sin80^{\circ}\approx 0.985$.
\begin{itemize}
\item[(a)] Find the time of flight.
\item[(b)] Find the range.
\item[(c)] Find the greatest height reached.
\item[(d)] At what other angle, with the same speed, would the ball have the same range?
\end{itemize}
\tcblower
\textbf{(a)} $T=\dfrac{2u\sin\alpha}{g}=\dfrac{2\times25\times0.643}{10}=3.22\ \mathrm{s}$.

\textbf{(b)} $R=\dfrac{u^2\sin2\alpha}{g}=\dfrac{625\times0.985}{10}\approx 61.6\ \mathrm{m}$.

\textbf{(c)} $H=\dfrac{u^2\sin^2\alpha}{2g}=\dfrac{625\times(0.643)^2}{20}=\dfrac{625\times0.4134}{20}\approx 12.9\ \mathrm{m}$.

\textbf{(d)} The complementary angle $90^{\circ}-40^{\circ}=50^{\circ}$ gives the same range, since $\sin(2\times50^{\circ})=\sin100^{\circ}=\sin80^{\circ}$. The ball would then go higher ($H\approx 18.3\ \mathrm{m}$) but stay longer in the air ($T\approx 3.83\ \mathrm{s}$).
\end{exoresolu}

% ============================================================
% 5. APPLICATIONS: NON-LEVEL GROUND AND TARGET PROBLEMS
% ============================================================
\courssub{Applications and Problem-Solving with Projectiles}
When the launch and landing heights differ, or a specific target must be hit, the standard formulas do not apply. We must work from the parametric equations or the trajectory equation.

\begin{methode}[Non-level problems: work from the equations]
\begin{enumerate}
\item Choose axes and write $x(t)=u\cos\alpha\, t$ and $y(t)=u\sin\alpha\, t-\tfrac12 gt^2$ with the correct \emph{signs} for the initial position (e.g. if launched from a height $h$ above the landing level, set $y(0)=h$ or $y(0)=0$ and landing at $y=-h$).
\item Impose the condition of the problem (e.g. $y=-h$ at landing, or the trajectory passes through $(x_0,y_0)$).
\item Solve the resulting equation — often a quadratic in $t$ or in $\tan\alpha$; keep the physically meaningful root ($t>0$, $\alpha$ acute usually).
\item For "passes through $(x_0,y_0)$" problems, use the trajectory equation $y_0=x_0\tan\alpha-\dfrac{gx_0^2}{2u^2}(1+\tan^2\alpha)$ to get a quadratic in $\tan\alpha$.
\end{enumerate}
\textbf{Reflex:} a quadratic in $\tan\alpha$ with two positive roots means the target can be hit by \emph{two} possible throws (a low one and a high one).
\end{methode}

\begin{exoresolu}[Ball thrown from a balcony]
From a balcony $15\ \mathrm{m}$ above the ground, a student throws a ball with speed $20\ \mathrm{m\,s^{-1}}$ at $30^{\circ}$ above the horizontal. Take $g=10\ \mathrm{m\,s^{-2}}$.
\begin{itemize}
\item[(a)] Find the time at which the ball hits the ground.
\item[(b)] Find the horizontal distance travelled before landing.
\item[(c)] Find the speed of the ball just before impact.
\end{itemize}
\tcblower
Take the point of projection as origin, $y$ upwards. Then $u_x=20\cos30^{\circ}=10\sqrt3$, $u_y=20\sin30^{\circ}=10$, and the ground is at $y=-15$.

\textbf{(a)} Solve $y(t)=-15$:
\[
10t-5t^2=-15\ \Longrightarrow\ 5t^2-10t-15=0\ \Longrightarrow\ t^2-2t-3=0.
\]
Factorising: $(t-3)(t+1)=0$, so $t=3\ \mathrm{s}$ (the root $t=-1$ is rejected).

\textbf{(b)} $x(3)=10\sqrt3\times3=30\sqrt3\approx 52.0\ \mathrm{m}$.

\textbf{(c)} At $t=3$: $v_x=10\sqrt3\ \mathrm{m\,s^{-1}}$, $v_y=10-10\times3=-20\ \mathrm{m\,s^{-1}}$.
\[
|\vec{v}|=\sqrt{(10\sqrt3)^2+(-20)^2}=\sqrt{300+400}=\sqrt{700}=10\sqrt7\approx 26.5\ \mathrm{m\,s^{-1}}.
\]
Note that the impact speed exceeds the launch speed: the ball has gained kinetic energy by falling $15\ \mathrm{m}$.
\end{exoresolu}

\begin{exoresolu}[Hitting a target on a wall]
A particle is projected from ground level with speed $25\ \mathrm{m\,s^{-1}}$ towards a vertical wall $30\ \mathrm{m}$ away. The particle just passes over the top of the wall which is $10\ \mathrm{m}$ high. Find the two possible angles of projection. Take $g=10\ \mathrm{m\,s^{-2}}$.
\tcblower
The trajectory must pass through $(30,10)$. Using $y=x\tan\alpha-\dfrac{gx^2}{2u^2}(1+\tan^2\alpha)$:
\[
10=30\tan\alpha-\frac{10\times30^2}{2\times25^2}(1+\tan^2\alpha)=30\tan\alpha-\frac{9000}{1250}(1+\tan^2\alpha)=30\tan\alpha-7.2(1+\tan^2\alpha).
\]
Rearrange: $7.2\tan^2\alpha-30\tan\alpha+17.2=0$. Multiply by 5: $36\tan^2\alpha-150\tan\alpha+86=0$, divide by 2: $18\tan^2\alpha-75\tan\alpha+43=0$.
\[
\tan\alpha=\frac{75\pm\sqrt{5625-3096}}{36}=\frac{75\pm\sqrt{2529}}{36}\approx\frac{75\pm50.29}{36}.
\]
$\tan\alpha_1\approx 3.48 \Rightarrow \alpha_1\approx 74^{\circ}$ (high trajectory), $\tan\alpha_2\approx 0.686 \Rightarrow \alpha_2\approx 34.5^{\circ}$ (low trajectory). Both are valid.
\end{exoresolu}

% ============================================================
% 6. FREE FALL AND VERTICAL PROJECTION
% ============================================================
\courssub{Free Fall, Vertical Projection and Motion Graphs}
When the motion is purely vertical, we have one-dimensional motion with constant acceleration $\pm g$. The vector notation reduces to a sign convention.

\begin{methode}[One-dimensional vertical motion]
All motion is on a vertical line; choose an origin and a positive direction, then use the constant-acceleration formulas with $a=\pm g$:
\[
v=u+at,\qquad s=ut+\tfrac12 at^2,\qquad v^2=u^2+2as.
\]
\begin{enumerate}
\item \textbf{Dropped object (free fall):} $u=0$; taking downwards positive, $v=gt$, $h=\tfrac12 gt^2$.
\item \textbf{Thrown upwards:} take upwards positive, $a=-g$. At the highest point $v=0$.
\item Time to the top: $t_{\text{up}}=\dfrac{u}{g}$; greatest height: $H=\dfrac{u^2}{2g}$.
\item Returning to the launch level, the object arrives with speed $u$ (same speed, opposite direction) after total time $\dfrac{2u}{g}$.
\end{enumerate}
\textbf{Warning:} fix the sign convention \emph{before} substituting, and keep it throughout the whole solution.
\end{methode}

\begin{exoresolu}[Stone thrown upwards from a well]
A stone is thrown vertically upwards from the top of a well with speed $14\ \mathrm{m\,s^{-1}}$. It falls into the well and reaches the water surface $20\ \mathrm{m}$ below the point of projection. Take $g=9.8\ \mathrm{m\,s^{-2}}$.
\begin{itemize}
\item[(a)] Find the greatest height reached above the point of projection.
\item[(b)] Find the total time of flight until the stone hits the water.
\item[(c)] Find the speed of impact with the water.
\end{itemize}
\tcblower
Take upwards as positive, origin at the point of projection; $u=14$, $a=-9.8$.

\textbf{(a)} At the top, $v=0$: $v^2=u^2+2as$ gives $0=14^2-2\times9.8\times H$, so
\[
H=\frac{196}{19.6}=10\ \text{m}.
\]

\textbf{(b)} The water is at $s=-20\ \mathrm{m}$. Using $s=ut+\tfrac12 at^2$:
\[
-20=14t-4.9t^2\ \Longrightarrow\ 4.9t^2-14t-20=0.
\]
By the quadratic formula:
\[
t=\frac{14\pm\sqrt{196+392}}{9.8}=\frac{14\pm\sqrt{588}}{9.8}\approx\frac{14\pm24.25}{9.8}.
\]
Keeping the positive root: $t\approx\dfrac{38.25}{9.8}\approx 3.90\ \mathrm{s}$.

\textbf{(c)} $v=u+at=14-9.8\times3.90\approx-24.2\ \mathrm{m\,s^{-1}}$. The impact speed is about $24.2\ \mathrm{m\,s^{-1}}$ downwards. \emph{Check with} $v^2=u^2+2as=196-19.6\times(-20)=588$, so $|v|=\sqrt{588}\approx 24.2\ \mathrm{m\,s^{-1}}$. \checkmark
\end{exoresolu}

% ============================================================
% 7. GRAPHICAL REPRESENTATION OF VERTICAL MOTION
% ============================================================
\courssub{Graphical Representation of Vertical Motion}
For an object thrown vertically (acceleration $-g$, constant), the graphs are powerful tools for visualising and solving problems.

\begin{methode}[Graphs of vertical motion]
For an object thrown vertically upwards with initial speed $u$ from a reference level:
\begin{enumerate}
\item \textbf{Velocity--time graph:} a straight line of gradient $-g$, starting at $v=u$. The area between the line and the time axis gives the displacement (areas below the axis count negatively).
\item \textbf{Displacement--time graph:} a parabola $s=ut-\tfrac12 gt^2$; its gradient at any point is the velocity; the vertex corresponds to the highest point ($v=0$).
\item \textbf{Speed--time graph:} the velocity graph \emph{reflected} above the axis where $v<0$ — it is V-shaped for an up--down flight. Speed is never negative.
\item Read off: time at the top (where $v=0$), landing time (where $s$ returns to its initial value), impact velocity (final value of $v$).
\end{enumerate}
\textbf{Reflex:} gradient of $s$--$t$ graph $=$ velocity; gradient of $v$--$t$ graph $=$ acceleration; area under $v$--$t$ graph $=$ displacement.
\end{methode}

\begin{exoresolu}[Graphs for a ball thrown upwards]
A ball is thrown vertically upwards with speed $20\ \mathrm{m\,s^{-1}}$ from ground level and is caught again at the same level. Take $g=10\ \mathrm{m\,s^{-2}}$.
\begin{itemize}
\item[(a)] Sketch the velocity--time graph for the whole flight.
\item[(b)] Use the graph to find the total time of flight and the greatest height.
\item[(c)] Describe the displacement--time graph.
\end{itemize}
\tcblower
\textbf{(a)} $v=20-10t$: a straight line from $(0,20)$ with gradient $-10$, crossing the axis at $t=2$ and reaching $v=-20$ at $t=4$.

\begin{popfigure}
\begin{center}
\begin{tikzpicture}[scale=0.9]
\draw[->,thick] (0,-2.6)--(0,2.8) node[above]{$v$ (m\,s$^{-1}$)};
\draw[->,thick] (0,0)--(4.8,0) node[right]{$t$ (s)};
\draw[very thick,popblue] (0,2)--(4,-2);
\fill[popblueL,opacity=0.6] (0,0)--(0,2)--(2,0)--cycle;
\fill[poppink,opacity=0.35] (2,0)--(4,-2)--(4,0)--cycle;
\foreach \x in {1,2,3,4}{\draw (\x,0.06)--(\x,-0.06) node[below=1pt]{\small $\x$};}
\draw (0.06,2)--(-0.06,2) node[left=1pt]{\small $20$};
\draw (0.06,-2)--(-0.06,-2) node[left=1pt]{\small $-20$};
\node[popdark] at (1,1.4) {\small $+A$};
\node[popdark] at (3.4,-1.2) {\small $-A$};
\end{tikzpicture}
\legende{Velocity--time graph of the ball: gradient $-g=-10\ \mathrm{m\,s^{-2}}$.}
\end{center}
\end{popfigure}

\textbf{(b)} The ball returns when the \emph{net} area is zero. The two triangles have equal area $A=\tfrac12\times2\times20=20$, one positive, one negative, so the displacement is zero at $t=4\ \mathrm{s}$: the total flight time is $4\ \mathrm{s}$. The greatest height is the positive area up to $t=2$: $H=20\ \mathrm{m}$. (Check: $H=\dfrac{u^2}{2g}=\dfrac{400}{20}=20\ \mathrm{m}$. \checkmark)

\textbf{(c)} The displacement--time graph is the parabola $s=20t-5t^2$, starting at the origin, rising to a maximum of $20\ \mathrm{m}$ at $t=2$ (where its gradient is zero, matching $v=0$), and returning to $s=0$ at $t=4\ \mathrm{s}$. Its gradient decreases steadily, reflecting the constant acceleration $-10\ \mathrm{m\,s^{-2}}$.
\end{exoresolu}

% ============================================================
% COMMON MISTAKES
% ============================================================
\begin{attention}[Frequent errors to avoid]
\begin{itemize}
\item Mixing sign conventions inside one calculation: decide the positive direction first and write it down.
\item Using $R=\dfrac{u^2\sin2\alpha}{g}$ when the landing point is not at the launch height — return to $x(t)$, $y(t)$.
\item Forgetting that $v_x$ is constant: at the top of a projectile path the velocity is \emph{not} zero, only $v_y$ is.
\item Confusing speed and velocity on graphs: speed is $|\vec{v}|$ and is never negative.
\item Dropping the constants of integration when working back from acceleration: they are fixed by the initial conditions.
\item Using $g=9.8$ and $g=10$ in the same problem; follow the instruction in the question.
\end{itemize}
\end{attention}

% ============================================================
% CONSOLIDATION EXERCISES
% ============================================================
\courssub{Consolidation Exercises}

\begin{exercice}[Exercise 1: Vector kinematics]
The position vector of a particle $P$ at time $t$ seconds is $\vec{r}(t)=(t^3-6t)\vec{i}+(4t^2-8t)\vec{j}$ metres.
\begin{enumerate}
\item Find the velocity and acceleration vectors at time $t$.
\item Find the time when $P$ is moving parallel to the $\vec{j}$-axis.
\item Find the speed of $P$ when $t=2$.
\end{enumerate}
\end{exercice}

\begin{exercice}[Exercise 2: Integration from acceleration]
A particle moves in the plane with acceleration $\vec{a}(t)=(4t-2)\vec{i}+6\vec{j}\ \mathrm{m\,s^{-2}}$. At $t=0$, its velocity is $3\vec{i}-2\vec{j}\ \mathrm{m\,s^{-1}}$ and its position vector is $5\vec{i}+4\vec{j}$ metres. Find expressions for $\vec{v}(t)$ and $\vec{r}(t)$.
\end{exercice}

\begin{exercice}[Exercise 3: Projectile on level ground]
A projectile is fired from ground level with speed $30\ \mathrm{m\,s^{-1}}$ at an angle $\alpha$ where $\sin\alpha=0.6$ and $\cos\alpha=0.8$. Take $g=10\ \mathrm{m\,s^{-2}}$.
\begin{enumerate}
\item Find the time of flight, the range and the maximum height.
\item Find the velocity of the projectile $2\ \mathrm{s}$ after projection.
\end{enumerate}
\end{exercice}

\begin{exercice}[Exercise 4: Projectile from a height]
A stone is thrown horizontally with speed $15\ \mathrm{m\,s^{-1}}$ from the top of a vertical cliff $80\ \mathrm{m}$ high. Take $g=10\ \mathrm{m\,s^{-2}}$.
\begin{enumerate}
\item Find the time taken to reach the sea.
\item Find the horizontal distance from the foot of the cliff where the stone hits the sea.
\item Find the speed and direction of motion at impact.
\end{enumerate}
\end{exercice}

\begin{exercice}[Exercise 5: Vertical motion graphs]
A ball is projected vertically upwards with speed $25\ \mathrm{m\,s^{-1}}$ from a point $5\ \mathrm{m}$ above the ground. Take $g=10\ \mathrm{m\,s^{-2}}$.
\begin{enumerate}
\item Sketch the velocity--time graph until the ball hits the ground.
\item Use the graph to find the total time of flight and the maximum height above the ground.
\item Sketch the corresponding displacement--time graph (displacement measured from the ground).
\end{enumerate}
\end{exercice}

\begin{corrige}[Exercise 1]
\begin{enumerate}
\item $\vec{v}(t)=(3t^2-6)\vec{i}+(8t-8)\vec{j}$; $\vec{a}(t)=6t\vec{i}+8\vec{j}$.
\item Parallel to $\vec{j}$-axis $\Rightarrow v_x=0 \Rightarrow 3t^2-6=0 \Rightarrow t=\sqrt{2}\ \mathrm{s}$ ($t>0$).
\item At $t=2$: $\vec{v}=(6)\vec{i}+(8)\vec{j}$; speed $=\sqrt{36+64}=10\ \mathrm{m\,s^{-1}}$.
\end{enumerate}
\end{corrige}

\begin{corrige}[Exercise 2]
$\vec{v}(t)=\int\vec{a}\,dt=(2t^2-2t+3)\vec{i}+(6t-2)\vec{j}$.
$\vec{r}(t)=\int\vec{v}\,dt=\left(\frac{2}{3}t^3-t^2+3t+5\right)\vec{i}+(3t^2-2t+4)\vec{j}$.
\end{corrige}

\begin{corrige}[Exercise 3]
$u=30$, $\sin\alpha=0.6$, $\cos\alpha=0.8$, $g=10$.
$T=\frac{2\times30\times0.6}{10}=3.6\ \mathrm{s}$.
$R=\frac{30^2\times2\times0.6\times0.8}{10}=86.4\ \mathrm{m}$.
$H=\frac{30^2\times0.6^2}{20}=16.2\ \mathrm{m}$.
At $t=2$: $v_x=24\ \mathrm{m\,s^{-1}}$, $v_y=18-20=-2\ \mathrm{m\,s^{-1}}$; $\vec{v}=24\vec{i}-2\vec{j}\ \mathrm{m\,s^{-1}}$.
\end{corrige}

\begin{corrige}[Exercise 4]
Horizontal throw $\Rightarrow \alpha=0$, $u_x=15$, $u_y=0$. Origin at cliff top, $y$ up, ground at $y=-80$.
$y(t)=-5t^2=-80 \Rightarrow t=4\ \mathrm{s}$.
$x(4)=15\times4=60\ \mathrm{m}$.
At $t=4$: $v_x=15$, $v_y=-40$; speed $=\sqrt{225+1600}=\sqrt{1825}\approx 42.7\ \mathrm{m\,s^{-1}}$; angle below horizontal $\theta=\arctan(40/15)\approx 69.4^{\circ}$.
\end{corrige}

\begin{corrige}[Exercise 5]
Upwards positive, origin at ground. Initial $s=5$, $u=25$, $a=-10$.
$v=25-10t$. Hits ground when $s=0$: $5+25t-5t^2=0 \Rightarrow t^2-5t-1=0 \Rightarrow t=\frac{5+\sqrt{29}}{2}\approx 5.19\ \mathrm{s}$.
Max height: $v=0 \Rightarrow t=2.5\ \mathrm{s}$, $s_{\max}=5+25(2.5)-5(2.5)^2=36.25\ \mathrm{m}$.
$v$-$t$ graph: line from $(0,25)$ to $(5.19,-26.9)$, crossing $t$-axis at $2.5$.
$s$-$t$ graph: parabola from $(0,5)$ to $(5.19,0)$, vertex at $(2.5, 36.25)$.
\end{corrige}

\newpage

\coursec{Exercises}

\courssub{I check my knowledge}

\begin{exercice}[Multiple choice questions]
For each question, choose the single correct answer, giving a brief justification.
\begin{enumerate}
\item A particle has position vector $\vec{r}(t) = (3t^2)\vec{i} + (4t - t^2)\vec{j}$ metres. Its velocity at time $t$ is:
\begin{enumerate}
\item $(6t)\vec{i} + (4 - 2t)\vec{j}$
\item $(3t^2)\vec{i} + (4 - 2t)\vec{j}$
\item $(6t)\vec{i} + (4t - 2t)\vec{j}$
\item $(6)\vec{i} + (-2)\vec{j}$
\end{enumerate}
\item A projectile is launched from level ground. At the highest point of its trajectory:
\begin{enumerate}
\item its velocity and acceleration are both zero;
\item its velocity is zero but its acceleration is $g$ downwards;
\item the horizontal component of its velocity is zero;
\item the vertical component of its velocity is zero and its acceleration is $g$ downwards.
\end{enumerate}
\item For a projectile launched at speed $u$ at angle $\theta$ on level ground, the range is:
\begin{enumerate}
\item $\dfrac{u^2\sin\theta}{g}$
\item $\dfrac{u^2\sin 2\theta}{g}$
\item $\dfrac{u^2\sin^2\theta}{2g}$
\item $\dfrac{2u\sin\theta}{g}$
\end{enumerate}
\item A stone is dropped from rest from a height of \SI{20}{m}. Taking $g = \SI{10}{m.s^{-2}}$, the time taken to reach the ground is:
\begin{enumerate}
\item \SI{1}{s}
\item $\sqrt{2}\ \mathrm{s}$
\item \SI{2}{s}
\item \SI{4}{s}
\end{enumerate}
\end{enumerate}
\end{exercice}

\begin{exercice}[True or false?]
State whether each assertion is true or false, justifying your answer carefully.
\begin{enumerate}
\item If the acceleration of a particle is constant, then its velocity is constant.
\item Two projectiles launched from the same point on level ground with the same speed at complementary angles $\theta$ and $90^\circ - \theta$ have the same range.
\item A ball thrown vertically upwards has zero acceleration at the top of its flight.
\item The horizontal component of the velocity of a projectile remains constant throughout the flight (air resistance neglected).
\end{enumerate}
\end{exercice}

\begin{exercice}[Key definitions]
\begin{enumerate}
\item Define the \cle{velocity} and the \cle{acceleration} of a particle whose position vector is $\vec{r}(t)$.
\item State the two components of the acceleration of a projectile in the standard model (no air resistance), and justify them.
\item Define the \cle{range}, the \cle{time of flight} and the \cle{maximum height} of a projectile.
\item What is meant by \cle{free fall}? State the value of the acceleration of a body in free fall near the Earth's surface.
\end{enumerate}
\end{exercice}

\begin{exercice}[Direct calculations]
Take $g = \SI{10}{m.s^{-2}}$ throughout.
\begin{enumerate}
\item A particle has velocity $\vec{v}(t) = (4t)\vec{i} + (3 - 2t)\vec{j}\ \mathrm{m\,s^{-1}}$. Find its acceleration and its speed at $t = 2$.
\item A stone is dropped from rest from a bridge and hits the water \SI{3}{s} later. Find the height of the bridge above the water and the speed of impact.
\item A projectile is launched at \SI{20}{m.s^{-1}} at $30^\circ$ above the horizontal. Find the horizontal and vertical components of its initial velocity.
\item A ball is thrown vertically upwards at \SI{15}{m.s^{-1}}. Find the greatest height reached and the time taken to reach it.
\end{enumerate}
\end{exercice}

\courssub{I practise}

\begin{exercice}[Differentiating a position vector]
A particle moves in a plane with position vector
\[ \vec{r}(t) = (2t^2 - 3t)\vec{i} + (t^3 - 6t)\vec{j} \quad \text{metres}. \]
\begin{enumerate}
\item Find its velocity $\vec{v}(t)$ and acceleration $\vec{a}(t)$.
\item Find its velocity and acceleration at $t = 2$, and its speed at that instant.
\item Find the times when the particle moves parallel to the $y$-axis.
\item Find the Cartesian equation of its path.
\end{enumerate}
\end{exercice}

\begin{exercice}[Integrating an acceleration]
A particle starts from the origin $O$ with velocity $(4\vec{i} + 3\vec{j})\ \mathrm{m\,s^{-1}}$ and moves with acceleration
\[ \vec{a}(t) = (2t\,\vec{i} - 6\,\vec{j})\ \mathrm{m\,s^{-2}}. \]
\begin{enumerate}
\item Find its velocity vector at time $t$.
\item Find its position vector at time $t$.
\item Find its distance from $O$ when $t = 3$.
\end{enumerate}
\end{exercice}

\begin{exercice}[Standard projectile on level ground]
A stone is projected at \SI{25}{m.s^{-1}} at $40^\circ$ above the horizontal from level ground. Taking $g = \SI{10}{m.s^{-2}}$, find:
\begin{enumerate}
\item the time of flight;
\item the maximum height reached;
\item the range on the horizontal plane through the point of projection.
\end{enumerate}
(Give your answers to 3 significant figures; you may use $\sin 40^\circ \approx 0.643$ and $\cos 40^\circ \approx 0.766$.)
\end{exercice}

\begin{exercice}[Horizontal projection from a cliff]
A ball is thrown horizontally at \SI{12}{m.s^{-1}} from the top of a cliff \SI{45}{m} high, and falls into the sea. Taking $g = \SI{10}{m.s^{-2}}$, find:
\begin{enumerate}
\item the time of flight;
\item the horizontal distance from the foot of the cliff to the point of impact;
\item the velocity (magnitude and direction) of the ball on hitting the water.
\end{enumerate}
\end{exercice}

\courssub{I prepare for the GCE}

\begin{exercice}[Clearing the wall and scoring (GCE style) \hfill (13 marks)]
A footballer kicks a ball from ground level with speed \SI{16}{m.s^{-1}} at an angle $\theta$ above the horizontal. Take $g = \SI{10}{m.s^{-2}}$ and neglect air resistance.
\begin{enumerate}
\item Write down the horizontal and vertical components of the displacement of the ball at time $t$, and hence show that the Cartesian equation of the trajectory is
\[ y = x\tan\theta - \frac{5x^2}{128}\left(1 + \tan^2\theta\right). \hfill \text{(4 marks)} \]
\item A vertical wall \SI{2}{m} high stands \SI{9}{m} from the kicker. Show that for the ball to clear the wall, $\tan\theta$ must satisfy a certain inequality. \hfill (3 marks)
\item The crossbar of the goal, \SI{2.4}{m} high, is \SI{18}{m} away. Show that for the ball to pass under the crossbar and above the goal line at that point, $\tan\theta$ satisfies a quadratic equation at the limiting case, and find the possible values of $\theta$ for which the ball both clears the wall and enters the goal. \hfill (6 marks)
\end{enumerate}
\end{exercice}

\begin{exercice}[Projectile from a rooftop (GCE style) \hfill (12 marks)]
A projectile is launched from the top of a building \SI{20}{m} high with speed \SI{15}{m.s^{-1}} at $30^\circ$ above the horizontal, and lands in the horizontal school field below. Take $g = \SI{10}{m.s^{-2}}$.
\begin{enumerate}
\item Write down the equations giving the horizontal and vertical displacements of the projectile at time $t$ after launch, measured from the foot of the building. \hfill (3 marks)
\item Show that the time of flight $T$ satisfies $5T^2 - 7.5T - 20 = 0$, and hence find $T$. \hfill (4 marks)
\item Find the horizontal distance travelled by the projectile before landing. \hfill (2 marks)
\item Find the speed of the projectile and the angle its velocity makes with the horizontal at the instant it lands. \hfill (3 marks)
\end{enumerate}
\end{exercice}

\begin{exercice}[Vertical motion and its graphs (GCE style) \hfill (14 marks)]
Take $g = \SI{10}{m.s^{-2}}$ throughout this question.
\begin{enumerate}
\item A mango falls from rest from a tree and takes \SI{1.5}{s} to reach the ground. From what height did it fall, and with what speed does it hit the ground? \hfill (3 marks)
\item A stone is thrown vertically upwards at \SI{20}{m.s^{-1}} from the edge of a well and falls to the bottom of the well, \SI{25}{m} below the point of projection. Find the total time of flight and the speed of the stone at the bottom. \hfill (4 marks)
\item An object is released from a balloon which is rising vertically at \SI{8}{m.s^{-1}} when the balloon is \SI{100}{m} above the ground. Explain why the object initially moves upwards, and find the time taken for the object to reach the ground and its speed on impact. \hfill (4 marks)
\item A ball is thrown vertically upwards at \SI{14}{m.s^{-1}} and returns to the thrower's hand. Sketch, on separate axes, the displacement--time, velocity--time and acceleration--time graphs from launch until the ball returns to the thrower's hand, marking all key values. \hfill (3 marks)
\end{enumerate}
\end{exercice}

\newpage

\coursec{Solutions to the exercises}

\coursec{Solutions d\'etaill\'ees des exercices du chapitre PMM11}

\begin{corrige}[Exercice 1 -- Questions \`a choix multiple]
\begin{enumerate}
\item \textbf{R\'eponse (a).} La vitesse est la d\'eriv\'ee de la position : $\vec{v}(t) = \dfrac{d\vec{r}}{dt} = (6t)\vec{i} + (4 - 2t)\vec{j}$. Les options (b) et (c) conservent des termes non d\'eriv\'es ; (d) est l'acc\'el\'eration, pas la vitesse.
\item \textbf{R\'eponse (d).} Au point le plus haut, la composante verticale de la vitesse est instantan\'ement z\'ero (le projectile cesse de monter et commence \`a descendre), mais la composante horizontale $u\cos\theta$ est inchang\'ee, et l'acc\'el\'eration est $\vec{g}$ vers le bas pendant tout le vol. Donc (a), (b) et (c) sont toutes fausses.
\item \textbf{R\'eponse (b).} La port\'ee sur un sol horizontal est $R = \dfrac{u^2\sin 2\theta}{g}$. L'option (c) est la hauteur maximale ; l'option (d) est le temps de vol.
\item \textbf{R\'eponse (c).} Pour un corps l\^ach\'e sans vitesse initiale, $h = \dfrac{1}{2}gt^2$, donc $20 = \dfrac{1}{2}(10)t^2$, d'o\`u $t^2 = 4$ et $t = \SI{2}{s}$.
\end{enumerate}
\end{corrige}

\begin{corrige}[Exercice 2 -- Vrai ou faux ?]
\begin{enumerate}
\item \textbf{Faux.} Une acc\'el\'eration constante signifie $\vec{v}(t) = \vec{u} + \vec{a}t$ : la vitesse change uniform\'ement. Une acc\'el\'eration constante implique une vitesse constante seulement si $\vec{a} = \vec{0}$. Contre-exemple : un projectile a une acc\'el\'eration constante $\vec{g}$ mais sa vitesse change continuellement.
\item \textbf{Vrai.} $R = \dfrac{u^2\sin 2\theta}{g}$, et $\sin 2(90^\circ - \theta) = \sin(180^\circ - 2\theta) = \sin 2\theta$. Donc des angles compl\'ementaires donnent la m\^eme port\'ee (mais des hauteurs maximales et des temps de vol diff\'erents).
\item \textbf{Faux.} Au sommet du vol, la \emph{vitesse} est z\'ero, mais l'acc\'el\'eration est $g \approx \SI{9.81}{m.s^{-2}}$ vers le bas pendant tout le vol -- la gravitation ne s'arr\^ete jamais. Si l'acc\'el\'eration \'etait nulle au sommet, le ballon y resterait.
\item \textbf{Vrai.} En n\'egligeant la r\'esistance de l'air, la seule force est le poids, qui est vertical. Donc l'acc\'el\'eration horizontale est nulle, donc $v_x = u\cos\theta$ est constante pendant tout le vol.
\end{enumerate}
\end{corrige}

\begin{corrige}[Exercice 3 -- D\'efinitions cl\'es]
\begin{enumerate}
\item La \cle{vitesse} est le taux de variation de la position : $\vec{v}(t) = \dfrac{d\vec{r}}{dt}$. L'\cle{acc\'el\'eration} est le taux de variation de la vitesse : $\vec{a}(t) = \dfrac{d\vec{v}}{dt} = \dfrac{d^2\vec{r}}{dt^2}$.
\item Dans le mod\`ele standard du projectile, $\vec{a} = -g\,\vec{j}$, c'est-\`a-dire composante horizontale $a_x = 0$ (aucune force horizontale, donc aucune acc\'el\'eration horizontale) et composante verticale $a_y = -g$ (le poids est la seule force, agissant verticalement vers le bas, et d'apr\`es la deuxi\`eme loi de Newton $m\vec{a} = m\vec{g}$).
\item Le \cle{temps de vol} est le temps total entre le lancement et l'atterrissage. La \cle{port\'ee} est la distance horizontale entre le point de lancement et le point d'atterrissage. La \cle{hauteur maximale} est le plus grand d\'eplacement vertical atteint au-dessus du point de lancement.
\item La \cle{chute libre} est le mouvement d'un corps soumis \`a son poids seul (r\'esistance de l'air n\'eglig\'ee). Pr\`es de la surface de la Terre, l'acc\'el\'eration de la chute libre est $g \approx \SI{9.81}{m.s^{-2}}$ (souvent prise comme \SI{10}{m.s^{-2}} dans les exercices), dirig\'ee verticalement vers le bas.
\end{enumerate}
\end{corrige}

\begin{corrige}[Exercice 4 -- Calculs directs]
\begin{enumerate}
\item Acc\'el\'eration : $\vec{a} = \dfrac{d\vec{v}}{dt} = 4\,\vec{i} - 2\,\vec{j}\ \SI{}{m.s^{-2}}$ (constante). \`a $t = 2$ : $\vec{v}(2) = 8\,\vec{i} - 1\,\vec{j}$, donc la vitesse (module) est $|\vec{v}(2)| = \sqrt{8^2 + (-1)^2} = \sqrt{65} \approx \SI{8.06}{m.s^{-1}}$.
\item L\^ach\'e sans vitesse initiale : $h = \dfrac{1}{2}gt^2 = \dfrac{1}{2}(10)(3^2) = \SI{45}{m}$. Vitesse \`a l'impact : $v = gt = 10 \times 3 = \SI{30}{m.s^{-1}}$.
\item $u_x = u\cos\theta = 20\cos 30^\circ = 20 \times \dfrac{\sqrt{3}}{2} = 10\sqrt{3} \approx \SI{17.3}{m.s^{-1}}$ ; $u_y = u\sin\theta = 20\sin 30^\circ = \SI{10}{m.s^{-1}}$.
\item Sens vers le haut positif : $v = u - gt$. \`a la plus grande hauteur $v = 0$, donc $t = \dfrac{u}{g} = \dfrac{15}{10} = \SI{1.5}{s}$. Hauteur : $h = ut - \dfrac{1}{2}gt^2 = 15(1.5) - 5(1.5)^2 = 22.5 - 11.25 = \SI{11.25}{m}$ (equivalentment $h = \dfrac{u^2}{2g} = \dfrac{225}{20}$).
\end{enumerate}
\end{corrige}

\begin{corrige}[Exercice 5 -- D\'erivation d'un vecteur position]
\begin{enumerate}
\item En d\'erivant composante par composante :
\[ \vec{v}(t) = \frac{d\vec{r}}{dt} = (4t - 3)\vec{i} + (3t^2 - 6)\vec{j}, \qquad \vec{a}(t) = \frac{d\vec{v}}{dt} = 4\,\vec{i} + 6t\,\vec{j}. \]
\item \`a $t = 2$ : $\vec{v}(2) = 5\,\vec{i} + 6\,\vec{j}\ \SI{}{m.s^{-1}}$ et $\vec{a}(2) = 4\,\vec{i} + 12\,\vec{j}\ \SI{}{m.s^{-2}}$. Vitesse : $|\vec{v}(2)| = \sqrt{25 + 36} = \sqrt{61} \approx \SI{7.81}{m.s^{-1}}$.
\item Le mouvement est parall\`ele \`a l'axe $y$ lorsque la composante horizontale de la vitesse est nulle (et la verticale non nulle) : $4t - 3 = 0 \Rightarrow t = \dfrac{3}{4}\ \SI{}{s}$. (V\'erification : $3t^2 - 6 = \frac{27}{16} - 6 \neq 0$, la particule bouge bien.)
\item On a $x = 2t^2 - 3t$ et $y = t^3 - 6t = t(t^2 - 6)$. Le param\`etre $t$ ne peut pas \^etre \'elimin\'e par une alg\`ebre \'el\'ementaire sous une forme close simple car $x$ est quadratique et $y$ cubique en $t$ ; la trajectoire est donn\'ee sous forme param\'etrique par
\[ x = 2t^2 - 3t, \qquad y = t^3 - 6t, \]
et l'\'elimination de $t$ (par exemple en r\'esolvant $2t^2 - 3t - x = 0$, $t = \dfrac{3 \pm \sqrt{9 + 8x}}{4}$, et en substituant dans $y$) donne la relation cart\'esienne
\[ y = \left(\frac{3 \pm \sqrt{9+8x}}{4}\right)\left(\left(\frac{3 \pm \sqrt{9+8x}}{4}\right)^{\!2} - 6\right), \]
qu'il est pr\'ef\'erable de laisser sous forme param\'etrique. (Au niveau GCE, \'ecrire les \'equations param\'etriques et la proc\'edure de substitution rapporte tous les points.)
\end{enumerate}
\end{corrige}

\begin{corrige}[Exercice 6 -- Int\'egration d'une acc\'el\'eration]
\begin{enumerate}
\item $\vec{v}(t) = \displaystyle\int \vec{a}\,dt = (t^2 + c_1)\vec{i} + (-6t + c_2)\vec{j}$. \`a $t = 0$, $\vec{v} = 4\vec{i} + 3\vec{j}$, donc $c_1 = 4$, $c_2 = 3$ :
\[ \vec{v}(t) = (t^2 + 4)\vec{i} + (3 - 6t)\vec{j}\ \SI{}{m.s^{-1}}. \]
\item $\vec{r}(t) = \displaystyle\int \vec{v}\,dt = \left(\frac{t^3}{3} + 4t + k_1\right)\vec{i} + \left(3t - 3t^2 + k_2\right)\vec{j}$. Comme $\vec{r}(0) = \vec{0}$, $k_1 = k_2 = 0$ :
\[ \vec{r}(t) = \left(\frac{t^3}{3} + 4t\right)\vec{i} + (3t - 3t^2)\vec{j}\ \SI{}{m}. \]
\item \`a $t = 3$ : $\vec{r}(3) = (9 + 12)\vec{i} + (9 - 27)\vec{j} = 21\vec{i} - 18\vec{j}$. Distance \`a $O$ :
\[ |\vec{r}(3)| = \sqrt{21^2 + 18^2} = \sqrt{441 + 324} = \sqrt{765} = 3\sqrt{85} \approx \SI{27.7}{m}. \]
\end{enumerate}
\end{corrige}

\begin{corrige}[Exercice 7 -- Projectile standard sur sol horizontal]
Composantes initiales : $u_x = 25\cos 40^\circ \approx 25(0.766) = \SI{19.15}{m.s^{-1}}$ ; $u_y = 25\sin 40^\circ \approx 25(0.643) = \SI{16.075}{m.s^{-1}}$.
\begin{enumerate}
\item Temps de vol : $T = \dfrac{2u\sin\theta}{g} = \dfrac{2 \times 16.075}{10} = \SI{3.215}{s} \approx \SI{3.22}{s}$ (3 ch. sign.).
\item Hauteur maximale : $H = \dfrac{u^2\sin^2\theta}{2g} = \dfrac{(16.075)^2}{20} = \dfrac{258.4}{20} \approx \SI{12.9}{m}$ (3 ch. sign.).
\item Port\'ee : $R = \dfrac{u^2\sin 2\theta}{g} = \dfrac{625 \times 2\sin 40^\circ \cos 40^\circ}{10} = \dfrac{625 \times 2(0.643)(0.766)}{10} \approx \dfrac{625 \times 0.985}{10} \approx \SI{61.6}{m}$ (3 ch. sign.).
\end{enumerate}
\end{corrige}

\begin{corrige}[Exercice 8 -- Lancer horizontal depuis une falaise]
On place l'origine au pied de la falaise : $x$ horizontal, $y$ vertical. Vitesse initiale : $u_x = 12$, $u_y = 0$ ; hauteur initiale $y_0 = 45$.
\begin{enumerate}
\item Mouvement vertical : $y = 45 - \dfrac{1}{2}gt^2 = 45 - 5t^2$. Impact quand $y = 0$ : $t^2 = 9$, donc $T = \SI{3}{s}$.
\item Distance horizontale : $x = u_x T = 12 \times 3 = \SI{36}{m}$ du pied de la falaise.
\item \`a l'impact : $v_x = 12$ (constant), $v_y = -gT = -30\ \SI{}{m.s^{-1}}$.
\[ |\vec{v}| = \sqrt{12^2 + 30^2} = \sqrt{144 + 900} = \sqrt{1044} \approx \SI{32.3}{m.s^{-1}}. \]
Direction sous l'horizontale : $\tan\alpha = \dfrac{30}{12} = 2.5$, donc $\alpha \approx 68.2^\circ$ sous l'horizontale.
\end{enumerate}
\end{corrige}

\begin{corrige}[Exercice 9 -- Franchir le mur et marquer (style GCE)]
\begin{enumerate}
\item Avec $u = 16$, $u_x = 16\cos\theta$, $u_y = 16\sin\theta$ :
\[ x = 16t\cos\theta, \qquad y = 16t\sin\theta - 5t^2. \]
On \'elimine $t$ : $t = \dfrac{x}{16\cos\theta}$, d'o\`u
\[ y = 16\sin\theta \cdot \frac{x}{16\cos\theta} - 5\left(\frac{x}{16\cos\theta}\right)^2 = x\tan\theta - \frac{5x^2}{256\cos^2\theta} = x\tan\theta - \frac{5x^2}{256}\left(1 + \tan^2\theta\right), \]
en utilisant $\dfrac{1}{\cos^2\theta} = \sec^2\theta = 1 + \tan^2\theta$. $\blacksquare$

\emph{Bar\`eme : M1 composantes du d\'eplacement ; M1 \'elimination de $t$ ; M1 utilisation de $\sec^2\theta = 1 + \tan^2\theta$ ; A1 r\'esultat.}

\item Le ballon franchit le mur si $y > 2$ quand $x = 9$ :
\[ 9\tan\theta - \frac{5 \times 81}{256}\left(1 + \tan^2\theta\right) > 2. \]
Posons $m = \tan\theta$ : $9m - \dfrac{405}{256}(1 + m^2) > 2$, soit
\[ 405m^2 - 2304m + 917 < 0 \quad (\text{apr\`es multiplication par } 256 \text{ et r\'earrangement}). \]
Donc $\tan\theta$ doit rendre le trin\^ome $405m^2 - 2304m + 917$ n\'egatif, c'est-\`a-dire que $m$ doit se trouver strictement entre ses deux racines. $\blacksquare$

\emph{Bar\`eme : M1 substitution $x = 9$, $y > 2$ ; M1 formation de l'in\'equation du second degr\'e ; A1 in\'equation correcte.}

\item Au but, $x = 18$ ; le ballon entre dans le but si $0 < y < 2.4$.
Cas limite pour la barre transversale ($y = 2.4$) :
\[ 18m - \frac{5 \times 324}{256}(1 + m^2) = 2.4 \;\Longrightarrow\; 18m - \frac{405}{64}(1 + m^2) = 2.4. \]
Multipliant par $64$ : $1152m - 405 - 405m^2 = 153.6$, soit
\[ 405m^2 - 1152m + 558.6 = 0. \]
Discriminant : $\Delta_c = 1152^2 - 4(405)(558.6) = 1327104 - 904884 = 422220 > 0$.
Racines : $m = \dfrac{1152 \pm \sqrt{422220}}{810} \approx \dfrac{1152 \pm 649.78}{810}$, donnant $m_1 \approx 0.620$, $m_2 \approx 2.224$.
Le trin\^ome \'etant convexe (coefficient de $m^2$ positif), $y < 2.4$ pour $m < 0.620$ \textbf{ou} $m > 2.224$.

Condition pour \^etre au-dessus du sol ($y > 0$ \`a $x = 18$) :
\[ 18m - \frac{405}{64}(1+m^2) > 0 \;\Longrightarrow\; 405m^2 - 1152m + 405 < 0. \]
Discriminant : $\Delta_g = 1152^2 - 4(405)^2 = 1327104 - 656100 = 671004 > 0$.
Racines : $m = \dfrac{1152 \pm \sqrt{671004}}{810} \approx \dfrac{1152 \pm 819.15}{810}$, donnant $m_3 \approx 0.411$, $m_4 \approx 2.433$.
Donc $y > 0$ pour $0.411 < m < 2.433$.

Condition du mur (d\'ej\`a trouv\'ee) : $0.4306 < m < 5.258$ (racines de $405m^2 - 2304m + 917 = 0$ : $m \approx 0.4306$ et $m \approx 5.258$).

Intersection des trois conditions :
\begin{itemize}
\item Branche basse : $\max(0.4306, 0.411) < m < \min(0.620, 2.433)$ $\Rightarrow$ $0.4306 < m < 0.620$.
\item Branche haute : $\max(2.224, 0.4306) < m < \min(2.433, 5.258)$ $\Rightarrow$ $2.224 < m < 2.433$.
\end{itemize}
En degr\'es : $\theta \in (23.3^\circ,\ 31.8^\circ) \cup (65.8^\circ,\ 67.6^\circ)$.
Pour ces angles, le ballon franchit le mur, passe sous la barre et au-dessus du sol \`a $x = 18$ : il marque.

\emph{Bar\`eme : M1 substitution $x = 18$ ; A1 trin\^omes en $\tan\theta$ ; M1 analyse des discriminants et des signes ; M1 r\'esolution de l'in\'equation du mur ; A1 racines ; A1 intervalles finaux pour $\theta$. Attente de l'examinateur : manipulation alg\'ebrique rigoureuse et \'enonc\'e clair de la gamme d'angles.}
\end{enumerate}
\end{corrige}

\begin{corrige}[Exercice 10 -- Projectile depuis un toit (style GCE)]
\begin{enumerate}
\item Origine au pied de l'immeuble, $y$ vers le haut. Point de lancement $(0, 20)$, $u_x = 15\cos 30^\circ = \dfrac{15\sqrt{3}}{2} \approx 12.99$, $u_y = 15\sin 30^\circ = 7.5$ :
\[ x = 7.5\sqrt{3}\,t \approx 12.99t, \qquad y = 20 + 7.5t - 5t^2. \]
\emph{Bar\`eme : M1 composantes ; A1 $x(t)$ ; A1 $y(t)$ avec le $+20$.}

\item Atterrissage quand $y = 0$ : $20 + 7.5T - 5T^2 = 0$, soit $5T^2 - 7.5T - 20 = 0$. $\blacksquare$
\[ T = \frac{7.5 \pm \sqrt{7.5^2 + 400}}{10} = \frac{7.5 \pm \sqrt{456.25}}{10} = \frac{7.5 \pm 21.36}{10}. \]
On prend la racine positive : $T \approx \SI{2.89}{s}$ (la racine n\'egative $-1.39$ s est rejet\'ee comme non physique).
\emph{Bar\`eme : M1 \'equation $y = 0$ ; A1 \'equation donn\'ee ; M1 formule quadratique ; A1 racine positive choisie avec justification.}

\item Distance horizontale : $x = 7.5\sqrt{3} \times 2.886 \approx 12.99 \times 2.886 \approx \SI{37.5}{m}$.
\emph{Bar\`eme : M1 $x = u_x T$ ; A1 valeur.}

\item \`a l'atterrissage : $v_x = 7.5\sqrt{3} \approx \SI{12.99}{m.s^{-1}}$ ; $v_y = 7.5 - 10T = 7.5 - 28.86 = -\SI{21.36}{m.s^{-1}}$.
\[ |\vec{v}| = \sqrt{12.99^2 + 21.36^2} = \sqrt{168.7 + 456.3} = \sqrt{625} = \SI{25}{m.s^{-1}}. \]
Angle sous l'horizontale : $\tan\alpha = \dfrac{21.36}{12.99} = 1.644$, donc $\alpha \approx 58.7^\circ$ sous l'horizontale.
\emph{Bar\`eme : M1 composantes \`a $T$ ; A1 vitesse ; A1 angle avec direction pr\'ecis\'ee.}
\end{enumerate}
\end{corrige}

\begin{corrige}[Exercice 11 -- Mouvement vertical et ses graphiques (style GCE)]
\begin{enumerate}
\item Chute sans vitesse initiale : $h = \dfrac{1}{2}gt^2 = \dfrac{1}{2}(10)(1.5)^2 = \SI{11.25}{m}$. Vitesse \`a l'impact : $v = gt = 10 \times 1.5 = \SI{15}{m.s^{-1}}$.
\emph{Bar\`eme : M1 formule correcte ; A1 hauteur ; A1 vitesse.}

\item On prend le sens vers le haut positif, origine au bord du puits. D\'eplacement au fond : $s = -25$ m. En utilisant $s = ut - \dfrac{1}{2}gt^2$ :
\[ -25 = 20t - 5t^2 \;\Longrightarrow\; 5t^2 - 20t - 25 = 0 \;\Longrightarrow\; t^2 - 4t - 5 = 0 \;\Longrightarrow\; (t - 5)(t + 1) = 0. \]
Donc $t = \SI{5}{s}$ (on rejette $t = -1$). Vitesse au fond : $v = u - gt = 20 - 50 = -30\ \SI{}{m.s^{-1}}$, soit \SI{30}{m.s^{-1}} vers le bas. (V\'erification avec $v^2 = u^2 + 2g(25) = 400 + 500 = 900$.)
\emph{Bar\`eme : M1 convention de signe et $s = -25$ ; A1 \'equation du second degr\'e ; A1 $t = 5$ s avec rejet de la racine n\'egative ; A1 vitesse.}

\item L'objet a la vitesse ascendante du ballon \SI{8}{m.s^{-1}} \`a l'instant du l\^acher : il n'est pas l\^ach\'e sans vitesse initiale. Il continue donc vers le haut, en d\'ec\'el\'erant de $g$, jusqu'\`a ce que sa vitesse devienne nulle, et seulement alors il retombe. Avec l'origine au point de l\^acher, $u = +8$, $s = -100$ :
\[ -100 = 8t - 5t^2 \;\Longrightarrow\; 5t^2 - 8t - 100 = 0 \;\Longrightarrow\; t = \frac{8 \pm \sqrt{64 + 2000}}{10} = \frac{8 \pm \sqrt{2064}}{10} \approx \frac{8 \pm 45.43}{10}. \]
Racine positive : $t \approx \SI{5.34}{s}$. Vitesse \`a l'impact : $v = 8 - 10(5.34) = -45.4\ \SI{}{m.s^{-1}}$, soit environ \SI{45.4}{m.s^{-1}} vers le bas.
\emph{Bar\`eme : B1 explication du mouvement initial vers le haut (inertie / vitesse initiale du ballon) ; M1 \'equation avec $u = 8$, $s = -100$ ; A1 temps ; A1 vitesse.}

\item Avec $u = 14$ : temps au sommet $t_1 = \dfrac{u}{g} = \SI{1.4}{s}$ ; temps total $T = \SI{2.8}{s}$ ; hauteur maximale $h = \dfrac{u^2}{2g} = \dfrac{196}{20} = \SI{9.8}{m}$.

\textbf{D\'eplacement--temps :} parabole $s = 14t - 5t^2$ de $(0,0)$, sommet $(1.4,\ 9.8)$, retour \`a $(2.8,\ 0)$.

\textbf{Vitesse--temps :} droite $v = 14 - 10t$ de $+14$ \`a $t = 0$, passant par $0$ \`a $t = 1.4$, jusqu'\`a $-14\ \SI{}{m.s^{-1}}$ \`a $t = 2.8$.

\textbf{Acc\'el\'eration--temps :} ligne horizontale $a = -10\ \SI{}{m.s^{-2}}$ pendant tout le vol.

\begin{popfigure}
\begin{center}
\begin{tikzpicture}[scale=0.9]
% s-t graph
\begin{scope}
\draw[->] (0,0) -- (3.4,0) node[right] {$t$};
\draw[->] (0,0) -- (0,2.4) node[above] {$s$ (m)};
\draw[popblue, thick, domain=0:2.8, samples=60] plot (\x, {2*(\x)*(2.8-\x)/1.96});
\draw[dashed, gray] (1.4,0) node[below] {$1.4$} -- (1.4,2) -- (0,2) node[left] {$9.8$};
\node[below] at (2.8,0) {$2.8$};
\node[below] at (1.7,-0.6) {$s = 14t - 5t^2$};
\end{scope}
% v-t graph
\begin{scope}[xshift=5.2cm]
\draw[->] (0,-1.6) -- (0,2.2) node[above] {$v$ (m/s)};
\draw[->] (0,0) -- (3.4,0) node[right] {$t$};
\draw[poporange, thick] (0,1.6) -- (2.8,-1.4);
\node[left] at (0,1.6) {$14$};
\node[left] at (0,-1.4) {$-14$};
\draw[dashed, gray] (1.4,0.1) -- (1.4,-0.1) node[below] {$1.4$};
\node[below] at (2.8,0) {$2.8$};
\node[below] at (1.7,-1.9) {$v = 14 - 10t$};
\end{scope}
% a-t graph
\begin{scope}[xshift=10.4cm]
\draw[->] (0,-1.8) -- (0,1.2) node[above] {$a$ (m/s$^2$)};
\draw[->] (0,0) -- (3.4,0) node[right] {$t$};
\draw[popgreen, thick] (0,-1.2) -- (2.8,-1.2);
\node[left] at (0,-1.2) {$-10$};
\node[below] at (2.8,0) {$2.8$};
\node[below] at (1.7,-2.1) {$a = -g$ (constant)};
\end{scope}
\end{tikzpicture}
\end{center}
\legende{Graphiques d\'eplacement--temps, vitesse--temps et acc\'el\'eration--temps pour la balle lanc\'ee vers le haut \`a \SI{14}{m.s^{-1}}.}
\end{popfigure}

\emph{Bar\`eme : B1 parabole avec sommet et z\'eros corrects ; B1 droite vitesse traversant z\'ero \`a $t = 1.4$ avec valeurs $\pm 14$ ; B1 ligne acc\'el\'eration constante \`a $-10$. Attente de l'examinateur : le graphique de la vitesse doit devenir \emph{n\'egatif} (la balle redescend) et le graphique de l'acc\'el\'eration doit montrer que $a$ ne s'annule jamais.}
\end{enumerate}
\end{corrige}

\newpage

\coursec{Summary sheet}

\begin{popfigure}
\begin{tikzpicture}[
  mindmap,
  grow cyclic,
  every node/.style={concept, text=white, font=\small},
  concept color=popblue,
  level 1/.append style={level distance=3.4cm, sibling angle=60},
  level 2/.append style={level distance=2.0cm, sibling angle=45, font=\scriptsize},
  text width=2.6cm
]
\node[concept color=popdark, text width=3.2cm, font=\small\bfseries]{Kinematics of Particles in a Plane}
  child[concept color=popblue]{ node {Position, velocity, acceleration as vectors}
    child { node[text width=2.2cm] {$\vec{v}=\dot{\vec{r}}$, $\vec{a}=\dot{\vec{v}}$} }
    child { node[text width=2.2cm] {Differentiate componentwise} }
  }
  child[concept color=popgreen]{ node {Integration: back to $\vec{v}$ and $\vec{r}$}
    child { node[text width=2.2cm] {Constants from initial conditions} }
  }
  child[concept color=poporange]{ node {Projectile model}
    child { node[text width=2.2cm] {$a_x=0$, $a_y=-g$} }
    child { node[text width=2.2cm] {No air resistance} }
  }
  child[concept color=popgold]{ node {Projectile equations}
    child { node[text width=2.2cm] {$x=ut\cos\theta$} }
    child { node[text width=2.2cm] {$y=ut\sin\theta-\tfrac12 gt^2$} }
    child { node[text width=2.2cm] {Trajectory: parabola} }
  }
  child[concept color=poppurple]{ node {Range, height, flight time}
    child { node[text width=2.2cm] {$R=\frac{u^2\sin2\theta}{g}$} }
    child { node[text width=2.2cm] {$H=\frac{u^2\sin^2\theta}{2g}$} }
    child { node[text width=2.2cm] {$T=\frac{2u\sin\theta}{g}$} }
  }
  child[concept color=poppink]{ node {Free fall and vertical motion}
    child { node[text width=2.2cm] {Thrown up: $v=u-gt$} }
    child { node[text width=2.2cm] {Graphs: $x$--$t$, $v$--$t$} }
  };
\end{tikzpicture}
\legende{Mind map of the chapter: from vector kinematics to projectile motion and free fall.}
\end{popfigure}

\begin{bilan}
\textbf{Key definitions}
\begin{itemize}
\item \cle{Position vector}: $\vec{r}(t)=x(t)\vec{i}+y(t)\vec{j}$ locates the particle at time $t$.
\item \cle{Velocity}: $\vec{v}=\dfrac{d\vec{r}}{dt}=\dot{x}\,\vec{i}+\dot{y}\,\vec{j}$; it is tangent to the path. Speed $=|\vec{v}|=\sqrt{\dot{x}^2+\dot{y}^2}$.
\item \cle{Acceleration}: $\vec{a}=\dfrac{d\vec{v}}{dt}=\ddot{x}\,\vec{i}+\ddot{y}\,\vec{j}$.
\item \cle{Projectile}: a body launched with initial velocity $\vec{u}$ and then moving under gravity alone (air resistance neglected).
\item \cle{Free fall}: motion with $\vec{a}=-g\vec{j}$, $g\approx 9.81\ \mathrm{m\,s^{-2}}$ (take $g=9.81$ or $10$ as instructed).
\end{itemize}

\textbf{Formulas and laws to know}
\begin{itemize}
\item Differentiation: $\vec{r}\ \xrightarrow{\ d/dt\ }\ \vec{v}\ \xrightarrow{\ d/dt\ }\ \vec{a}$. Integration reverses the chain; each integration introduces a constant fixed by initial conditions.
\item Projectile with launch speed $u$ at angle $\theta$ from the origin:
\[ x=ut\cos\theta, \qquad y=ut\sin\theta-\tfrac12 gt^2. \]
\[ v_x=u\cos\theta \ (\text{constant}), \qquad v_y=u\sin\theta-gt. \]
\item Cartesian equation of the \cle{trajectory} (a parabola):
\[ y=x\tan\theta-\frac{g\,x^2}{2u^2\cos^2\theta}=x\tan\theta-\frac{g x^2}{2u^2}\bigl(1+\tan^2\theta\bigr). \]
\item On horizontal ground:
\[ \text{Time of flight } T=\frac{2u\sin\theta}{g},\qquad \text{Range } R=\frac{u^2\sin 2\theta}{g},\qquad \text{Max height } H=\frac{u^2\sin^2\theta}{2g}. \]
\item Maximum range: $R_{\max}=\dfrac{u^2}{g}$ at $\theta=45^\circ$; complementary angles $\theta$ and $90^\circ-\theta$ give the same range.
\item At the highest point: $v_y=0$ but $v_x=u\cos\theta\neq 0$; time to reach it is $T/2$.
\item Vertical throw upwards (taking up as positive): $v=u-gt$, $y=ut-\tfrac12 gt^2$, $v^2=u^2-2gy$; greatest height $\dfrac{u^2}{2g}$; the body returns with speed $u$ (symmetry of free fall).
\item Graphs: gradient of $x$--$t$ graph $=$ velocity; gradient of $v$--$t$ graph $=$ acceleration; area under $v$--$t$ graph $=$ displacement. For vertical motion the $v$--$t$ graph is a straight line of gradient $-g$.
\end{itemize}

\textbf{Methods}
\begin{itemize}
\item \emph{Vector kinematics:} differentiate or integrate \textbf{each component separately}; use the given position/velocity at $t=0$ to find constants.
\item \emph{Projectile problems:} (1) resolve $\vec{u}$ into $u\cos\theta$ and $u\sin\theta$; (2) write $x(t)$ and $y(t)$; (3) translate the question into a condition ($y=0$: landing; $v_y=0$: summit; $y=h$: target height); (4) solve for $t$, then substitute.
\item \emph{Finding $\theta$ or $u$ from the trajectory:} substitute the coordinates of a known point into the Cartesian equation; use $1/\cos^2\theta=1+\tan^2\theta$ to get a quadratic in $\tan\theta$.
\item \emph{Vertical throws:} choose a positive direction once and keep it; $a=-g$ throughout, even at the top.
\end{itemize}

\textbf{Common mistakes}
\begin{itemize}
\item Mixing degrees and radians, or using $\sin\theta$ where $\cos\theta$ is needed when resolving $\vec{u}$.
\item Believing velocity is zero at the highest point: only the \emph{vertical} component vanishes.
\item Using the range formula when launch and landing are at \textbf{different heights} — go back to $y(t)=0$ instead.
\item Forgetting the constant of integration or the units ($\mathrm{m}$, $\mathrm{m\,s^{-1}}$, $\mathrm{m\,s^{-2}}$).
\item Sign errors in vertical motion: acceleration is $-g$ on the way up \emph{and} on the way down.
\item Confusing distance travelled with displacement on a $v$--$t$ graph (areas below the axis count negatively for displacement).
\end{itemize}

\textbf{GCE tips}
\begin{itemize}
\item State your model clearly: ``air resistance neglected, $g=9.81\ \mathrm{m\,s^{-2}}$'' earns marks.
\item Show the resolved components $u\cos\theta$, $u\sin\theta$ before writing equations — method marks are awarded there.
\item Keep exact values ($\sin 30^\circ=\tfrac12$, etc.) until the final line, then round to 3 significant figures.
\item For ``show that'' questions, derive from $x(t)$, $y(t)$ — never quote the result you are asked to prove.
\item Sketch the parabola and mark the given data; it prevents sign and angle errors and helps the examiner follow your work.
\end{itemize}
\end{bilan}

\findecours

\end{document}
